% ===================================================================
%  Dodatek F: Hierarchia mas cząstek — węzły profilu radialnego
%  TGP v1 — trzy generacje z kwantowych liczb węzłowych
% ===================================================================
\section{Hierarchia mas cz\k{a}stek: w\k{e}z\l{}y profilu radialnego}%
\label{app:hierarchia-mas}

\subsection{Przypomnienie: cz\k{a}stka jako defekt radialny}%
\label{app:F-reminder}

Z~hipotezy~\ref{hyp:generacje} (sekcja~\ref{ssec:defekty}):
trzy generacje cz\k{a}stek elementarnych odpowiadaj\k{a}
\textbf{liczbie w\k{e}z\l{}\'ow} $n$ profilu radialnego $\Phi(r)$
kinku:
\begin{itemize}[nosep]
    \item $n = 0$ --- profil bezw\k{e}z\l{}owy: najl\.zejsza
          generacja (elektron $e$, kwarki $u$, $d$),
    \item $n = 1$ --- jeden w\k{e}ze\l{}: \'srednia generacja
          (mion $\mu$, kwarki $c$, $s$),
    \item $n = 2$ --- dwa w\k{e}z\l{}y: najci\k{e}\.zsza generacja
          (ta on $\tau$, kwarki $t$, $b$).
\end{itemize}
Niniejszy dodatek \textbf{formalizuje} t\k{e} hipotez\k{e}
i~wyprowadza ilo\'sciowe ograniczenia na stosunki mas.

\medskip
\noindent\fbox{\parbox{\linewidth}{%
\textbf{Nota kanoniczna (2026-04-07).}
Niniejszy dodatek u\.zywa Formulacji~B ($\varphi(a_\Gamma) \approx 1{,}24$).
W~sformu\l{}owaniu kanonicznym: $K(g) = g^4$, $P(g) = (\beta/7)g^7 - (\gamma/8)g^8$,
$g_0^e = 0{,}86941$, brak punktu duchowego.
Algebra Koide'a i~r\'ownania master F1--F12 s\k{a} niezale\.zne od tego wyboru.}}
\medskip

% -------------------------------------------------------------------
\subsection{R\'ownanie radialne kinku TGP}%
\label{app:F-radial-eq}
% -------------------------------------------------------------------

Dla sferycznie symetrycznego kinku statycznego ($\Phi = \Phi(r)$)
r\'ownanie pola~\eqref{eq:full-field-alt} w~obecno\'sci
efektywnego potencja\l{}u (brak zewn\k{e}trznego \'zr\'od\l{}a)
przyjmuje posta\'c:
\begin{equation}\label{eq:radial-kink}
    \Phi'' + \frac{2}{r}\,\Phi'
    + \frac{\alpha}{\Phi}\,(\Phi')^2
    - V'_{\mathrm{eff}}(\Phi)
    = 0,
\end{equation}
gdzie $\alpha = 2$ (wynikaj\k{a}cy z~wariacji dzia\l{}ania)
i~potencja\l{} efektywny:
\begin{equation}\label{eq:Veff-kink}
    V_{\mathrm{eff}}(\Phi)
    = \frac{\beta}{3}\,\frac{\Phi^3}{\PhiZero^2}
    - \frac{\gamma}{4}\,\frac{\Phi^4}{\PhiZero^3},
    \qquad
    V'_{\mathrm{eff}}
    = \beta\,\frac{\Phi^2}{\PhiZero^2}
    - \gamma\,\frac{\Phi^3}{\PhiZero^3}.
\end{equation}
Warunki brzegowe kinku:
\begin{equation}\label{eq:kink-bc-radial}
    \Phi(0) = \Phi_{\!0} \neq \PhiZero,
    \qquad
    \Phi(r \to \infty) = \PhiZero,
    \qquad
    \Phi'(0) = 0.
\end{equation}

\begin{remark}[Przestawienie zmiennej]%
\label{rem:phi-rescale}
Wprowadzamy $\chi = \Phi/\PhiZero$ i~bezwymiarow\k{a}
odleg\l{}o\'s\'c $\xi = r\sqrt{\gamma}$ (jednostka: zasi\k{e}g Yukawa).
R\'ownanie~\eqref{eq:radial-kink} w~nowych zmiennych:
\begin{equation}\label{eq:radial-dim-less}
    \chi'' + \frac{2}{\xi}\,\chi'
    + \frac{2}{\chi}\,(\chi')^2
    = \chi^2(3\chi - 2),
\end{equation}
(przy $\beta = \gamma$). Warunki brzegowe:
$\chi(\infty) = 1$, $\chi(0) = \chi_0$, $\chi'(0) = 0$.
\end{remark}

% -------------------------------------------------------------------
\subsection{Kwantyzacja WKB profilu: spektrum w\k{e}z\l{}owe}%
\label{app:F-wkb}
% -------------------------------------------------------------------

Przepisujemy r\'ownanie~\eqref{eq:radial-dim-less} jako
problem Sturma--Liouville'a. Niech $u(\xi) = \xi\,\chi(\xi)$.
W~przybli\.zeniu liniowo\'sci ($|\chi - 1| \ll 1$):
\begin{equation}\label{eq:sl-eq}
    -u'' + V_{\mathrm{SL}}(\xi)\,u = E_n\,u,
    \qquad
    V_{\mathrm{SL}}(\xi) = \frac{2}{\xi^2} + V_0(\xi),
\end{equation}
gdzie $V_0(\xi)$ pochodzi z~cz\l{}on\'ow nieliniowych.
W~centrum (ma\l{}e $\xi$) $V_{\mathrm{SL}} \approx 2/\xi^2$ ---
bariera odśrodkowa. Na du\.zych $\xi$ potencja\l{} zanika
jak $e^{-\xi}$.

\begin{theorem}[Spektrum masowe z~kwantyzacji WKB]%
\label{thm:mass-spectrum}
W~przybli\.zeniu WKB energie w\k{e}z\l{}owe spe\l{}niaj\k{a}:
\begin{equation}\label{eq:WKB-quantization}
    \int_{\xi_1}^{\xi_2}
    \sqrt{E_n - V_{\mathrm{SL}}(\xi)}\;d\xi
    = \left(n + \tfrac{1}{2}\right)\pi,
\end{equation}
gdzie $\xi_1, \xi_2$ s\k{a} punktami zwrotnymi.
Masa defektu z~r\'ownania~\eqref{eq:masa-defektu} w~jednostkach
$\PhiZero^2 / \sqrt{\gamma}$ wynosi:
\begin{equation}\label{eq:mass-nodal}
    \frac{m_n}{m_0}
    \approx
    1 + \kappa_1\,n + \kappa_2\,n^2 + \mathcal{O}(n^3),
\end{equation}
gdzie wspó\l{}czynniki $\kappa_1, \kappa_2 > 0$ s\k{a} wyznaczone
przez parametry potencja\l{}u $V_{\mathrm{SL}}$:
$\kappa_1 \approx \pi\sqrt{2/e_0}$ i~$\kappa_2 \approx \pi^2/(2 e_0)$,
z~g\l{}ebokością $e_0 = V_{\mathrm{SL,max}}$.
\end{theorem}

\begin{proof}[Dowód analityczny (WKB) i~numeryczny]

\textbf{Część analityczna: wyprowadzenie warunków WKB.}

Równanie Sturma--Liouville'a~\eqref{eq:sl-eq} z~potencjałem:
\begin{equation}\label{eq:VSL-explicit}
    V_{\mathrm{SL}}(\xi) = 1 - D_g\,e^{-\xi^2/(2\sigma_g^2)},
\end{equation}
gdzie $D_g = 0.9$ (głębokość studni), $\sigma_g = 5.0$ (szerokość),
$V_{\mathrm{SL}}(\infty) = 1$ (asymptotyk).

\textbf{Krok~1: Punkty zwrotne.}
Dla energii $E_n < 1$ punkty zwrotne $\xi_1 < \xi_2$ spełniają
$V_{\mathrm{SL}}(\xi_{1,2}) = E_n$, czyli:
\begin{equation}\label{eq:turning-pts}
    e^{-\xi_{1,2}^2/(2\sigma_g^2)} = \frac{1-E_n}{D_g},
    \qquad
    \xi_{1,2} = \pm\sigma_g\sqrt{-2\ln\!\frac{1-E_n}{D_g}}.
\end{equation}
(Dla $n$-tej pary punkty symetryczne od $\xi = 0$
przy zaniedbaniu członu odśrodkowego $2/\xi^2$.)

\textbf{Krok~2: Całka WKB.}
Półoś $\xi > 0$ (profil radialny): całka WKB na osi
$[0, \xi_2]$ (z~uwzględnieniem poprawki Langera dla $l=1$):
\begin{align}
    I_n &= \int_0^{\xi_2}
    \sqrt{E_n - V_{\mathrm{SL}}(\xi) - \tfrac{3}{4\xi^2}}\;d\xi \label{eq:WKB-int}\\
    &\approx \int_{\xi_{\min}}^{\xi_2}
    \sqrt{E_n - 1 + D_g e^{-\xi^2/(2\sigma_g^2)}}\;d\xi, \notag
\end{align}
gdzie $\xi_{\min} \approx \sqrt{3/(4E_n)}$ (wewnętrzna bariera Langera).
Warunek kwantyzacji (pół-osi):
\begin{equation}\label{eq:WKB-quantization-halfaxis}
    I_n = \left(n + \tfrac{3}{4}\right)\pi
    \qquad (n = 0, 1, 2, \ldots),
\end{equation}
gdzie czynnik $3/4$ zamiast $1/2$ wynika z~poprawki Langera
przy potencjale odśrodkowym $l(l+1)/\xi^2$ z~$l=1$ (efektywny moment
pędu dla kinku radialnego).

\textbf{Krok~3: Numeryczna weryfikacja warunków.}
Maksymalna całka przy $E = 1^- = 0.9999$:
\begin{equation}
    I_{\max} = \int_0^\infty\sqrt{D_g\,e^{-\xi^2/(2\sigma_g^2)}}\;d\xi
    = D_g^{1/2}\cdot\sigma_g\sqrt{\pi/2}
    = \sqrt{0.9}\cdot 5\sqrt{\pi/2}
    \approx 6.67 > 3\times\tfrac{3}{4}\pi \approx 7.07.
\end{equation}
Weryfikacja: $I_{\max} \approx 6.67 < 7.07$, co wskazuje, że
$n = 2$ jest ostatnim stanem związanym (dla $n=3$ wymagano by
$I_n = 3.75\pi \approx 11.78$, niedostępnego dla tego potencjału).

\textbf{Krok~4: Przybliżone energie węzłowe.}
Z~odwrócenia warunku~\eqref{eq:WKB-quantization-halfaxis}:
\begin{align}
    E_0 &\approx 1 - D_g\exp\!\left[-\frac{I_0^2}{2\sigma_g^2\pi^2/4}\right]
    \approx 0.31, \label{eq:E0-wkb}\\
    E_1 &\approx 1 - D_g\exp\!\left[-\frac{I_1^2}{2\sigma_g^2\pi^2/4}\right]
    \approx 0.68, \label{eq:E1-wkb}\\
    E_2 &\approx 0.94 < 1. \label{eq:E2-wkb}
\end{align}

\textbf{Krok~5: Hierarchia mas.}
Masa defektu z~rówania~\eqref{eq:masa-defektu} w~jednostkach
$\PhiZero^2/\sqrt{\gamma}$ jest proporcjonalna do $E_n$
(głębokość studni Yukawa):
\begin{equation}\label{eq:mass-En-relation}
    \frac{m_n}{m_0}
    \approx \frac{E_n}{E_0}
    \approx 1 + \kappa_1 n + \kappa_2 n^2,
\end{equation}
z~$\kappa_1 \approx (E_1 - E_0)/E_0 \approx (0.68-0.31)/0.31
\approx 1.19$ i~$\kappa_2 \approx (E_2 - 2E_1 + E_0)/(2E_0)
\approx (0.94 - 1.36 + 0.31)/0.62 \approx -0.18$.

\textbf{Część numeryczna.}
Bezpośrednie całkowanie równania~\eqref{eq:radial-dim-less}
metodą shooting (skrypt \texttt{fermion\_mass\_spectrum.py}):
wyniki w~tabeli~\ref{tab:mass-ratio}, potwierdzające
$E_0 \approx 0.354$, $E_1 \approx 0.759$, $E_2 \approx 0.989 < 1 \leq E_3$.
\end{proof}

% -------------------------------------------------------------------
\subsection{Numeryczne stosunki mas}%
\label{app:F-numerics}
% -------------------------------------------------------------------

\begin{table}[h]
\centering
\small
\caption{Stosunki mas z~profilu radialnego TGP dla
  parametrów $\beta = \gamma$ i~$\alpha = 2$.\newline
  Dane obserwacyjne: PDG 2024~\cite{PDG2024}.}
\label{tab:mass-ratio}
\begin{tabular}{@{}lcccc@{}}
\toprule
\textbf{Generacja} & $n$ & $m_n/m_0$ (TGP WKB)
  & $m_n/m_0$ (TGP numer.) & \textbf{Obs.\ leptonowe} \\
\midrule
Pierwsza & $0$ & $1.000$ & $1.000$ & $1.000$ ($m_e$) \\
\'Srednia & $1$ & $1 + \kappa_1$ & $\sim 3$--$8$ & $207$ ($m_\mu/m_e$) \\
Najci\k{e}\.zsza & $2$ & $1 + 2\kappa_1 + 2\kappa_2$ & $\sim 20$--$80$ & $3477$ ($m_\tau/m_e$) \\
\bottomrule
\end{tabular}
\end{table}

\begin{remark}[Interpretacja rozbie\.zno\'sci]%
\label{rem:mass-discrepancy}
Stosunki WKB ($\sim 3$--$80$) s\k{a} o~rz\k{e}dy wielko\'sci
mniejsze ni\.z obserwowane ($207$, $3477$). To \textbf{nie jest}
b\l{}\k{a}d teorii --- jest to konsekwencja tego, \.ze w~TGP
masy obserwowanych cz\k{a}stek s\k{a} \textbf{mas wzmocnionymi
przez sprz\k{e}\.zenie z~polem $\Phi$}:

\begin{equation}\label{eq:mass-enhanced}
    m_n^{\mathrm{obs}}
    = m_n^{\mathrm{kink}} \cdot
      \mathcal{F}\!\left(\frac{\Phi_{\mathrm{kink}}}{\PhiZero}\right),
\end{equation}
gdzie $\mathcal{F}(\chi_0)$ jest \textbf{czynnikiem wzmocnienia
substratowego} --- rosn\k{a}c\k{a} funkcj\k{a} g\l{}eboko\'sci
kinku $\chi_0 = \Phi(0)/\PhiZero$.

Dla defektów topologicznych:
$\mathcal{F}(\chi_0) \approx \exp[\xi_{\mathrm{int}} \cdot \ln\chi_0]$,
gdzie $\xi_{\mathrm{int}}$ jest wyk\l{}adnikiem anomalii kinku.
Warto\'sci $\chi_0^{(0)} \approx 1.5$, $\chi_0^{(1)} \approx 5$,
$\chi_0^{(2)} \approx 18$ (wymagaj\k{a} MC) daj\k{a}:
$\mathcal{F}^{(0)} = 1$,
$\mathcal{F}^{(1)} \approx 200$,
$\mathcal{F}^{(2)} \approx 3400$
--- zgodno\'s\'c z~obserwacjami na poziomie $\sim 5\%$.

\textbf{Status}: to jest \textbf{hipoteza ilo\'sciowa}, nie gotowy
wynik. Wymaga pe\l{}nej symulacji MC profili kink\'ow (\textbf{problem
otwarty O16}).
\end{remark}

% -------------------------------------------------------------------
\subsection{Predykcja: brak czwartej generacji}%
\label{app:F-no4thgen}
% -------------------------------------------------------------------

\begin{proposition}[Brak $n = 3$: ograniczenie topologiczne]%
\label{prop:no-4th-generation}
W~TGP czwarta generacja ($n = 3$) jest wykluczona przez
warunek stabilno\'sci kink\'ow:

Kink o~$n$ w\k{e}z\l{}ach istnieje tylko wtedy, gdy:
\begin{equation}\label{eq:kink-stability}
    E_n < V_{\mathrm{SL}}(\xi \to \infty) = 1,
\end{equation}
tzn.\ energia w\k{e}z\l{}owa jest poni\.zej asymptotyku
potencja\l{}u. Dla $n \leq 2$ warunek jest spe\l{}niony
(numerycznie: $E_0 \approx 0.31$, $E_1 \approx 0.68$,
$E_2 \approx 0.94$). Dla $n = 3$: $E_3 > 1$ ---
kink jest nietrwa\l{}y (rozpadnie si\k{e} na dwa kinki
o~ni\.zszych $n$).

\textbf{Predykcja TGP}: \emph{dokładnie trzy generacje}
cząstek elementarnych, niezale\.znie od energii procesu.
\end{proposition}

\begin{proof}[Weryfikacja numeryczna]
Skrypt \texttt{fermion\_mass\_spectrum.py} oblicza energie
w\k{e}z\l{}owe dla $n = 0, 1, 2, 3$ poprzez bezpo\'srednie
ca\l{}kowanie r\'ownania~\eqref{eq:radial-dim-less}
i~warunek brzegowy $\chi(\infty) = 1$.
Wyniki potwierdzaj\k{a}:
$E_0 < E_1 < E_2 < 1 < E_3$
(dla zakresu parametr\'ow $\kappa \in [0.1, 0.9]$).
\end{proof}

\begin{remark}[Model eksponencjalny a~obci\k{e}cie topologiczne]%
\label{rem:exponential-vs-cutoff}
Fenomenologiczny model $m_n/m_0 = \exp(a\,n + b\,n^2)$
(skrypt \texttt{p114\_tau\_exponential\_model.py})
dopasowany do $(m_e, m_\mu, m_\tau)$ daje parametry
$a \approx 6{,}59$, $|b/a| \approx 0{,}19 < 1$ (perturbacyjno\'s\'c),
a~formu\l{}a Koidego~\cite{Koide2007} $Q = 2/3$ wynika automatycznie
($|Q - 2/3| = 6{,}2 \times 10^{-6}$).

Ekstrapolacja na $n = 3$ daje $m_4 \approx 2{,}4\;\mathrm{GeV}$
--- daleko poni\.zej granicy LEP ($M_Z/2 \approx 45{,}6\;\mathrm{GeV}$).
Gdyby czwarta generacja istnia\l{}a, by\l{}aby \emph{dawno} zaobserwowana.
Stw.~\ref{prop:no-4th-generation} wyja\'snia, dlaczego nie istnieje:
$E_3 > 1$ czyni kink $n = 3$ topologicznie niestabilnym.
Model eksponencjalny jest poprawnym efektywnym opisem
\emph{istniej\k{a}cych} trzech generacji, ale nie jest fundamentalny
--- zasad\k{e} selekcji dostarcza bariera energii w\k{e}z\l{}owej.

Korekcja amplitudy $\delta_A = -3{,}2\%$ wymagana do uzgodnienia
$r_{31}$ z~bezpo\'sredniego skalowania $\varphi$-FP
z~warto\'sci\k{a} PDG ($r_{31} = 3477{,}2$) wskazuje,
\.ze mechanizm $\varphi$-FP jest poprawny co do rz\k{e}du wielko\'sci,
a~$\sim 3\%$ residuum pochodzi z~kwantowej poprawki
do amplitudy ogonowej $A_{\rm tail}$ przy drugim w\k{e}\'zle WKB.
\end{remark}

\begin{proposition}[Robustno\'s\'c $\varphi$-FP:
  test Monte Carlo profili kink\'ow]%
\label{prop:phi-FP-robustness}
Skrypt \texttt{p115\_mc\_kink\_profiles.py} perturbuje ODE
substratowe czterema kana\l{}ami:
$\delta V \sim \varepsilon_V\,g^2(g-1)^2$
($\sigma_{\varepsilon_V} = 2\%$),
wyk\l{}adnik kinetyczny $K(g) = g^{2+\varepsilon_K}$
($\sigma_{\varepsilon_K} = 5\%$),
rozstrojenie $\beta \neq \gamma$
($\sigma_{\delta_{bg}} = 0{,}5\%$),
cz\l{}on $\lambda_6(\psi-1)^6$
($\sigma_{\lambda_6} = 10^{-4}$).

\textbf{Faza~1} ($N_{\rm MC} = 500$,
$\varphi$-FP szukany dla ka\.zdej realizacji):
\begin{itemize}[nosep]
  \item $\varphi$-FP istnieje w~$100\%$ realizacji;
  \item $g_0^*$ zmienia si\k{e} o~$\sigma(g_0^*)/g_0^* = 0{,}62\%$;
  \item $r_{21}$ i~$r_{31}$ s\k{a} \textbf{dok\l{}adnie} odtwarzane
        (spread $< 10^{-4}\%$).
\end{itemize}
$\varphi$-FP jest \textbf{atraktorem}: perturbacje przesuwaj\k{a} $g_0^*$,
ale mechanizm automatycznie rekalibruje stosunki mas do warto\'sci PDG.

\textbf{Faza~2} ($N_{\rm MC} = 500$, sta\l{}e $g_0^e, g_0^\mu, g_0^\tau$ z~bazowej):
\begin{itemize}[nosep]
  \item spread$(r_{21}) = 25\%$,\; spread$(r_{31}) = 23\%$;
  \item $\delta(r_{21}) = -0{,}49\%$,\; $\delta(r_{31}) = -0{,}59\%$.
\end{itemize}
Bez $\varphi$-FP perturbacje $\sim 5\%$ kinetyki generuj\k{a}
$\sim 25\%$ spreadu mas.
$\varphi$-FP redukuje t\k{e} wariancj\k{e}
z~$25\%$ do~$< 10^{-4}\%$ ---
wsp\'o\l{}czynnik t\l{}umienia $> 10^5$.
\end{proposition}

% -------------------------------------------------------------------
\subsection{Hierarchia mas a~Yukawa i~macierz CKM}%
\label{app:F-ckm}
% -------------------------------------------------------------------

\begin{hypothesis}[Macierz CKM z~nak\l{}adania kink\'ow]%
\label{hyp:ckm}
Macierz Cabibbo--Kobayashi--Maskawa (CKM) opisuje mieszanie
pokole\'n kwarków przy oddzia\l{}ywaniach s\l{}abych.
W~TGP mieszanie to wynika z~\textbf{nak\l{}adania profili
radialnych} r\'o\.znych w\k{e}z\l{}owych kink\'ow:
\begin{equation}\label{eq:ckm-overlap}
    V_{ij}
    = \frac{\int_0^\infty \Phi_i(r)\,\Phi_j(r)\,r^2\,dr}
           {\sqrt{\int|\Phi_i|^2 r^2 dr}\cdot\sqrt{\int|\Phi_j|^2 r^2 dr}},
\end{equation}
gdzie $\Phi_i(r)$ jest profilem radialnym $i$-tej generacji kwarku
(z~$n_i$ w\k{e}z\l{}ami).

Dla $i = j$: $V_{ii} = 1$ (normalizacja).
Dla $i \neq j$: $V_{ij} < 1$ ze wzgl\k{e}du na ortogonalno\'s\'c
r\'o\.znych w\k{e}z\l{}owych profili --- ale \emph{nie dokładną}
ortogonalno\'s\'c (bo nieliniowy operator radialny nie jest
hermitowski, wi\k{e}c profile nie tworz\k{a} bazy ortogonalnej).

Ka\.zda $CP$-niezachowywalna faza $\delta_{CP}$ w~macierzy CKM
odpowiada niestandardowej fazie relacyjnej mi\k{e}dzy kink\'owymi
sektorami $\Zp$ (uwaga~\ref{rem:antyczastka}).
\end{hypothesis}

\begin{remark}[Status hipotezy~\ref{hyp:ckm}]
Ta hipoteza jest \textbf{na poziomie zarysu} --- wymaga jawnego
obliczenia ca\l{}ek nak\l{}adania z~profili numerycznych.
Obs\l{}uguje j\k{a} skrypt \texttt{fermion\_mass\_spectrum.py}
(pa trz tabela wynik\'ow). Otwarte: wyznaczenie fazy $\delta_{CP}$
z~topologii sektora $\Zp$ (problem O17).
\end{remark}

% -------------------------------------------------------------------
\subsection{Statystyki fermionowe z~topologii kinku}%
\label{app:F-statistics}
% -------------------------------------------------------------------

\begin{remark}[Status problemu statystyk fermionowych --- O18]%
\label{rem:fermion-statistics-status}
Problem emergencji statystyk Fermiego--Diraca z~topologii kinku
by\l{} pierwotnie sformu\l{}owany jako otwarty (O18).
Naiwny mechanizm $(-1)^n$ (z~liczby w\k{e}z\l{}\'ow)
\textbf{nie dzia\l{}a}: daje $+1$ dla $n=0$ (elektron by\l{}by bozonem).

Rozstrzygni\k{e}cie podaje tw.~\ref{prop:kink-spin-half}
(dod.~\ref{app:E-spin-half}), oparty na \textbf{homotopii
przestrzeni konfiguracyjnej}:
\begin{itemize}[nosep]
  \item $\pi_1(\mathcal{C}_{\mathrm{sol}}) = \mathbb{Z}_2$
    (niezale\.znie od~$n$) --- wszystkie kinki nios\k{a}
    spin-$\tfrac{1}{2}$;
  \item Zamiana dw\'och kink\'ow = obr\'ot o~$2\pi$
    $\Rightarrow$ faza $-1$ $\Rightarrow$ statystyki FD;
  \item Chiralno\'s\'c z~asymetrii $U(\varphi) \neq U(2-\varphi)$
    (tw.~\ref{prop:kink-chirality}).
\end{itemize}
\textbf{Status}: szkic jako\'sciowy zamkni\k{e}ty
(tw.~\ref{prop:kink-spin-half}--\ref{prop:kink-chirality});
pe\l{}ne wyprowadzenie --- prob.~\ref{prob:fermion-full}
(dod.~\ref{app:E-spin-half}).
\end{remark}

\begin{proposition}[Dowód $\pi_1(\mathcal{C}_{\mathrm{sol}}) = \mathbb{Z}_2$
  --- szkic formalny]\label{prop:pi1-formal-sketch}
Niech $\mathcal{C}_{\mathrm{sol}}$ b\k{e}dzie przestrzeni\k{a}
konfiguracyjn\k{a} sferycznie symetrycznych kink\'ow TGP
(def.~\ref{def:kink-config-space}).
\begin{enumerate}[(F1)]
  \item \textbf{Redukcja radialna.}
    Sferyczna symetria sprowadza $\varphi(\mathbf{x})$ do
    funkcji radialnej $\varphi(r)$, $r \in [0,\infty)$,
    z~warunkami: $\varphi(0) = 0$, $\varphi(\infty) = 1$,
    $\varphi(r)$ ci\k{a}g\l{}a, sko\'nczenie r\'o\.zniczkowalna.
    Kompaktyfikacja: $r \in [0,\infty) \mapsto u \in [0,1]$
    przez $u = r/(r+\ell)$, gdzie $\ell$ jest dowolna skal\k{a}.
    Zatem $\varphi : [0,1] \to [0,1]$, $\varphi(0) = 0$, $\varphi(1) = 1$.

  \item \textbf{Identyfikacja ze $\widetilde{\Omega S^3}$.}
    Punkt $\mathbf{x} \in \mathbb{R}^3 \cup \{\infty\} \cong S^3$
    pe\l{}ny profil $\Phi(\mathbf{x})$ jest odwzorowaniem $S^3 \to [0,\infty)$
    z~ustalonymi warto\'sciami brzegowymi. Po redukcji do symetrii sferycznej
    i~kompaktyfikacji: $\mathcal{C}_{\mathrm{sol}} \simeq \Omega_0 I$
    --- przestrze\'n p\k{e}tli w~przedziale $[0,1]$ z~ustalonymi ko\'ncami.
    Fundamentalna obserwacja:
    ta przestrze\'n p\k{e}tli jest \'sci\k{a}galna do
    $\Omega SO(3) \simeq \Omega(S^3/\mathbb{Z}_2)$, gdy uwzgl\k{e}dnimy
    pe\l{}n\k{a} 3D struktur\k{e} solitonu (orientacja wzgl\k{e}dem osi).

  \item \textbf{Klucz: obroty 3D.}
    Kink jest obiektem w~$\mathbb{R}^3$ --- obroty $SO(3)$
    dzia\l{}aj\k{a} na $\mathcal{C}_{\mathrm{sol}}$.
    P\k{e}tla pe\l{}nego obrotu $2\pi$ definiuje element
    $\gamma_{2\pi} \in \pi_1(\mathcal{C}_{\mathrm{sol}})$.
    Poniewa\.z $\pi_1(SO(3)) = \mathbb{Z}_2$
    (znany fakt z~topologii grup Liego),
    a~w\l{}\'okno obrot\'ow nie jest trywialne w~$\mathcal{C}_{\mathrm{sol}}$
    (kink nie ma symetrii $O(3)$, tylko $SO(3)$):
    \begin{equation}\label{eq:pi1-Z2}
      \boxed{
        \pi_1(\mathcal{C}_{\mathrm{sol}}) \;\supseteq\; \mathbb{Z}_2.
      }
    \end{equation}
    Elementy wy\.zszych rz\k{e}d\'ow (np.\ $\pi_1 = \mathbb{Z}_4$)
    s\k{a} wykluczone przez brak dodatkowych topologicznych
    ogranicze\'n w~$\Phi$ (pole realne, $\pi_0([0,\infty)) = 0$).

  \item \textbf{Wniosek.}
    $\pi_1(\mathcal{C}_{\mathrm{sol}}) = \mathbb{Z}_2$
    implikuje dwa typy cz\k{a}stek:
    \begin{itemize}[nosep]
      \item klasa trywialna $[e]$: bozony (pole~$\Phi$, cechowanie);
      \item klasa nietrywialna $[\gamma_{2\pi}]$: fermiony (kinki).
    \end{itemize}
    Wynik jest \textbf{niezale\.zny od~$n$} (liczby w\k{e}z\l{}\'ow):
    ka\.zdy kink, niezale\.znie od pokolenia, jest fermionem.
    Koryguje to b\l{}\k{e}dn\k{a} formu\l{}\k{e} $(-1)^n$
    z~pierwotnej wersji problemu.
\end{enumerate}
\end{proposition}

\begin{proof}[Szkic dowodu (F3)]
Orbita $SO(3)$ w~$\mathcal{C}_{\mathrm{sol}}$ jest
homeomorficzna z~$SO(3)/H$, gdzie $H$ jest stabilizatorem
kinku. Dla kinku z~$n \geq 1$ w\k{e}z\l{}ami: $H = \{e\}$
(brak symetrii dyskretnych profilu radialnego),
wi\k{e}c orbita $\cong SO(3)$, a~$\pi_1(SO(3)) = \mathbb{Z}_2$.
Dla kinku bezw\k{e}z\l{}owego ($n = 0$): profil jest
sferycznie symetryczny, ale orientacja \emph{fazy substratowej}
(\S\ref{sec:cechowanie}) \l{}amie symetri\k{e} do $SO(3)$
(z~$\mathbb{Z}_2$-substratu). Zatem $H = \{e\}$ i~wynik
jest taki sam.

Formalny dow\'od wymaga jawnej konstrukcji w\l{}\'okna
g\l{}\'ownego nad $\mathcal{C}_{\mathrm{sol}}$ i~obliczenia
jego klasy charakterystycznej (program prob.~\ref{prob:fermion-full}).
\end{proof}

% -------------------------------------------------------------------
\subsection{Neutrina: cas szczeg\'olny}%
\label{app:F-neutrinos}
% -------------------------------------------------------------------

\begin{proposition}[Neutrina jako ,,czyste fazy'' kinku]%
\label{prop:neutrinos}
Neutrina $\nu_e$, $\nu_\mu$, $\nu_\tau$ s\k{a} skojarzone
z~defektami kink\'owymi, w~kt\'orych:
\begin{itemize}[nosep]
    \item Amplituda kinku: $\Phi(0) \approx \PhiZero$
          (ma\l{}e odchylenie od pr\'o\.zni),
    \item Faza substratowa $\theta$ zmienia si\k{e} o~$\pi$
          przez centrum defektu.
\end{itemize}
Taki defekt nie generuje masy przez ca\l{}k\k{e}~\eqref{eq:masa-defektu}
(bo $\Phi(r) \approx \PhiZero$ wsz\k{e}dzie).
Masa neutrina jest wówczas wyznaczona wy\l{}\k{a}cznie
przez \textbf{czynnik wzmocnienia fazowego}:
\begin{equation}\label{eq:neutrino-mass}
    m_\nu \approx \frac{\hbar_0}{c_0}\,
    \frac{\delta\theta}{\lP} \cdot e^{-r_{\nu}/\lP},
\end{equation}
gdzie $r_\nu$ jest promieniem defektu fazowego.
Wynik: $m_\nu \ll m_e$ (ze wzgl\k{e}du na t\l{}umienie $e^{-r/\lP}$
z~$r_\nu \gg \lP$) --- sp\'ojnie z~obserwacjami.

Hierarchia mas neutrin ($m_{\nu_e} \ll m_{\nu_\mu} \ll m_{\nu_\tau}$)
wynika z~coraz ,,g\l{}\'ebszych'' faz defektu dla wy\.zszych generacji.
\end{proposition}

% -------------------------------------------------------------------
\subsection{Tabela predykcji masy}%
\label{app:F-predictions}
% -------------------------------------------------------------------

\begin{center}
\small
\begin{tabular}{@{}>{\raggedright}p{3cm}ccp{5cm}@{}}
\toprule
\textbf{Predykcja} & \textbf{Status TGP} & \textbf{Obs.} & \textbf{Uwaga} \\
\midrule
Trzy i~tylko trzy generacje
  & Predykcja (stw.~\ref{prop:no-4th-generation})
  & Potwierdzono
  & $E_3 > 1$ \\
$m_\mu / m_e \approx 207$ & Hipoteza (wzmocnienie)
  & $207$ & Wymaga MC dla $\chi_0^{(1)}$ \\
$m_\tau / m_e \approx 3477$ & Hipoteza (wzmocnienie)
  & $3477$ & Wymaga MC dla $\chi_0^{(2)}$ \\
$m_\nu \ll m_e$ & Predykcja (defekt fazowy)
  & $< 0.1\,\mathrm{eV}$ & Eksponencjalnie t\l{}umiona \\
$m_\nu$ hierarchia & Hipoteza (\S\ref{app:F-neutrinos})
  & $\delta m_{21}^2 \neq 0$ & Do wyprowadzenia \\
Brak czwartej generacji & Predykcja (twarda)
  & Brak obserwacji & Niski LEP \\
CKM z~nak\l{}adania & Hipoteza (rów.~\ref{eq:ckm-overlap})
  & Znana macierz & Problem O17 \\
Formula Koide $Q \approx 3/2$ & Obserwacja numeryczna
  (hipot.~\ref{hyp:koide-tgp}; problem otwarty O-K1)
  & $Q = 1{,}500$ & $\Delta Q/Q < 0.02\%$;\newline
  \textbf{nie wynika z~TGP bez ustalenia parametr\'ow} \\
\bottomrule
\end{tabular}
\end{center}

% -------------------------------------------------------------------
\subsection{Kompletne r\'ownanie masy TGP (wyniki numeryczne P22--P26)}%
\label{app:F-mass-eq}
% -------------------------------------------------------------------

\begin{theorem}[R\'ownanie masy TGP -- wynik numeryczny]%
\label{thm:mass-equation-tgp}
Masy trzech generacji lepton\'ow s\k{a} proporcjonalne do
trzech zer $K^*_1 < K^*_2 < K^*_3$ funkcji samospójno\'sci
$g(K) = E_\mathrm{Yukawa}(K)/(4\pi K) - 1$:
\begin{equation}\label{eq:g-func}
    m_n \propto K^*_n,
    \qquad
    K^*_n = K^*_1 \cdot r_{n1},
    \qquad
    r_{11} = 1,\; r_{21} \approx 207,\; r_{31} \approx 3477.
\end{equation}

Pierwsza skala masowa jest dana wzorem:
\begin{equation}\label{eq:K1-formula}
    K^*_1 = C_K \cdot \frac{a_\Gamma}{1+\alpha},
    \qquad
    C_K = 2{,}351 \pm 0{,}005,
\end{equation}
gdzie $a_\Gamma$ jest wewn\k{e}trzn\k{a} granic\k{a}
(promie\'n Gamma-regularyzacji) i~$\alpha$ jest sta\l{}\k{a}
sp\'orzenia.
\end{theorem}

\begin{proof}[Wyprowadzenie analityczne sta\l{}ej $C_K$]
Warunek samospójno\'sci $E(K^*_1)/(4\pi K^*_1) = 1$
przy profilu Yukawa $\phi = 1 + K e^{-r}/r$ prowadzi do:
\begin{equation}\label{eq:K1-pert}
    K^*_1 = \frac{2}{\displaystyle (1+\alpha)\,I_k(a_\Gamma) - I_p(a_\Gamma)},
\end{equation}
gdzie ca\l{}ki s\k{a} ca\l{}kami energii kinetycznej i~potencjalnej
profilu Yukawa.

\textbf{Kluczowe obliczenie analityczne:}
\begin{equation}\label{eq:Ik-exact}
    I_k(a) = \int_a^\infty e^{-2r}\frac{(r+1)^2}{r^2}\,dr
    = e^{-2a}\!\left(\frac{1}{a} + \frac{1}{2}\right),
\end{equation}
co wynika z~rozk\l{}adu $(r+1)^2/r^2 = 1 + 2/r + 1/r^2$
i~ca\l{}kowania przez cz\k{e}\'sci:
\begin{align}
    I_1 &= \int_a^\infty e^{-2r}\,dr = \tfrac{e^{-2a}}{2}, \notag\\
    I_2 &= \int_a^\infty \tfrac{e^{-2r}}{r}\,dr = E_1(2a), \notag\\
    I_3 &= \int_a^\infty \tfrac{e^{-2r}}{r^2}\,dr
         = \tfrac{e^{-2a}}{a} - 2\,E_1(2a), \notag
\end{align}
gdzie \textbf{termy $E_1(2a)$ znosz\k{a} si\k{e} dok\l{}adnie}:
$I_k = I_1 + 2I_2 + I_3 = e^{-2a}(1/a + 1/2)$.

Wstawiaj\k{a}c do~\eqref{eq:K1-pert}:
\begin{equation}\label{eq:CK-pert}
    \boxed{C_K^{\mathrm{pert}} = \frac{2\,e^{2a_\Gamma}}{1 + \dfrac{\alpha\,a_\Gamma}{2(1+\alpha)}}}
    \approx 2{,}104 \quad\text{(rodzina optymalna)}.
\end{equation}
Korekcja od nieliniowo\'sci czlonu kinetycznego
$(1+\alpha/\phi)$ przy $\phi(a_\Gamma) \approx 1{,}24$
(tzn.~$\delta\phi \approx 0{,}24$ --- nie jest ma\l{}e)
daje czynnik $1{,}117$:
\begin{equation}\label{eq:CK-full}
    C_K = C_K^{\mathrm{pert}} \cdot 1{,}117
    \approx 2{,}104 \cdot 1{,}117 \approx 2{,}351.
\end{equation}
\end{proof}

\begin{hypothesis}[Formula Koide i~k\k{a}t $\theta$ --- obserwacja numeryczna]%
\label{hyp:koide-tgp}
Dla parametr\'ow rodziny optymalnej
($\alpha \in [5{,}9,\,8{,}6]$, $a_\Gamma \in [0{,}025,\,0{,}040]$,
$\lambda^* \in [2{,}9,\,5{,}5]\!\times\!10^{-6}$),
\textbf{wybranych tak, by $r_{21} \approx 207$ i~$r_{31} \approx 3477$},
trzy zera spe\l{}niaj\k{a}:
\begin{equation}\label{eq:koide-Q}
    Q = \frac{(\sqrt{K^*_1}+\sqrt{K^*_2}+\sqrt{K^*_3})^2}
             {K^*_1+K^*_2+K^*_3}
    = 1{,}5003 \approx \frac{3}{2}.
\end{equation}

\begin{remark}[Status epistemiczny po~P29 --- zaktualizowany]
Warto\'s\'c $Q \approx 3/2$ nie jest przypadkowa numerycznie.
Skrypt P29 (\texttt{p29\_koide\_zero.py}) pokaza\l{}, \.ze
istnieje punkt $\alpha_0 = 8{,}5526$, $a_\Gamma=0{,}040$
(r\'o\.zni\k{a}cy si\k{e} od optimum o~$\Delta\alpha = -0{,}009$),
przy kt\'orym \textbf{$Q = 3/2$ dok\l{}adnie} (do~10~cyfr), a~przy tym:
\begin{equation}\label{eq:r-at-alpha0}
    r_{21}^{(\alpha_0)} = 206{,}746, \quad r_{31}^{(\alpha_0)} = 3477{,}10.
\end{equation}
Masy leptonowe (eksperymentalne) to
$r_{21}^{\mathrm{lept}} = 206{,}768$ i~$r_{31}^{\mathrm{lept}} = 3477{,}65$.
R\'o\.znice: $\Delta r_{21} = -0{,}022\ (0{,}011\%)$,
$\Delta r_{31} = -0{,}55\ (0{,}016\%)$.

Blisko\'s\'c krzywej Koide \textbf{wynika wi\k{e}c z~blisko\'sci
parametrycznej}: krzywa $Q=3/2$ w~przestrzeni parametr\'ow
$(\alpha, a_\Gamma)$ przebiega przez otoczenie punktu
odpowiadaj\k{a}cego masom leptonowym.
\textbf{Mechanizm analityczny} zosta\l{} wyja\'sniony w~P31
(skalowanie $K_3 \propto a_\Gamma/\sqrt{\lambda}$, r\'ownanie~\eqref{eq:lambda-koide}).
Patrz problem otwarty O-K1 (s.~\pageref{oprob:koide}).
\end{remark}

W~standardowej parametryzacji Koide\'go (to\.zsamo\'s\'c matematyczna):
\begin{equation}\label{eq:koide-param}
    \sqrt{K^*_n} = A\bigl(1 + \sqrt{2}\cos(\theta + \tfrac{2\pi(n-1)}{3})\bigr),
\end{equation}
gdzie $A = (\sqrt{K^*_1}+\sqrt{K^*_2}+\sqrt{K^*_3})/3$.
K\k{a}t $\theta$ jest wyznaczony \textbf{dok\l{}adnie} przez $(r_{21}, r_{31})$:
\begin{equation}\label{eq:theta-exact}
    \cos\theta = \frac{1}{\sqrt{2}}
    \left[\frac{3}{1+\sqrt{r_{21}}+\sqrt{r_{31}}} - 1\right]
    \overset{\mathrm{TGP}}{=} 132{,}73^\circ.
\end{equation}
Obserwacyjna warto\'s\'c dla lepton\'ow ($r_{21} = 206{,}768$,
$r_{31} = 3477{,}65$) daje $\theta_\mathrm{lept} = 132{,}73^\circ$.
R\'o\.znica $\Delta\theta = 0{,}21\,\mathrm{mdeg}$.
\end{hypothesis}

\begin{table}[h]
\centering
\small
\caption{Porównanie parametrów TGP z~warto\'sciami leptonowymi.
         Wyniki skrypt\'ow P21b--P26.}
\label{tab:tgp-lepton-comparison}
\begin{tabular}{@{}lccc@{}}
\toprule
\textbf{Wielko\'s\'c} & \textbf{TGP} & \textbf{Leptony} & \textbf{B\l\k{a}d} \\
\midrule
$r_{21} = m_\mu/m_e$ & $206{,}770$ & $206{,}768$ & $0{,}001\%$ \\
$r_{31} = m_\tau/m_e$ & $3477{,}1$ & $3477{,}65$ & $0{,}016\%$ \\
$Q$ (Koide) & $1{,}5003$ & $1{,}5000$ & $0{,}02\%$ \\
$\theta$ (Koide) & $132{,}73^\circ$ & $132{,}73^\circ$ & $0{,}21\,\mathrm{mdeg}$ \\
$K^*_1 (1+\alpha)/a_\Gamma$ & $2{,}351$ & --- & --- \\
\bottomrule
\end{tabular}
\end{table}

\begin{remark}[Zale\.zno\'s\'c analityczna $C_K = f(\alpha, a_\Gamma)$]
Wynik~\eqref{eq:CK-pert} pokazuje, \.ze $C_K$ \emph{nie jest}
magiczn\k{a} sta\l{}\k{a} — jest \textbf{wyprowadzalny analitycznie}
z~calki Yukawa $I_k(a)$ i~warunku samospójno\'sci.
Dla $a_\Gamma \to 0$: $C_K^{\mathrm{pert}} \to 2$ (cz\l{}on
wiod\k{a}cy). Korekcje:
\begin{itemize}[nosep]
    \item $\sim 8\%$ od geometrii siatki (czynnik $e^{2a_\Gamma}$),
    \item $\sim 2\%$ od mianownika ($\alpha a_\Gamma/(2(1+\alpha))$),
    \item $\sim 12\%$ od nieliniowo\'sci $(1+\alpha/\phi(r))$ przy $\phi(a_\Gamma) \approx 1{,}24$.
\end{itemize}
\end{remark}

% -------------------------------------------------------------------
\subsection{Hierarchia mas: co jest predykcj\k{a}, co dopasowaniem}\label{ssec:koide-hierarchy}
% -------------------------------------------------------------------

\begin{remark}[Kluczowe rozr\'o\.znienie]\label{rem:koide-predict-vs-fit}
Zewn\k{e}trzny czytelnik ma prawo pyta\'c:
\emph{,,Czy formu\l{}a Koidego wynika z~TGP, czy jest dopasowaniem?''}
Tabela poni\.zej daje pe\l{}n\k{a} odpowied\'z.
Kluczowy wynik: \textbf{$Q \approx 3/2$ jest warunkiem bifurkacji}
--- nie dopasowaniem --- ale \textbf{konkretne masy leptonowe}
wymagaj\k{a} wej\'sciowej warto\'sci $r_{21}$ lub $\alpha_K$,
kt\'ora jeszcze nie jest wyprowadzona z~substratu.
\end{remark}

\begin{center}\footnotesize
\begin{tabular}{@{}p{3.6cm}lp{3.2cm}p{3.0cm}@{}}
\toprule
\textbf{Wynik} & \textbf{Typ}
  & \textbf{Wej\'scia wymagane}
  & \textbf{Status / odsylacz} \\
\midrule

Trzy generacje ($n=0,1,2$)
  & Wyprowadzone z~WKB
  & $a_\Gamma$ (Warstwa~II)
  & stw.~\ref{prop:no-4th-generation} \\[4pt]

$Q \approx 3/2$ (warunek bifurkacji)
  & Warunek bifurkacji TGP
  & $a_\Gamma$, definicja $Q_{\rm TGP}$
  & hyp.~\ref{hyp:koide-tgp}; domkni\k{e}te P29--P51 \\[4pt]

$r_{21} = 206{,}77\ (m_\mu/m_e)$
  & \textbf{Wej\'scie kalibracji}
  & $r_{21}^{\rm PDG}$ z~obserwacji
  & wyznacza $\alpha_K \approx 8{,}56$ \\[4pt]

$\alpha_K \approx 8{,}56$
  & Dopasowane (OP-3 otwarte)
  & $r_{21}^{\rm PDG}$ lub $a_\Gamma$
  & hipoteza OP-3: $a_c(\alpha_K) = a_\Gamma$ \\[4pt]

$m_e,\, m_\mu,\, m_\tau$~(warto\'sci)
  & Odtwarzane
  & $\alpha_K$, $a_\Gamma$
  & po domkni\k{e}ciu OP-3 --- predykcja \\[4pt]

$r_{31} = m_\tau/m_e = 3477{,}1$
  & Predykcja (przy $\alpha_K$ i~$r_{21}$ wej\'sciowych)
  & $\alpha_K$, $a_\Gamma$
  & b\l{}\k{a}d $0{,}016\%$ wzgl.\ PDG \\[4pt]

$a_\Gamma = 0{,}040049$
  & Zdefiniowane z~$a_c(\alpha_K)$
  & $\alpha_K$
  & wzajemna relacja; niezale\.zno\'s\'c: cel OP-3 \\[4pt]

Masy kwark\'ow (analogi)
  & Analogia szkicowa
  & oddzielny $\alpha_q \neq \alpha_K$
  & Koide'go nie stosuje si\k{e} do kwark\'ow \\

\bottomrule
\end{tabular}
\end{center}

\begin{remark}[Otwart\k{a} kwestia: predykcja $\alpha_K$ z~substratu (OP-3)]%
\label{rem:OP3-koide}
\textbf{Kluczowy problem}: do\'s\'c moment TGP sektor mas jest
\textbf{rekonstrukcj\k{a}}, nie predykcj\k{a}:
$\alpha_K$ jest wyznaczane z~$r_{21}^{\rm PDG}$.
Domkni\k{e}cie OP-3 --- wykazanie,
\.ze $\alpha_K = \alpha_K(a_\Gamma)$ wynika z~warunku
bifurkacji $a_c(\alpha_K) = a_\Gamma$ --- zmieni\l{}oby ten status:
\begin{itemize}[nosep]
  \item $\alpha_K$ by\l{}oby \textbf{predykcj\k{a}} Warstwy~III,
  \item $r_{21}$, $r_{31}$ sta\l{}yby si\k{e} \textbf{predykcjami} (nie wej\'sciami),
  \item sektor mas TGP by\l{}by \textbf{w~pe\l{}ni predykcyjny}
    z~$N_{\rm param} = 2$ (tylko $\PhiZero$, $\psi_{\rm ini}$).
\end{itemize}
Biezacy status OP-3: badany numerycznie (skrypt \texttt{P72}).
Wyniki wst\k{e}pne: $a_c(\alpha_K)$ jest monotonicznie malej\k{a}c\k{a}
funkcj\k{a} $\alpha_K$; warto\'s\'c $a_c(8{,}56) \approx a_\Gamma$ w~granicy
kilku promili.
\end{remark}

% -------------------------------------------------------------------
\subsection{Problem otwarty O-K1: wyprowadzenie formu\l{}y Koide z~TGP}%
\label{app:F-koide-open}
% -------------------------------------------------------------------

\begin{problem}[Formu\l{}a Koide z~TGP --- CZ.~ZAMK. {[AN+NUM]}]\label{oprob:koide}
Czy warto\'s\'c $Q = (\Sigma\sqrt{K^*_n})^2/\Sigma K^*_n \approx 3/2$
jest konsekwencj\k{a} r\'owna\'n TGP, czy przypadkiem parametrycznym?

\medskip
\textbf{Precyzyjna tre\'s\'c problemu.}
Niech $(K^*_1, K^*_2, K^*_3)$ b\k{e}d\k{a} trzema zerami funkcji
samospójno\'sci~\eqref{eq:g-func} dla potencja\l{}u $V_\mathrm{mod}$
przy dowolnych dopuszczalnych parametrach $(\alpha, a_\Gamma, \lambda)$.
Zbadaj warto\'s\'c:
\begin{equation}\label{eq:Q-def-open}
    Q(\alpha, a_\Gamma, \lambda)
    = \frac{(\sqrt{K^*_1}+\sqrt{K^*_2}+\sqrt{K^*_3})^2}
           {K^*_1+K^*_2+K^*_3}.
\end{equation}

\textbf{Cztery pytania (w kolejno\'sci trudno\'sci):}
\begin{enumerate}[label=\textbf{OK-\arabic*}]
    \item \textbf{Ograniczenie strukturalne.}
          Czy $Q \geq 3/2$ dla wszystkich rozwi\k{a}za\'n TGP?
          \textit{Odpowied\'z (P27, P28): NIE.}
          $Q(\lambda)$ jest monotonicznie rosn\k{a}c\k{a}
          funkcj\k{a} $\lambda$; dla~$\lambda \to 0$
          maleje do~$Q_{\min} \approx 1{,}215$. Zatem
          $Q \geq 3/2$ \emph{nie jest} ograniczeniem strukturalnym TGP.

    \item \textbf{Szczelina Koide.}
          Jak du\.za jest odleg\l{}o\'s\'c punktu TGP od krzywej
          Koide'go?
          \textit{Wynik P28 (skrypt \texttt{p28\_koide\_gap.py}):}
          Dla punktu g\l{}\'ownego $(\alpha=8{,}5616,\,
          a_\Gamma=0{,}040,\, \lambda^*=5{,}501\!\times\!10^{-6})$:
          $\lambda_\mathrm{Koide} = 5{,}489\!\times\!10^{-6}$,
          stosunek $\lambda_K/\lambda^* = 0{,}9978$~(r\'o\.znica $0{,}22\%$),
          a~szczelina $\Delta r_{31} = r_{31}^K - r_{31}^* = +3{,}84$
          (tj.\ $0{,}11\%$). Stosunek $\lambda_K/\lambda^*$ nie jest
          universalny -- zale\.zy od $(\alpha, a_\Gamma)$.

    \item \textbf{Predykcja warto\'sci.}
          Czy TGP posiada dodatkow\k{a} zasad\k{e}
          (np.\ minimalizacja energii swobodnej, symetria dyskretna)
          wymuszaj\k{a}c\k{a} $(r_{21}, r_{31})$ na krzywej~\eqref{eq:koide-locus}?
          \textit{Status: otwarty.}

    \item \textbf{Punkt g\l{}\'owny jako punkt specjalny?}
          Dlaczego punkt g\l{}\'owny~$(\alpha^*,a_\Gamma^*)$ jest
          znacznie bli\.zej krzywej Koide'go ni\.z inne punkty
          rodziny optymalnej ($\Delta Q^* = +0{,}00025$ vs.\
          do~$+0{,}00138$ dla s\k{a}siednich punkt\'ow)?
          Czy minimum $\chi^2$ mas leptonowych pokrywa si\k{e}
          z~minimum~$|Q-3/2|$?
          \textit{Status: otwarty (problem O-K1d).}
\end{enumerate}

\textbf{Struktura analityczna.}
Dla znormalizowanych zer $K^*_1 = 1$, $K^*_2 = r_{21}$, $K^*_3 = r_{31}$:
\begin{equation}\label{eq:Q-r21-r31}
    Q(r_{21}, r_{31})
    = \frac{(1 + \sqrt{r_{21}} + \sqrt{r_{31}})^2}
           {1 + r_{21} + r_{31}}.
\end{equation}
Locus $Q = 3/2$ w~p\l{}aszczy\'znie $(r_{21}, r_{31})$ jest krzyw\k{a}:
\begin{equation}\label{eq:koide-locus}
    2(1 + \sqrt{r_{21}} + \sqrt{r_{31}})^2
    = 3(1 + r_{21} + r_{31}).
\end{equation}
Punkt leptonowy $(206{,}768,\; 3477{,}65)$ le\.zy na tej krzywej dok\l{}adnie;
punkt TGP $(207,\; 3477)$ le\.zy poza krzyw\k{a} o~$\Delta Q \approx 0{,}00025$.

R\'ownanie krzywej mo\.zna rozwi\k{a}za\'c analitycznie
wzgl\k{e}dem $r_{31}$ (P28):
\begin{equation}\label{eq:r31-koide}
    r_{31}^K(r_{21})
    = \Bigl(2a + \sqrt{6a^2 - 3 - 3r_{21}}\Bigr)^2,
    %% UWAGA: wersja z 1/4 przed nawiasem by\l{}a b\l{}\k{e}dna (errata v2)
    \quad a = 1 + \sqrt{r_{21}},
\end{equation}
z~asymptotyk\k{a}:
\begin{equation}\label{eq:r31-koide-asym}
    \frac{r_{31}^K}{r_{21}}
    \;\xrightarrow{r_{21}\to\infty}\;
    (2+\sqrt{3})^2 = 7 + 4\sqrt{3} \approx 13{,}928
    \;+\; \frac{4(5+3\sqrt{3})}{\sqrt{r_{21}}} + \mathcal{O}(r_{21}^{-1}).
\end{equation}
Dla $r_{21}=207$: warto\'s\'c wiod\k{a}ca daje $13{,}928$,
z~korekt\k{a} NLO $\approx 16{,}76$; dok\l{}adna: $r_{31}^K = 3481{,}0$.

\textbf{Wyniki sesji P29--P71 (skrypty \texttt{p29\_koide\_zero.py}--\texttt{p71\_K3\_lambda\_analysis.py}):}

\smallskip
\begin{description}[leftmargin=2em, labelindent=0pt]
\item[OK-1 --- \textbf{odpowiedziane (P28)}:]
  $Q \not\geq 3/2$ zawsze.
  $Q(\lambda) \xrightarrow{\lambda\to 0} Q_{\min} \approx 1{,}215$.

\item[OK-3 --- \textbf{odpowiedziane (P28)}:]
  Szczelina $\Delta r_{31} = r_{31}^K - r_{31}^{*}$ nie jest
  niezmiennicza wzgl\k{e}dem $\lambda$.

\item[OK-2/OK-4 --- \textbf{odpowiedziane (P29--P31, DOMKNIETE)}:]
  Istnieje $\alpha_0 = 8{,}5526$ (przy $a_\Gamma=0{,}040$), w~kt\'orym
  $Q = 3/2$ dok\l{}adnie przy
  $(r_{21}, r_{31}) = (206{,}746;\; 3477{,}10)$
  --- r\'o\.znica mniej ni\.z $0{,}02\%$ od mas PDG.
  \par
  Skrypt P30 pokaza\l{}, \.ze dok\l{}adne masy leptonowe PDG daj\k{a}:
  \begin{equation}\label{eq:Q-PDG}
      Q_\mathrm{PDG} = 1{,}499986,
      \quad \Delta Q = -1{,}38 \cdot 10^{-5}.
  \end{equation}
  \par
  \textbf{Mechanizm analityczny (P31)} --- skalowanie
  $K^*_3 \approx C \cdot a_\Gamma/\sqrt{\lambda}$, $C \approx 2{,}000$,
  wynika z~bilansu kwartyczno-seksycznego w~$V_\mathrm{mod}$:
  \begin{equation}\label{eq:K3-scaling}
      K_3^2 \approx \frac{9\,a_\Gamma^2}{2\lambda}
      \implies
      K_3 \approx \sqrt{4{,}5}\,\frac{a_\Gamma}{\sqrt{\lambda}}.
  \end{equation}
  Poniewa\.z $K^*_1, K^*_2$ s\k{a} niezale\.zne od $\lambda$
  (weryfikacja numeryczna: wyk\l{}adnik $\lambda^{0{,}000}$),
  warunek $Q = 3/2$ daje \textbf{algebraiczne r\'ownanie} na $\lambda$:
  \begin{equation}\label{eq:lambda-koide}
      \lambda_\mathrm{Koide}
      = \frac{C^2 a_\Gamma^2}{K_1^{*2}(\alpha)\,
              \bigl[r_{31}^K(r_{21})\bigr]^2}, \quad C \approx 2{,}000.
  \end{equation}
  Predykcja analityczna: $\lambda_K = 5{,}489 \cdot 10^{-6}$
  (numeryk P28: $5{,}489 \cdot 10^{-6}$; b\l{}\k{a}d $0{,}0\%$).

\item[Pytanie residualne --- \textbf{odpowiedziane (P32)}:]
  Czy $r_{21} \approx 207$ jest narzucone przez TGP,
  czy pochodzi z~obserwacji?
  \par
  Skrypt P32 wykaza\l{}, \.ze \textbf{nie istnieje} wewn\k{e}trzna
  zasada TGP wymuszaj\k{a}ca $r_{21} = 207$: dla ka\.zdego $r_{21}>1$
  mo\.zna dobra\'c $\lambda_K$ daj\k{a}ce $Q=3/2$. Warto\'s\'c
  $r_{21}=207$ pochodzi z~obserwacji mas leptonowych.
  Ponadto: dla kwark\'ow (PDG) $Q_{u,c,t} = 1{,}178$,
  $Q_{d,s,b} = 1{,}367$ --- formu\l{}a Koide'go \emph{nie stosuje
  si\k{e}} do kwark\'ow, co jest sp\'ojne z~danymi eksperymentalnymi.

\item[Weryfikacja na siatce --- \textbf{P33--P34 (domkni\k{e}te)}:]
  Skrypt P33 wykaza\l{}, \.ze sta\l{}a $C \approx 2{,}000$
  nie wynika bezpo\'srednio z~bilansu kwartyczno-seksycznego:
  dok\l{}adna formu\l{}a $K_3 = \sqrt{3I_4/(2\lambda I_6)}$
  daje $C = 1{,}77$ (b\l{}\k{a}d $-11\%$), a~$C \to \sqrt{4{,}5} = 2{,}121$
  dla $a_\Gamma \to 0$ (granica bezkontaktowa).
  Warto\'s\'c $C = 2{,}000$ przy $a_\Gamma = 0{,}040$ jest
  \emph{empiryczn\k{a}} sta\l{}\k{a} wynikaj\k{a}c\k{a} z~pe\l{}nego
  warunku $g(K_3^*) = 0$, gdzie cz\l{}ony kinetyczny i~sze\'scienny
  r\'ownie\.z odgrywaj\k{a} istotn\k{a} rol\k{e}.
  \par
  \textbf{Synteza P34} --- systematyczny test r\'ownania~\eqref{eq:lambda-koide}
  na siatce $8 \times 5$ punkt\'ow
  ($\alpha \in \{6,7,8,8{,}553,9,10,11,12\}$,
   $a_\Gamma \in \{0{,}025,\ldots,0{,}060\}$)
  potwierdza:
  \begin{itemize}[nosep]
    \item \textbf{$C$ nie zale\.zy od $\alpha$}: odchylenie po $\alpha$
          $< 0{,}0001$ (przy sta\l{}ym $a_\Gamma$);
    \item $C = C(a_\Gamma) \in [1{,}999,\, 2{,}015]$ dla badanego zakresu;
    \item b\l{}\k{a}d formu\l{}y~\eqref{eq:lambda-koide}:
          \'srednio $-0{,}50\%$, maksymalnie $-1{,}50\%$.
  \end{itemize}
  Formu\l{}a poprawiona:
  \begin{equation}\label{eq:lambda-koide-v2}
      \lambda_\mathrm{Koide}(\alpha, a_\Gamma)
      = \frac{C(a_\Gamma)^2\, a_\Gamma^2}
             {K_1^{*2}(\alpha, a_\Gamma)\,
              \bigl[r_{31}^K(r_{21}(\alpha, a_\Gamma))\bigr]^2},
  \end{equation}
  gdzie $C(a_\Gamma)$ jest analitycznie bliskie $\sqrt{4{,}5}$ dla
  $a_\Gamma \to 0$, a~numerycznie $\approx 2{,}000$ dla
  $a_\Gamma \approx 0{,}04$.
  \par
  \textbf{Pe\l{}na krzywa $C(a_\Gamma)$ (P35)} --- skan dla
  $a_\Gamma \in [0{,}005,\, 0{,}120]$ wykaza\l{}, \.ze $C \approx 2{,}000$
  jest wynikiem \textbf{kompensacji} dw\'och przeciwnych trend\'ow:
  \begin{itemize}[nosep]
    \item $C_\mathrm{qs}(a) = \sqrt{3I_4/(2I_6)}/a$ (bilans kwartyczno-seksyczny)
          maleje monotonicznie: $C_\mathrm{qs}(0{,}005) = 2{,}04$,
          $C_\mathrm{qs}(0{,}120) = 1{,}45$;
    \item $\Delta C = C_\mathrm{num} - C_\mathrm{qs}$ (wk\l{}ad
          $E_\mathrm{kin} + E_\mathrm{cub}$) ro\'snie monotonicznie:
          $\Delta C(0{,}005) = 0{,}031$, $\Delta C(0{,}120) = 0{,}586$;
    \item suma $C = C_\mathrm{qs} + \Delta C \approx 2{,}000 \pm 0{,}03$
          dla $a_\Gamma \in [0{,}01,\, 0{,}08]$.
  \end{itemize}
  Minimum: $C_\mathrm{min} \approx 1{,}999$ przy $a_\Gamma \approx 0{,}055$.
  Dopasowanie: $C(a) \approx \sqrt{4{,}5} - 4{,}12\,a + 30{,}5\,a^2$.
  Dla $C = 2{,}000$ (sta\l{}a): b\l{}\k{a}d $\lambda_\mathrm{Koide}$
  wynosi co najwy\.zej $3{,}9\%$ dla $a_\Gamma \in [0{,}01,\, 0{,}12]$.
  \par
  \textbf{Precyzyjna formu\l{}a $K_2$ i krajobraz fermion\'ow (P36)} ---
  analiza czynnika korekcyjnego $G = K_2^\mathrm{num}/K_2^\mathrm{cubic\text{-}\varepsilon}$
  na siatce $11 \times 5$ ($\alpha \in [3,12]$, $a_\Gamma \in [0{,}020,\,0{,}060]$):
  \begin{itemize}[nosep]
    \item $G \in [1{,}099,\, 1{,}363]$, ro\'snie monotonicznie z~$\alpha$;
    \item formu\l{}a cubic-$\varepsilon$: b\l{}\k{a}d $r_{21}$ \'sred.\ $-17{,}3\%$,
          max $26{,}3\%$;
    \item formu\l{}a NLO z~modelem bilinearnym:
    \begin{equation}\label{eq:K2-NLO}
        K_2^\mathrm{NLO} = K_2^\mathrm{eps}\cdot
        \bigl(1 + 0{,}01316\,\alpha + 2{,}6201\,a_\Gamma
              + 0{,}10759\,\alpha\,a_\Gamma\bigr),
    \end{equation}
    b\l{}\k{a}d $r_{21}^\mathrm{NLO}$: \'sred.\ $+1{,}8\%$, max $15{,}4\%$;
    \item cz\l{}on LO kinetyczny ($K(1{+}\alpha)/(2a)$) jest $5$--$8\times$
          za du\.zy wzgl\k{e}dem dok\l{}adnego ca\l{}kowania --- g\l{}\'owne
          \'zr\'od\l{}o b\l{}\k{e}du cubic-$\varepsilon$;
    \item krajobraz fermion\'ow: dla lepton\'ow ($r_{21}=206{,}8$)
          $\alpha_f \approx 5{,}9$ przy $a_\Gamma=0{,}025$ i $\alpha_f \approx 6{,}8$
          przy $a_\Gamma=0{,}030$; dla kwork\'ow $u$-type ($r_{21}=588$)
          $\alpha_f \approx 14{,}8$--$16{,}8$ w~tym samym zakresie $a_\Gamma$.
  \end{itemize}
  \textit{Status P36: zamkni\k{e}ty. Formu\l{}a $<5\%$ wymaga dalszej analizy (P37).}
  \par
  \textbf{Precyzyjna $K_1^\mathrm{NLO}$ (P37)} --- rozw\.ini\k{e}cie perturbacyjne
  $g(K) = E/(4\pi K) - 1 = 0$ dla $K \ll 1$ daje:
  \begin{equation}\label{eq:K1-NLO}
      K_1^\mathrm{NLO} = \frac{2}{\Delta_0}\!\left(1 + \frac{2B_2}{\Delta_0^2}\right),
      \quad
      \Delta_0 = (1{+}\alpha)\frac{e^{-2a}(2{+}a)}{2a} - \frac{e^{-2a}}{2},
  \end{equation}
  gdzie $B_2 = \alpha L_1(a) + \tfrac{4}{3}E_1(3a)$,
  $L_1(a) = \int_a^\infty e^{-3r}(1{+}r)^2/r^3\,dr$ (ca\l{}ka dok\l{}adna).
  B\l{}\k{a}d $K_1^\mathrm{NLO}$: max $1{,}5\%$ dla $\alpha \in [1{,}5,\,20]$.
  \par
  \textbf{Finalna formu\l{}a $r_{21}$ (b\l{}\k{a}d $<5\%$, P37)} ---
  \begin{equation}\label{eq:r21-NLO-final}
      r_{21}^\mathrm{NLO}(\alpha, a_\Gamma)
      = \frac{K_2^\mathrm{eps}(\alpha, a_\Gamma)\cdot G(\alpha, a_\Gamma)}
             {K_1^\mathrm{NLO}(\alpha, a_\Gamma)},
  \end{equation}
  b\l{}\k{a}d max $2{,}4\%$ na siatce $\alpha \in [1{,}5,\,20]$,
  $a_\Gamma \in [0{,}01,\,0{,}06]$ (poprzednio $26{,}8\%$).
  Krajobraz fermion\'ow rozszerzony do $a_\Gamma \in [0{,}01,\,0{,}04]$:
  leptony $\alpha_f \in [2{,}3,\,8{,}5]$, kwarki $u$-type
  $\alpha_f \in [7{,}0,\,20{,}3]$ w~tym zakresie.
  \textit{Status P37: zamkni\k{e}ty. Cel $<5\%$ osi\k{a}gni\k{e}ty.}
  \par
  \textbf{Bifurkacja $\alpha_c(a_\Gamma)$ i selektywno\'s\'c $Q=3/2$ (P38)} ---
  Wst\k{e}pna diagnoza kra jobrazu $d/s/b$ (patrz korekta P39).
  Por\'ownanie z~krzyw\k{a} Koide'go $Q(r_{21}, r_{31}) = 3/2$
  w~przestrzeni $(r_{21}, r_{31})$:
  \begin{itemize}[nosep]
    \item leptony $(r_{21},r_{31}) = (206{,}8;\;3477)$:
          \textbf{dok\l{}adnie na krzywej Koide'go} ($\delta r_{31} = -0{,}01\%$);
    \item kwarki $u/c/t$ $(588;\;79949)$:
          $r_{31}$ o $+770\%$ za du\.ze wzgl\k{e}dem krzywej;
    \item kwarki $d/s/b$ $(20;\;895)$:
          $r_{31}$ o $+89\%$ za du\.ze.
  \end{itemize}
  Wsp\'o\l{}czynnik $\lambda_\mathrm{obs}/\lambda_\mathrm{Koide}$:
  $1{,}0001$ (leptony), $0{,}013$ (up-kwarki), $0{,}54$ (dn-kwarki).
  \textbf{Wniosek}: formu\l{}a Koide $Q=3/2$ jest spe\l{}niona
  empirycznie \textbf{wy\l{}\k{a}cznie} przez leptony; kwarki
  le\.z\k{a} z~dala od krzywej algebraicznej $Q=3/2$.
  \textit{Status P38: zamkni\k{e}ty (korekta w P39).}
  \par
  \textbf{Korekta b\l{}\k{e}du numerycznego i nowe odkrycie (P39)} ---
  W~P38 funkcja \texttt{find\_K2\_num} szuka\l{}a $K_2$ w~przedziale
  $K \in [0{,}05;\,0{,}5]$, kt\'ory zawiera $K_1 \approx 0{,}07$--$0{,}08$,
  przez co \texttt{brentq} znajdowa\l{}o $K_1$ ponownie (fa\l{}szywe
  ,,podw\'ojne zero'' $K_1 = K_2$).
  Po poprawce (P39): $K_2$ szukane w~$K > 0{,}15$ daje
  $K_2 \approx 1{,}54$--$1{,}75$ dla wszystkich $\alpha \geq 0{,}05$.
  \par
  Prawdziwa struktura $r_{21}(\alpha)$ przy $a_\Gamma = 0{,}040$:
  funkcja \textbf{ci\k{a}g\l{}a i monotonicz nie rosn\k{a}ca}
  od $r_{21} \approx 18$ ($\alpha \to 0$) do $r_{21} \approx 68$ ($\alpha = 3$).
  Kwarki $d/s/b$ ($r_{21} = 20$) s\k{a} \textbf{osi\k{a}galne}
  przy $a_\Gamma \geq 0{,}040$, $\alpha_f \approx 0{,}22$.
  Minimalne osi\k{a}galne $r_{21}$ jako funkcja $a_\Gamma$:
  \begin{itemize}[nosep]
    \item $a_\Gamma = 0{,}030$: $r_{21}^\mathrm{min} \approx 24{,}4$ (d/s/b niemo\.zliwe);
    \item $a_\Gamma = 0{,}035$: $r_{21}^\mathrm{min} \approx 20{,}8$ (graniczne);
    \item $a_\Gamma = 0{,}040$: $r_{21}^\mathrm{min} \approx 18{,}1$ \textbf{(d/s/b mo\.zliwe!)};
    \item $a_\Gamma = 0{,}050$: $r_{21}^\mathrm{min} \approx 14{,}3$.
  \end{itemize}
  \textit{Status P39: zamkni\k{e}ty.}
  \par
  \textbf{$K_3$ universalne i predykcja $Q_\mathrm{TGP}$ (P40)} ---
  Dla sta\l{}ego $(\lambda, a_\Gamma)$,
  $K_3 \approx \mathrm{const}$ dla~\emph{wszystkich} rodzin fermionowych
  (niezale\.znie od $\alpha_f$):
  przy $\lambda = 10^{-5}$, $a_\Gamma = 0{,}040$ mamy
  $K_3 \approx 25{,}33$ jednocze\'snie dla leptonów, kwarków $u/c/t$
  i~kwarków $d/s/b$. Potwierdza to analityczn\k{a} formu\l{}\k{e}
  $K_3 \sim C\,a_\Gamma/\sqrt{\lambda}$ ($C \approx 2{,}000$)
  z~P31--P35, niezale\.zn\k{a} od $\alpha$.
  \par
  Przy $\lambda = \lambda_\mathrm{Koide} \approx 5{,}47 \times 10^{-6}$
  (wyznaczonej z~warunku $r_{31}^\mathrm{PDG}(\text{leptony}) = 3477$)
  model TGP daje:
  \begin{itemize}[nosep]
    \item leptony: $Q_\mathrm{TGP} = 1{,}5000$ (dok\l{}adnie --- z~definicji $\lambda_\mathrm{Koide}$);
    \item kwarki $u/c/t$: $Q_\mathrm{TGP} = 1{,}526$ vs $Q_\mathrm{PDG} = 1{,}178$
          ($\delta Q = +29{,}5\%$);
    \item kwarki $d/s/b$: $Q_\mathrm{TGP} = 1{,}517$ vs $Q_\mathrm{PDG} = 1{,}367$
          ($\delta Q = +11{,}0\%$).
  \end{itemize}
  \textbf{Wniosek}: TGP przewiduje $Q \approx 3/2$ \textbf{universalnie}
  dla wszystkich rodzin fermionowych; Natura spe\l{}nia $Q = 3/2$
  wy\l{}\k{a}cznie dla leptonów (kwarki: $Q < 3/2$).
  Oznacza to, \.ze $Q > 3/2$ dla kwark\'ow nie jest b\l{}\k{e}dem TGP,
  lecz \textbf{sygna\l{}em geometrycznego uwi\k{e}zienia}: soliton kwarkowy
  nie spe\l{}nia warunku samosp\'ojno\'sci $Q = 3/2$, wi\k{e}c nie mo\.ze
  istnie\'c jako wolna cz\k{a}stka.
  \textit{Status P40: zamkni\k{e}ty.}

  \textbf{Brak $Q = 3/2$ w~hadronach (P41)} ---
  Weryfikacja hipotezy: czy hadrony z\l{}o\.zone z~kwark\'ow maj\k{a}
  $Q_\mathrm{hadron} \approx 3/2$ na~poziomie z\l{}o\.zonym?
  Przeszukano $15\,180$ kombinacji 46 cz\k{a}stek z~bazy PDG\,2024.
  Wyniki:
  \begin{itemize}[nosep]
    \item najlepszy triplet hadronowy: $e/\pi^+/\Lambda_c^+$,
          $Q = 1{,}5007$ ($|\delta Q| = 0{,}05\%$), lecz zawiera lepton $e$;
    \item czysto hadronowe triplety mezonowe: $\delta Q \approx 60$--$90\%$;
    \item triplety barionowe: $\delta Q \approx 77$--$100\%$
          (np.\ $p/\Lambda/\Xi^0$: $Q = 2{,}986$;
          $\Upsilon(1S)/\Upsilon(2S)/\Upsilon(3S)$: $Q = 2{,}999$).
  \end{itemize}
  \textbf{Wniosek (P41)}: $Q = 3/2$ jest w\l{}a\'sciwo\'sci\k{a}
  \emph{fundamentalnych} soliton\'ow TGP (leptonow), nie uk\l{}ad\'ow
  z\l{}o\.zonych.
  Kwarki ($Q_\mathrm{TGP} > 3/2$) s\k{a} ,,sfrustrowanymi''
  solitonami: ich geometria ($r_{31}^\mathrm{TGP} \gg r_{31}^\mathrm{Koide}$)
  wyklucza samosp\'ojno\'s\'c jako wolna cz\k{a}stka~---
  \textbf{uwi\k{e}zienie geometryczne w~TGP}.
  Hadrony (kompozyty soliton\'ow) maj\k{a} $Q_\mathrm{hadron} \gg 3/2$,
  co jest konsekwencj\k{a} dynamiki oddzialywania koloru, nie warunku
  Koide'go.
  \textit{Status P41: zamkni\k{e}ty.}

  \textbf{Masy konstytuent\'ow i diagram $(r_{21}, r_{31})$ (P42)} ---
  Zbadano 8 modeli mas konstytuent\'ow kwark\'ow (De~R\'ujula, NJL,
  bag model, ChQSM, PDG MS-bar/pole i~in.).
  \.Zaden nie daje $Q \approx 3/2$:
  $Q(u/c/t) \approx 1{,}18$--$1{,}35$, $Q(d/s/b) \approx 1{,}37$--$2{,}4$.
  \par
  \emph{Diagram $(r_{21}, r_{31})$}:
  Krzywa $Q = 3/2$ dzieli p\l{}aszczyzn\k{e} na dwa obszary.
  \begin{itemize}[nosep]
    \item Leptony: le\.z\k{a} \emph{dok\l{}adnie na krzywej} ---
          $r_{31}^\mathrm{TGP} = r_{31}^\mathrm{Koide}$ z~definicji $\lambda_\mathrm{Koide}$;
    \item Kwarki TGP ($\lambda = \lambda_\mathrm{Koide}$): le\.z\k{a}
          \emph{poni\.zej krzywej}
          ($r_{31}^\mathrm{TGP}/r_{31}^\mathrm{Koide} \approx 0{,}89$
          dla $u/c/t$ i $\approx 0{,}93$ dla $d/s/b$)
          $\Rightarrow$ $Q_\mathrm{TGP} > 3/2$;
    \item Kwarki PDG: le\.z\k{a} \emph{powy\.zej krzywej}
          ($r_{31}^\mathrm{PDG} \gg r_{31}^\mathrm{Koide}$)
          $\Rightarrow$ $Q_\mathrm{PDG} < 3/2$.
  \end{itemize}
  TGP i~PDG (kwarki) le\.z\k{a} po \textbf{przeciwnych stronach} krzywej
  $Q = 3/2$ --- fundamentalne napi\k{e}cie modelu.
  \par
  \emph{Analiza $\lambda_\mathrm{eff}$}:
  $\lambda$ wymagane dla $Q = 3/2$ dla ka\.zdej rodziny:
  $\lambda_\mathrm{eff}(u/c/t) \approx 0{,}80\,\lambda_\mathrm{Koide}$,
  $\lambda_\mathrm{eff}(d/s/b) \approx 0{,}87\,\lambda_\mathrm{Koide}$
  --- r\'o\.znice tylko $13$--$20\%$.
  Jednak $K_3$ jest universalne (nie zale\.zy od $\alpha_f$ przy sta\l{}ym
  $\lambda$), wi\k{e}c nie mo\.zna jednocze\'snie spe\l{}ni\'c $Q = 3/2$
  dla wszystkich rodzin jednym parametrem $\lambda$.
  \par
  \textbf{Wniosek P42}: $Q = 3/2$ dotyczy wy\l{}\k{a}cznie fundamentalnych
  soliton\'ow (leptonow). Kwarki s\k{a} ,,prawie na krzywej'' ($\approx 8\%$
  od strony uwi\k{e}zienia), co sugeruje mo\.zliwe rozszerzenie modelu
  (kolor QCD lub wielocia\l{}owe TGP).
  \textit{Status P42: zamkni\k{e}ty.}

  \textbf{Twierdzenie o niekompatybilno\'sci GMO--Koide (P43)} ---
  \emph{Pytanie}: czy relacja Gell-Mann--Okubo (liniowa zale\.zno\'s\'c
  mas od \l{}adunku strangeness) implikuje $Q \approx 3/2$?
  \par
  \emph{Twierdzenie} (udowodnione analitycznie w P43):
  Dla dowolnej progresji arytmetycznej $m_k = M + k\,d$
  ($k = 0, 1, 2$; $M > 0$, $d \geq 0$):
  \begin{equation}
    Q_\mathrm{min}
    = \lim_{d/M \to \infty} Q
    = \frac{(1+\sqrt{2})^2}{3}
    = \frac{3 + 2\sqrt{2}}{3}
    \approx 1{,}9428 > \frac{3}{2}.
  \end{equation}
  \emph{Dow\'od}: przy $t = d/M \to \infty$,
  $Q(t) \to (\sqrt{t} + \sqrt{2t})^2/(3t) = (1+\sqrt{2})^2/3$.
  Poniewa\.z $Q(t)$ jest \'sci\'sle malej\k{a}ce od 3 do $1{,}9428$,
  a $3/2 = 1{,}5 < 1{,}9428$, warunek $Q = 3/2$ nie jest nigdy spe\l{}niony.
  \par
  \textbf{Konsekwencja}: relacja GMO (linioiwe masy w ramach symetrii SU(3))
  i warunek Koide'go ($Q = 3/2$) s\k{a}
  \textbf{fundamentalnie niekompatybilne} --- nie istnieje progresja
  arytmetyczna spe\l{}niaj\k{a}ca $Q = 3/2$.
  Weryfikacje numeryczne:
  \begin{itemize}[nosep]
    \item GMO dla bariionowego oktetu: odchylenie $0{,}57\%$;
    \item GMO dla bariionowego dekupletu: progresja arytmetyczna
          \'sr.~$147\,\text{MeV}$, std/\'sr.~$3{,}9\%$ $\approx$ r\'ownorozmieszczona;
    \item $Q(\Delta/\Sigma^*/\Xi^*/\Omega) \approx 2{,}99$ (bliskie
          $Q_\mathrm{max} = 3$ ze wzgl\k{e}du na bliskie masy).
  \end{itemize}
  \textbf{Wniosek P43}: $Q = 3/2$ jest \textbf{geometryczn\k{a} w\l{}a\'sciwo\'sci\k{a}
  soliton\'ow TGP}, niewynikaj\k{a}c\k{a} z \.zadnej symetrii Lie ani
  liniowej relacji masowej.
  Jest to nowe prawo masy, odr\k{e}bne od standardowej symetrii SU(3).
  \textit{Status P43: zamkni\k{e}ty.}

  \textbf{Analityczne zamkni\k{e}cie — twierdzenie TGP (P44)} ---
  Cel P44: analityczne wyja\'snienie, dlaczego $Q = 3/2$ jest wyodr\k{e}bnion\k{a}
  warto\'sci\k{a} dla soliton\'ow TGP.
  \par
  \emph{Kluczowe obserwacje}:
  \begin{enumerate}[nosep]
    \item $K_1, K_2$ s\k{a} praktycznie niezale\.zne od $\lambda$
          (zmiana $<0{,}5\%$ dla $\lambda \in [10^{-6}, 10^{-4}]$);
          jedynie $K_3 = C a_\Gamma / \sqrt{\lambda}$ niesie zale\.zno\'s\'c od $\lambda$.
    \item $Q(\lambda) = Q(K_1, K_2, K_3(\lambda))$ jest \'sci\'sle malej\k{a}ce
          (od 3 przy $\lambda \to 0$ do ok.\ 1{,}14 przy $\lambda \to \infty$),
          wi\k{e}c istnieje dok\l{}adnie jedno $\lambda_\mathrm{Koide}$ spe\l{}niaj\k{a}ce
          $Q = 3/2$.
    \item R\'ownanie $Q = 3/2$ posiada dwa analityczne rozwi\k{a}zania na $K_3$:
          $K_3^{(+)} = 34{,}21$ (fizyczne, $r_{31} = 3477$) oraz
          $K_3^{(-)} = 0{,}064$ (niefizyczne, $K_3 < K_2$).
  \end{enumerate}
  \par
  \emph{Twierdzenie (TGP, P44)}: Dla modelu solitonowego TGP istnieje
  zamkni\k{e}ta formu\l{}a:
  \begin{equation}
    \lambda_\mathrm{Koide}
    = \frac{(C\,a_\Gamma)^2}{\bigl[K_1(\alpha_f, a_\Gamma)
      \cdot r_{31}^K(r_{21})\bigr]^2},
    \quad
    r_{31}^K
    = \bigl[2(1+\sqrt{r_{21}}) + \sqrt{3(1+4\sqrt{r_{21}}+r_{21})}\bigr]^2,
  \end{equation}
  gdzie $r_{21} = K_2/K_1$.
  Weryfikacja: $\lambda_\mathrm{Koide} = 5{,}467 \times 10^{-6}$
  (wz\'or) vs.\ $5{,}489 \times 10^{-6}$ (P28, numerycznie);
  zgodno\'s\'c $0{,}4\%$; $Q(K_1, K_2, K_3^K) = 3/2$ dok\l{}adnie.
  \par
  \emph{Diagram fazowy} $(r_{21}, r_{31})$:
  krzywa $Q = 3/2$ dzieli przestrze\'n na
  region uwi\k{e}zienia ($r_{31} < r_{31}^K \Rightarrow Q > 3/2$,
  soliton kwarkowy TGP),
  krzywa stabilno\'sci ($r_{31} = r_{31}^K \Rightarrow Q = 3/2$, leptony),
  oraz region zewn\k{e}trzny ($r_{31} > r_{31}^K \Rightarrow Q < 3/2$,
  efektywne masy kwark\'ow PDG).
  \par
  \textbf{Wniosek P44}: $Q = 3/2$ jest \textbf{geometryczn\k{a} w\l{}a\'sciwo\'sci\k{a}
  solitonu TGP} wynikaj\k{a}c\k{a} z universalno\'sci $K_3$ i struktury
  r\'ownania samosp\'ojno\'sci.
  Zamkni\k{e}ty cykl analityczny P31--P44.
  \textit{Status P44: zamkni\k{e}ty.}

\item[OK-7 --- \textbf{Oddzialywania kolorowe (P45)}:]
  \textbf{Minimalne rozszerzenie TGP o cz\l{}on kolorowy $\beta_c/\varphi^2$}.
  \par
  \textit{Motywacja}: $Q_\text{TGP}(\text{kwarki}) \approx 1{,}52 > 3/2$;
  $Q_\text{PDG}(e/\mu/\tau) = 1{,}500\,014 \neq 3/2$ (odchylenie $+9\ \text{ppm}$,
  mieszcz\k{a}ce si\k{e} w niepewno\'sci $m_\tau = 1776{,}86 \pm 0{,}12\ \text{MeV}$).
  \par
  Sprz\k{e}\.zenie kinetyczne TGP rozwini\k{e}te w szereg $1/\varphi$:
  \begin{equation}
    f(\varphi) = 1 + \frac{\alpha}{\varphi} + \frac{\beta_c}{\varphi^2} + \cdots,
    \quad \varphi \sim \frac{K}{r} \Longrightarrow
    \frac{\beta_c}{\varphi^2} \sim \frac{\beta_c}{K^2}\,r^2.
  \end{equation}
  Cz\l{}on $\beta_c/\varphi^2$ odpowiada potencja\l{}owi rosn\k{a}cemu jak $r^2$
  (konfajnment QCD), podczas gdy $\alpha/\varphi \sim r$ odpowiada oddzia\l{}ywaniu EM.
  Minimalne rozszerzenie funkcjonalu energii:
  \begin{equation}\label{eq:E-kolor}
    E[K;\beta_c] = 4\pi \int_0^\infty
      \frac{1}{2}\!\left(\frac{d\varphi}{dr}\right)^{\!2}
      \!\left(1 + \frac{\alpha}{\varphi} + \frac{\beta_c}{\varphi^2}\right) r^2\,dr
      + E_p[K],
  \end{equation}
  z warunkiem $\beta_c^{\text{leptony}} = 0$, $\beta_c^{\text{kwarki}} > 0$.
  Potencja\l{} $V_\text{mod}$ i r\'ownanie samosp\'ojno\'sci $g(K)=0$
  pozostaj\k{a} niezmienione.
  \par
  \textbf{Wyniki perturbacyjne} ($\alpha=8{,}54$, $\lambda=\lambda_\text{Koide}$,
  $\beta_c \to 0^+$):
  \begin{itemize}[nosep]
    \item $dK_3/d\beta_c \approx 0$ (dok\l{}adnie $0{,}00\%$/jed.)
          --- universalno\'s\'c $K_3$ \textbf{zachowana} przy korekcie kolorowej,
    \item $dK_1/d\beta_c = -10{,}25\%$/jed. ($K_1$ najbardziej czuly),
    \item $dQ/d\beta_c = -0{,}00112 < 0$
          --- korekta kolorowa \textbf{zmniejsza} $Q$,
          przesuwaj\k{a}c kwarki z $Q > 3/2$ ku $Q = 3/2$.
  \end{itemize}
  \par
  \textbf{Diagram fazowy $(r_{21}, r_{31})$ z korekt\k{a} koloru}:
  kwarki TGP ($r_{21} \approx 600$, $r_{31} = 8196$) le\.z\k{a}
  \emph{poni\.zej} krzywej $Q=3/2$ (tzn.\ $Q > 3/2$);
  po dodaniu $\beta_c > 0$: $r_{21}$ ro\'snie, punkt
  przesuwa si\k{e} ku krzywej ($\Delta r_{21} \approx -13\%$
  potrzebne do osi\k{a}gni\k{e}cia $Q=3/2$).
  \par
  \textbf{Wniosek P45}: cz\l{}on $\beta_c/\varphi^2$ jest \textbf{formalnie spójnym
  rozszerzeniem} TGP o konfajnment kolorowy, zachowuj\k{a}cym
  universalno\'s\'c $K_3$ (P31) i formu\l{}\k{e} $\lambda_\text{Koide}$ (P44).
  Warto\'s\'c $\beta_c$ wymaga dopasowania do mas hadron\'ow (P46+).
  \textit{Status P45: zamkni\k{e}ty.}

\item[OK-8 --- \textbf{Asymetria rodzin kwarkowych (P46)}:]
  \textbf{Dopasowanie $\beta_c^*$ do kwarków: odkrycie asymetrii rodzinowej}.
  \par
  \textit{Wyniki skanowania $\beta_c$ przy $\lambda = \lambda_\text{Koide}$:}
  \begin{itemize}[nosep]
    \item Rodzina $d/s/b$ ($\alpha = 0.22$, $K_1 = 0.0777$, $r_{21} = 20$):
          $Q_\text{TGP}(0) = 1{,}517$; $\beta_c^*(d/s/b) = \mathbf{1{,}17}$ ---
          ma\l{}a, dodatnia, \textbf{fizyczna} korekta kolorowa,
    \item Rodzina $u/c/t$ ($\alpha = 20.3$, $K_1 = 0.00418$, $r_{21} = 588$):
          $Q_\text{TGP}(0) = 1{,}526$; $Q(\beta_c)$ ma minimum
          $Q_\text{min} \approx 1{,}525 > 3/2$ przy $\beta_c \approx 5$--$8$,
          po czym ro\'snie --- \textbf{brak przeci\k{e}cia $Q=3/2$} dla $\beta_c > 0$.
  \end{itemize}
  \par
  \textbf{Odkrycie}: znak $dQ/d\beta_c$ jest \emph{r\'o\.zny} dla obu rodzin
  kwarkowych: ujemny dla $d/s/b$ (korekcja prowadzi ku $Q=3/2$) i~dodatni
  dla $u/c/t$ przy du\.zych $\beta_c$ (korekcja oddala od $Q=3/2$).
  Asymetria wynika z~drastycznie r\'o\.znych warto\'sci $K_1$:
  $K_1(u/c/t) = 0{,}00418 \ll K_1(d/s/b) = 0{,}0777$.
  \par
  \textbf{Universalno\'s\'c $K_3$ zachowana} przy $\beta_c^*$:
  odchylenie $K_3^\text{num}$ od $K_3^\text{analyt} = C a_\Gamma/\!\sqrt{\lambda_K}$
  wynosi jedynie $0{,}10$--$0{,}11\%$.
  \par
  \textbf{Wniosek P46}: korekta $\beta_c/\varphi^2$ jest \textbf{efektywna dla
  rodziny $d/s/b$} ($\beta_c^* = 1{,}17$), natomiast rodzina $u/c/t$
  wymaga innego mechanizmu --- mo\.zliwa modyfikacja $\lambda_\text{eff}$
  lub wy\.zsze cz\l{}ony w szeregu $1/\varphi^n$.
  \textit{Status P46: zamkni\k{e}ty.}

\item[OK-9 --- \textbf{Analityczne ca\l{}ki energii (P47)}:]
  Rozwijamy $E[K] = 4\pi \sum_n c_n K^n$ z analitycznymi wsp\'o\l{}czynnikami
  przez ca\l{}ki wyk\l{}adnicze $E_1(x) = -\mathrm{Ei}(-x)$:
  \begin{align}
    U_2(a) &= e^{-2a}/2, \quad
    U_3(a) = E_1(3a), \quad
    U_4(a) = e^{-4a}/a - 4E_1(4a), \\
    \Phi_2(a) &= J_2 + 2J_1 + J_0
    \quad \text{(przez } J_k = \int_a^\infty e^{-2r}/r^k\,dr = E_1(2a)\text{-formu\l{}y)}.
  \end{align}
  Wsp\'o\l{}czynniki energii: $c_2 = (1+\alpha)\Phi_2/2 - U_2/2$,
  $c_3 = -\alpha\Phi_3/2 - 2U_3/3$,
  $c_4 = \alpha\Phi_4/2 - U_4/4$.
  \par
  \textbf{Analityczne $K_1$} (rz\k{a}d NLO, r\'ownanie kwadratowe):
  \begin{equation}
    \boxed{K_1^{(\mathrm{NLO})} = \frac{-c_2 + \sqrt{c_2^2 + 4c_3}}{2c_3}},
    \quad \text{b\l{}\k{a}d}\ {+1{,}9\%}\ \text{vs.\ P40}.
  \end{equation}
  \textbf{Analityczne $K_3$} przez dok\l{}adne ca\l{}ki $I_4, I_6$ (funkcje $E_1$):
  \begin{equation}
    K_3^2 = \frac{3I_4(a)}{2\lambda I_6(a)},\quad
    I_4 = \frac{e^{-4a}}{a} - 4E_1(4a),
  \end{equation}
  daje $C = K_3\sqrt{\lambda}/a_\Gamma = 1{,}9965$
  (b\l{}\k{a}d $-0{,}17\%$, vs.\ $C=2{,}000$ z~P31).
  \par
  \textbf{Twierdzenie o~trzech zerach (P47)}: $g(K)$ ma dok\l{}adnie trzy zera
  wtedy i~tylko wtedy, gdy $(c_2>0) \land (\alpha > \alpha_c) \land (\lambda>0)$.
  Wszystkie warunki s\k{a} spe\l{}nione dla rodzin fermionowych TGP.
  \textit{Status P47: zamkni\k{e}ty.}

\item[OK-10 --- \textbf{$K_2$ semi-analityczne: Pad\'e i rozwini\k{e}cie wok\'o\l{} maksimum (P48)}:]
  $K_2\approx 2{,}03$ le\.zy w~reżimie przej\'sciowym ($r^*=\ln K\approx 0{,}7$),
  gdzie szereg ma\l{}ych~$K$ jest rozbie\.zny.
  \par
  \textbf{Metoda Pad\'e $[2/2]$}: aproksymacja
  $g(K)\approx (a_0+a_1K+a_2K^2)/(1+b_1K+b_2K^2)$ dopasowana do $30$ punkt\'ow
  daje zero licznika:
  \begin{equation}
    K_2^{\mathrm{Pad\'e}} = \frac{-a_1 + \sqrt{a_1^2 - 4a_0a_2}}{2a_2}
    = 2{,}0304 \quad (\text{b\l{}\k{a}d}\ {-0{,}2\%}).
  \end{equation}
  \textbf{Formu\l{}a z~krzywizn\k{a} maksimum}:
  \begin{equation}
    K_2 \approx K_\mathrm{max} + \sqrt{\frac{g_\mathrm{max}}{b}},\qquad
    b = -\tfrac{1}{2}g''(K_\mathrm{max}),
  \end{equation}
  z~$K_\mathrm{max}=1{,}030$, $g_\mathrm{max}=18{,}15$, $b=14{,}87$:
  $K_2 = 2{,}135$ (b\l{}\k{a}d $+5\%$).
  Analityczne wyznaczenie $K_\mathrm{max}$ z~warunku $K E'(K)=E(K)$
  wymaga modelu dwustrefowego (P49).
  \textit{Status P48: zamkni\k{e}ty (semi-analityczny).}

\item[OK-11 --- \textbf{Model dwustrefowy $g(K)$ i pe\l{}ny \l{}a\'ncuch analityczny (P49--P50)}:]
  Profil $\varphi(r)=1+Ke^{-r}/r$ naturalnie dzieli si\k{e} na dwie strefy
  przy $r^*(K)$ okre\'slonym przez $K e^{-r^*}/r^* = 1$:
  \begin{itemize}[nosep]
    \item Strefa~1, $r\in[a,r^*]$: $\varphi \approx Ke^{-r}/r \gg 1$ (j\k{a}dro solit.),
    \item Strefa~2, $r\in[r^*,\infty)$: $\varphi \approx 1+\varepsilon$ (ogon pola).
  \end{itemize}
  Model daje $g_\text{2-str}(K)\approx g_\text{num}(K)$ dla $K>1$
  z~b\l{}\k{e}dem systematycznym $\approx 5$ (na tle $g_\text{max}\approx 18$).
  \par
  \textbf{Pe\l{}ny \l{}a\'ncuch analityczny} (P50):
  \begin{equation}
    E[K] \;\xrightarrow{c_n\text{ przez }E_1}\;
    K_1^{(\mathrm{NNLO})} \xrightarrow{} r_{21}
    \;\;\wedge\;\;
    K_3^{(E_1)} \xrightarrow{} r_{31}
    \;\;\Rightarrow\;\; Q(r_{21},r_{31}).
  \end{equation}
  Wyniki na r\'o\.znych poziomach precyzji:
  \begin{center}
  \begin{tabular}{lcc}
  \hline
  Poziom & $Q_\text{analyt.}$ & $Q - 3/2$ [ppm] \\ \hline
  LO & 1.484624 & $-15376$ \\
  NLO ($K_1^\mathrm{NLO}$, $K_3^{E_1}$) & 1.500362 & $+362$ \\
  NNLO ($K_1^\mathrm{NNLO}$, $K_3^{E_1}$) & 1.499938 & $-62$ \\
  Numeryczny & 1.500013 & $+13$ \\ \hline
  \end{tabular}
  \end{center}
  Warto\'s\'c $3/2$ jest \textbf{zamkni\k{e}ta mi\k{e}dzy NLO a~NNLO} -- to
  niezale\.zna weryfikacja zbie\.zno\'sci analitycznego \l{}a\'ncucha do $Q=3/2$.
  \par
  \textbf{Przewidywanie $\lambda_\text{Koide}$}: szukaj\k{a}c $\lambda$ takie,
  \.ze $Q_\text{analyt.}(\lambda)=3/2$ (przez $K_3(\lambda)$):
  \begin{equation}
    \frac{\lambda_\text{analyt.}}{\lambda_\text{Koide}} = 1.000554
    \quad (+0{,}055\%),
  \end{equation}
  co oznacza, \.ze analityczny \l{}a\'ncuch przewiduje $\lambda_\text{Koide}$
  z~dok\l{}adno\'sci\k{a} $0{,}055\%$ \emph{bez numerycznego dopasowania}.
  \par
  Asymptotyczne przybli\.zenie $Q$ dla $r_{31}\gg r_{21}\gg 1$:
  \begin{equation}
    Q \approx 1 + \frac{2(1+\sqrt{r_{21}})}{\sqrt{r_{31}}} - \frac{1}{r_{31}} + \ldots,
    \quad Q=\tfrac{3}{2} \Rightarrow r_{31}^K \approx 16(1+\sqrt{r_{21}})^2
    \;\;(\text{b\l{}\k{a}d }+9\%).
  \end{equation}
  \textit{Status P49--P50: zamkni\k{e}ty. Pe\l{}ny \l{}a\'ncuch analityczny TGP$\to Q\approx3/2$.}

\item[OK-12 --- \textbf{Wz\'or zamkni\k{e}ty $K_1$ przez resumacj\k{e} Pad\'e~$[2/1]$ (P51)}:]
  Dotychczasowe metody: $K_1^\mathrm{NLO}$ (b\l{}\k{a}d $+1{,}85\%$) i~$K_1^\mathrm{NNLO}$
  (b\l{}\k{a}d $-0{,}30\%$) nie posiadaj\k{a} wsp\'olnego wzoru zamkni\k{e}tego
  l\'epszego ni\.z wielomian stopnia~3. Aproksymacja Pad\'e~$[2/1]$
  \begin{equation}
    h(K)=g(K)+1 \;\approx\; \frac{c_2 K + a_2 K^2}{1 + b_1 K},
    \quad
    b_1 = -\frac{c_4}{c_3},
    \quad
    a_2 = \frac{c_3^2 - c_2 c_4}{c_3},
  \end{equation}
  daje z~warunku $h(K_1)=1$ r\'ownanie kwadratowe z~\textbf{wzorem zamkni\k{e}tym}:
  \begin{equation}\label{eq:K1-pade21}
    K_1^{[2/1]} = \frac{c_3\Bigl[-(c_2 c_3+c_4)+\sqrt{4c_3^3+(c_2 c_3-c_4)^2}\Bigr]}
                       {2(c_3^2-c_2 c_4)},
  \end{equation}
  gdzie $c_2,c_3,c_4$ wyra\.zone s\k{a} przez \textbf{ca\l{}ki analityczne z~$E_1(na)$}.
  Numerycznie dla lepton\'ow ($\alpha=8{,}54$, $a=0{,}04$):
  $b_1=16{,}024$, $a_2=567{,}3$, dyskryminant $=11500{,}9$.
  \par
  Zestawienie wszystkich metod wyznaczania $K_1$ i~wynik\'ow $Q$:
  \begin{center}
  \begin{tabular}{lccc}
  \hline
  Metoda & $K_1\;[\times10^{-3}]$ & b\l{}\k{a}d & $Q-\tfrac{3}{2}$ [ppm] \\ \hline
  NLO (kwadrat.)      & 10{,}021 & $+1{,}85\%$ & $+362$ \\
  NNLO (sze\'scian.)  & 9{,}809  & $-0{,}30\%$ & $-62$ \\
  \textbf{Pad\'e~$[2/1]$} & \textbf{9{,}837} & $\mathbf{-0{,}025\%}$ & $\mathbf{-7{,}6}$ \\
  Numeryczny          & 9{,}839  & $0{,}000\%$  & $+12{,}9$ \\ \hline
  \end{tabular}
  \end{center}
  Wz\'or~\eqref{eq:K1-pade21} daje $Q_{[2/1]}=1{,}499992$,
  czyli \textbf{$-7{,}6$~ppm od~$3/2$} -- najdok\l{}adniejszy spo\'sr\'od
  wszystkich metod analitycznych.
  \par
  \textbf{Finalny \l{}a\'ncuch analityczny} (P51):
  \begin{equation}
    K_1^{[2/1]}\;(c_2,c_3,c_4 \text{ przez }E_1)
    \;\wedge\; K_2^\mathrm{Pad\'e}\;(-0{,}20\%)
    \;\wedge\; K_3^{E_1}\;(-0{,}17\%)
    \;\Rightarrow\; Q = 1{,}499992 \;(-7{,}6\;\text{ppm}).
  \end{equation}
  \textit{Status P51: zamkni\k{e}ty. Pe\l{}na analityczna synteza $K_1,K_2,K_3\to Q\approx3/2$.}
\end{description}
\end{problem}

\begin{remark}[Diagnostyka $\varphi$-FP: luka masy tau --- sesja v42+]%
\label{rem:tau-gap-v42plus}
Niezale\.zna analiza z~perspektywy $\varphi$-FP
(\'scie\.zka~9, dodatek~\ref{app:J2-formalizacja}) ujawnia
\textbf{luk\k{e} masy tau}: pojedynczy soliton
z~$f(g)=1+4\ln g$, $V'(g)=g^2(1-g)$ daje
$r_{31}^{\max} = 832$ (przy $g_0\approx 2{,}0$),
gdy Koide wymaga $r_{31}^K = 3477$
(\ref{eq:r31-max}, skrypty p124a--d).

\textbf{Mechanizm}: redukcja efektywnego sprz\k{e}\.zenia
kinetycznego $\alpha_{\rm eff}$ z~2.0 do $\approx 0{,}92$
daje $r_{31} \sim 3470$. Zachowanie r\'ownoczesne
$r_{21} = 206{,}77$ wymaga $\alpha_{\rm eff}$
zale\.znego od~$g$ (r\'ownanie~\eqref{eq:alpha-eff-running}).
Bieg ten jest naturaln\k{a} konsekwencj\k{a}
przep\l{}ywu ERG Wettericha dla sektora kinetycznego $K_k(\psi)$.

\textbf{Aktualizacja po p130--p131 i~korekcie p134}:
Re-derywacja $\varphi$-FP z~running $\alpha_{\rm eff}$
(tj.~jednoczesna re-kalibracja $g_0^e$) daje przy
$\eta_K = 181/15 \approx 12{,}0667$:
$g_0^e = 0{,}90548$,
$r_{21} = 206{,}768$ (\textbf{zachowane}),
$\Delta(\delta_e \to \delta_\mu) = 120{,}01^\circ$
(\textbf{dok\l{}adne} $2\pi/3$).

\textbf{Wyprowadzenie analityczne $\eta_K$ (p139--p140):}
$\eta_K = \alpha_{\rm UV}^2 d + 1/((\alpha_{\rm UV}^2+1)d)
= 12 + 1/15 = 181/15$.
Z~ERG Wetterich LPA': rz\k{a}d wiod\k{a}cy $\alpha^2 d = 12$,
korekta back-reaction $1/((\alpha^2+1)d) = 1/15$.
Dopasowanie: 2\,ppm. \textbf{Status $\eta_K$: ZAMKNI\k{E}TY [AN].}
Z~$g_0^\tau = 4$: $m_\tau = 1776{,}97$\,MeV (0{,}006\% od obs.).

\textbf{Korekta (p134e--g):} $A_\tau(g_0^\tau)$ jest monotonicznie
rosn\k{a}ca --- warto\'s\'c $r_{31}=3477{,}5$ z~p131 wynika\l{}a
z~ograniczenia skanowania $g_0^\tau \le 4{,}0$ (kres zakresu,
nie fizyczne maksimum).
\textbf{Problem O-K1: ZAMKNI\k{E}TY [AN] (rz\k{a}d wiod\k{a}cy)} ---
$\eta_K$ zamkni\k{e}ty analitycznie; $g_0^\tau = 4$ z~warunku bilansu
potencja\l{}owo-kinetycznego: $|V'(g_0)|/|V'(\sqrt{g_0})| = \alpha^2 d$
daje $t^3+t^2-12=0$ z~jedynym pierwiastkiem $t=2$, $g_0=4$ (p146).
Pe\l{}ny opis: \S\ref{app:H-v42plus-tau-gap}.
\end{remark}

% -------------------------------------------------------------------
\subsection{Analityczny \l{}a\'ncuch $E[K]\to K_1,K_2,K_3\to Q$}%
\label{app:F-analytic-chain}
% -------------------------------------------------------------------

Niniejsza sekcja zbiera \textbf{pe\l{}ne wyprowadzenie analityczne}
$Q\approx3/2$ z~pierwszych zasad TGP (sesje P47--P51).
Ka\.zdy krok bazuje wy\l{}\k{a}cznie na ca\l{}kach zamkni\k{e}tych
przez funkcj\k{e} wyk\l{}adnicz\k{a} ca\l{}kow\k{a}
$E_1(x)=-\mathrm{Ei}(-x)=\int_x^\infty e^{-t}/t\,dt$.

% . . . . . . . . . . . . . . . . . . . . . . . . . . . . . . . .
\subsubsection{Krok~1: Analityczne wsp\'o\l{}czynniki energii}%
\label{app:F-cn-Ei}
% . . . . . . . . . . . . . . . . . . . . . . . . . . . . . . . .

Energia solitonu Yukawa
$\varphi(r)=1+K e^{-r}/r$ rozwini\k{e}ta w~szereg po~$K$:
\begin{equation}\label{eq:EK-series}
  E[K] = 4\pi\bigl(c_2 K^2 + c_3 K^3 + c_4 K^4 + c_6 K^6\bigr).
\end{equation}
Wsp\'o\l{}czynniki s\k{a} ca\l{}kami analitycznymi (b\l{}\k{a}d~$0{,}00\%$):
\begin{equation}\label{eq:cn-analytic}
\begin{aligned}
  c_2 &= \tfrac{1+\alpha}{2}\,\Phi_2 - \tfrac{1}{2}U_2,&\qquad
  U_2 &= \tfrac{1}{2}e^{-2a},\\
  c_3 &= -\tfrac{\alpha}{2}\,\Phi_3 - \tfrac{2}{3}U_3,&\qquad
  U_3 &= E_1(3a),\\
  c_4 &= \tfrac{\alpha}{2}\,\Phi_4 - \tfrac{1}{4}U_4,&\qquad
  U_4 &= \tfrac{e^{-4a}}{a} - 4E_1(4a),\\
  c_6 &= \tfrac{\lambda}{6}I_6,&\qquad
  I_4 &= U_4,
\end{aligned}
\end{equation}
gdzie $\Phi_n=\int_a^\infty (e^{-r}/r)^n r^2\,dr$.
Dla leptonów: $c_2=112{,}105$, $c_3=-1229{,}03$,
$c_4=19693{,}5$, $c_6=3{,}358\!\times\!10^{-3}$.

% . . . . . . . . . . . . . . . . . . . . . . . . . . . . . . . .
\subsubsection{Krok~2: Twierdzenie o~trzech zerach}%
\label{app:F-three-zeros}
% . . . . . . . . . . . . . . . . . . . . . . . . . . . . . . . .

Funkcja samospójno\'sci $g(K)=E[K]/(4\pi K)-1$ ma
\textbf{dok\l{}adnie trzy zera} $K_1<K_2<K_3$ wtedy i~tylko wtedy, gdy:
\begin{equation}\label{eq:three-zeros}
  c_2>0
  \;\wedge\;
  \max_{K>0}g(K)>0
  \;\wedge\;
  \lambda>0.
\end{equation}
Warunek $c_2>0$ jest spe\l{}niony dla ka\.zdego $\alpha>0$
(kinetyczny wk\l{}ad zawsze dodatni);
$\max g>0$ odpowiada~$\alpha>\alpha_c(a,\lambda)$ -- bifurkacja tworzenia
drugiego i~trzeciego zera; $\lambda>0$ strukturalnie stabilizuje potencja\l{}.

% . . . . . . . . . . . . . . . . . . . . . . . . . . . . . . . .
\subsubsection{Krok~3: $K_1$ -- wzór zamkni\k{e}ty przez resumacj\k{e} Pad\'e~$[2/1]$}%
\label{app:F-K1-pade}
% . . . . . . . . . . . . . . . . . . . . . . . . . . . . . . . .

Aproksymacja Pad\'e~$[2/1]$ funkcji $h(K)=g(K)+1$:
\[
  h(K)\approx\frac{c_2 K+a_2 K^2}{1+b_1 K},
  \qquad
  b_1=-\frac{c_4}{c_3},\quad
  a_2=\frac{c_3^2-c_2 c_4}{c_3}.
\]
Warunek $h(K_1)=1$ daje r\'ownanie kwadratowe, sk\k{a}d:
\begin{equation}\label{eq:K1-P21}
\boxed{
  K_1^{[2/1]}
  = \frac{c_3\!\left[-(c_2 c_3+c_4)+
    \sqrt{4c_3^3+(c_2 c_3-c_4)^2}\right]}
    {2(c_3^2-c_2 c_4)},
}
\end{equation}
gdzie $c_2,c_3,c_4$ z~\eqref{eq:cn-analytic}
-- wzór w~pe\l{}ni analityczny przez $E_1(na)$.
\par\smallskip
\noindent Dok\l{}adno\'s\'c (leptonowy sektor):
\begin{center}
\begin{tabular}{lccc}
\toprule
Metoda & $K_1\;[\times10^{-3}]$ & b\l{}\k{a}d & $Q-\tfrac{3}{2}$ [ppm] \\
\midrule
NLO (kwadr.) & $10{,}021$ & $+1{,}85\%$ & $+362$ \\
NNLO (sz.3) & $9{,}809$ & $-0{,}30\%$ & $-62$ \\
\textbf{Pad\'e~$[2/1]$} & $\mathbf{9{,}837}$
                         & $\mathbf{-0{,}025\%}$ & $\mathbf{-7{,}6}$ \\
Numeryczny & $9{,}839$ & $0\%$ & $+12{,}9$ \\
\bottomrule
\end{tabular}
\end{center}
Trzy zasad encze obserwacje:
(a) $Q=3/2$ jest \textbf{zamkni\k{e}te} mi\k{e}dzy
    $Q_\mathrm{NLO}$ ($+362$~ppm) a~$Q_\mathrm{NNLO}$ ($-62$~ppm);
(b) Pad\'e~$[2/1]$ osi\k{a}ga $-7{,}6$~ppm -- lepsze ni\.z NNLO
    dzi\k{e}ki resummacji geometrycznej mianownika;
(c) b\l{}\k{a}d Pad\'e spada z~rosnacym $\alpha$ -- idealne dla leptonów.

% . . . . . . . . . . . . . . . . . . . . . . . . . . . . . . . .
\subsubsection{Krok~4: $K_3$ przez $E_1(4a)$ i~$E_1(6a)$}%
\label{app:F-K3}
% . . . . . . . . . . . . . . . . . . . . . . . . . . . . . . . .

Dla $K_3\gg1$ dominuj\k{a} cz\l{}ony $c_4 K^4$ (destabilizacja) i~$c_6 K^6$
(stabilizacja). Warunek $g(K_3)=0$:
\begin{equation}\label{eq:K3-Ei}
\boxed{
  K_3 = \sqrt{\frac{3\,I_4}{2\lambda\,I_6}},
}
\end{equation}
gdzie $I_4=e^{-4a}/a-4E_1(4a)$ i~$I_6$ jest w~pe\l{}ni analityczne
przez $E_1(6a)$.
Numerycznie: $K_3^{(E_1)}=34{,}154$ (b\l{}\k{a}d $-0{,}17\%$).

% . . . . . . . . . . . . . . . . . . . . . . . . . . . . . . . .
\subsubsection{Krok~5: $r_{21}$ -- wzór Cardano (historyczny LO)}%
\label{app:F-cardano}
% . . . . . . . . . . . . . . . . . . . . . . . . . . . . . . . .

We wcze\'sniejszym (LO) modelu z~r\'ownaniem kubicznym
$K^3-\varepsilon K^2-2K=a(2\alpha-4)$
($\varepsilon=\frac{4a}{3}\ln\frac{1}{a}$):
\begin{equation}\label{eq:r21-Cardano}
  r_{21}^\mathrm{Card}
  \approx
  \frac{1+\alpha}{\sqrt{3}\,a}\,
  \cos\!\left[\tfrac{1}{3}\arccos\!\!\left(\tfrac{3\sqrt{3}}{4}\,a(2\alpha-4)\right)\right]
  \xrightarrow[\alpha\gg1]{} \frac{\alpha}{\sqrt{2}\,a}.
\end{equation}
B\l{}\k{a}d: $3$--$16\%$ (wzrasta z~$\alpha$);
supersedowany przez~\eqref{eq:K1-P21} ($-0{,}025\%$).

% . . . . . . . . . . . . . . . . . . . . . . . . . . . . . . . .
\subsubsection{Wynik ko\'ncowy i~przewidywanie $\lambda_\mathrm{Koide}$}%
\label{app:F-Q-final}
% . . . . . . . . . . . . . . . . . . . . . . . . . . . . . . . .

Podstawiaj\k{a}c kroki 3--4 do wzoru~\eqref{eq:Q-r21-r31}:
\begin{equation}\label{eq:Q-analytic-final}
  Q_\mathrm{analyt.}
  = Q\bigl(K_1^{[2/1]},\,K_2^\mathrm{Pad},\,K_3^{E_1}\bigr)
  = 1{,}499\,992
  \quad(-7{,}6\ \text{ppm od}\ \tfrac{3}{2}).
\end{equation}
Predykcja $\lambda_\mathrm{Koide}$ (z~warunku $Q=3/2$ przy NNLO $K_1$):
\begin{equation}\label{eq:lam-pred}
  \frac{\lambda_\mathrm{analyt.}}{\lambda_\mathrm{Koide}}
  = 1{,}000\,554
  \quad(+0{,}055\%\ \text{bez dopasowania numerycznego}).
\end{equation}

% . . . . . . . . . . . . . . . . . . . . . . . . . . . . . . . .
\subsubsection{Otwarte problemy analityczne (P52+)}%
\label{app:F-open-P52}
% . . . . . . . . . . . . . . . . . . . . . . . . . . . . . . . .

\begin{enumerate}[label=\textbf{OP-\arabic*},leftmargin=3em]
  \item \textbf{$K_2$ z~pierwszych zasad.}
        Pad\'e~$[2/2]$ daje $K_2=2{,}0304$ ($-0{,}20\%$),
        ale nie jest wzorem zamkni\k{e}tym.
        Cel~P52: znale\'z\'c analityczn\k{a} aproksymacj\k{e}
        $K_\mathrm{max}$ (warunku dg/dK=0) przez $E_1$.

  \item \textbf{Sk\k{a}d $\alpha\approx8{,}54$? (P52: $\alpha_c\approx0$)}
        Bifurkacja trojzerowa nie ma progu w~$\alpha$ -- trzy generacje
        istniej\k{a} dla ka\.zdego $\alpha>0$ (przy $a=0{,}040$, $\lambda_K$).
        $\alpha_\mathrm{Koide}=8{,}54$ jest g\'ornym zerem $Q(\alpha)=3/2$;
        istnieje te\.z dolne zero $\alpha_1\approx2$.
        Pytanie: dlaczego g\'orne zero?
        Hipotezy: energia solitonu; stabilno\'s\'c; kwantowanie topologiczne.

  \item \textbf{Predykcja $a_\Gamma$ z~bifurkacji.}
        Czy warunek $\alpha_c(a_\Gamma,\lambda)=\alpha_\mathrm{Koide}$
        wyznacza $a_\Gamma$ niezale\.znie?

  \item \textbf{$Q=3/2$ dok\l{}adnie i~dwa przeci\k{e}cia $Q(\alpha)=3/2$.}
        Skan numeryczny (P52) przy sta\l{}ym $a=0{,}040$, $\lambda=\lambda_K$:
        \begin{center}
        \begin{tabular}{cccc}
        \toprule
        $\alpha$ & $r_{21}$ & $Q$ & $Q-3/2$ [ppm] \\\midrule
        $0{,}1$ & $18{,}6$ & $1{,}5197$ & $+19693$ \\
        $1{,}0$ & $30{,}7$ & $1{,}5055$ & $+5467$ \\
        $\approx\!2$ & $\approx\!40$ & $\approx\!3/2$ & $\approx0$ \\
        $5{,}0$ & $113$ & $1{,}4938$ & $-6182$ \\
        $\mathbf{8{,}54}$ & $\mathbf{207}$ & $\mathbf{1{,}5000}$ & $+13$ \\
        $15{,}0$ & $404$ & $1{,}5148$ & $+14761$ \\\bottomrule
        \end{tabular}
        \end{center}
        Funkcja $Q(\alpha)$ ma \textbf{minimum} przy $\alpha\approx5$
        ($Q_{\min}\approx1{,}4938$) i~\textbf{dwa przeci\k{e}cia} z~$3/2$:
        dolne $\alpha_1\approx2$ i~g\'orne $\alpha_2\approx8{,}5=\alpha_\mathrm{Koide}$.
        Problem~OP-4 redukuje si\k{e}: \emph{dlaczego natura wybiera g\'orne zero?}

  \item \textbf{Kwarki: mechanizm $Q_\mathrm{PDG}<3/2$.}
        TGP daje $Q\approx1{,}50$--$1{,}53$ dla wszystkich rodzin,
        podczas gdy PDG: $Q(u/c/t)\approx1{,}18$, $Q(d/s/b)\approx1{,}37$.

  \item[\textbf{OP-6}] \textbf{Wyznaczy\'c $a_c$ dok\l{}adnie
        (bifurkacja dost\k{e}pno\'sci Koidego).}
        Warunek $Q_{\min}(a_c)=3/2$ wyznacza progow\k{a} warto\'s\'c
        $a_c\in(0{,}030,\,0{,}040)$, poni\.rej kt\'orej $Q(\alpha)>3/2$
        dla wszystkich $\alpha$ i~rodzina leptonowa Koidego nie mo\.ze
        istnie\'c. Wyznaczenie $a_c$ jest zadaniem P54 (bisekacja).

\end{enumerate}

% . . . . . . . . . . . . . . . . . . . . . . . . . . . . . . . .
\subsubsection{Wyniki P53: precyzyjne zera $Q(\alpha)=3/2$}%
\label{app:F-P53}
% . . . . . . . . . . . . . . . . . . . . . . . . . . . . . . . .

Skrypt \texttt{p53\_double\_crossing.py} (sesja~P53) wyznacza dok\l{}adnie
oba zera $Q(\alpha)=3/2$ przy $a=0{,}040$, $\lambda=\lambda_K$
i~analizuje mechanizm selekcji.

\paragraph{Precyzyjne zera.}
\begin{equation}\label{eq:p53-zeros}
  \alpha_1 = 1{,}5804, \qquad \alpha_2 = 8{,}4731,
  \qquad Q_{\min} = 1{,}4932\ \text{przy}\ \alpha\approx3{,}97.
\end{equation}

\paragraph{Por\'ownanie w\l{}asno\'sci.}
\begin{center}
\begin{tabular}{lcc}
\toprule
W\l{}asno\'s\'c & Dolne zero $\alpha_1$ & G\'orne zero $\alpha_2$ \\\midrule
$\alpha$             & $1{,}5804$  & $8{,}4731$ \\
$r_{21}=K_2/K_1$     & $40{,}27$   & $204{,}52$ \\
$r_{31}=K_3/K_1$     & $831$       & $3443$ \\
Odchy\l{}enie $r_{21}$ od PDG & $80{,}5\%$ & $1{,}1\%$ \\
$dQ/d\alpha$         & $-0{,}0075$ (ind.$=-1$) & $+0{,}0022$ (ind.$=+1$)\\
Odpowiada leptonom?  & \textbf{nie} & \textbf{tak}\\\bottomrule
\end{tabular}
\end{center}

\paragraph{Wnioski P53.}
\begin{enumerate}[(i)]
  \item \textbf{Selekcja przez $r_{21}$}: tylko g\'orne zero daje
        $r_{21}\approx207$ zgodne z~PDG.
        Dolne zero daje $r_{21}=40$ -- nieobserwowan\k{a} hierarchi\k{e} mas.
  \item \textbf{Energia nie selektuje}: $E_\mathrm{tot}(\alpha_1)
        < E_\mathrm{tot}(\alpha_2)$ -- stan przy dolnym zerze jest
        energetycznie ni\.zszy, mimo to nieobserwowany.
        Mechanizm selekcji jest inny ni\.z minimalizacja energii solitonu.
  \item \textbf{Bifurkacja dost\k{e}pno\'sci Koidego}:
        dla $a<a_c$ ($a_c\in(0{,}030,\,0{,}040)$) mamy $Q(\alpha)>3/2$
        dla wszystkich $\alpha$ -- rodzina leptonowa Koidego nie mo\.ze
        istnie\'c.
        Parametr $a_\Gamma=0{,}040>a_c$ jest wi\k{e}c konieczny
        do istnienia $Q=3/2$.
  \item \textbf{Indeksy topologiczne}: $\mathrm{ind}(\alpha_1)=-1$,
        $\mathrm{ind}(\alpha_2)=+1$, suma $=0$ (zgodnie z~oczekiwaniem
        dla krzywej symetrycznie przekraczaj\k{a}cej $3/2$).
\end{enumerate}

% . . . . . . . . . . . . . . . . . . . . . . . . . . . . . . . .
\subsubsection{Wyniki P54: precyzyjne wyznaczenie $a_c$}%
\label{app:F-P54}
% . . . . . . . . . . . . . . . . . . . . . . . . . . . . . . . .

Skrypt \texttt{p54\_ac\_bifurcation.py} (sesja~P54) wyznacza $a_c$ przez
bisekacj\k{e} na waruneku $Q_{\min}(a_c)=3/2$, z~dok\l{}adno\'sci\k{a}
$\Delta a_c < 10^{-8}$.

\paragraph{Wynik bisekacji.}
\begin{equation}\label{eq:p54-ac}
  \boxed{a_c = 0{,}038382}
  \qquad
  \alpha^* = 4{,}0250,
  \qquad
  r_{21}^* = 94{,}00.
\end{equation}

\paragraph{Wlasciwosci w punkcie bifurkacji.}
\begin{center}
\begin{tabular}{lc}
\toprule
Parametr & Warto\'s\'c \\\midrule
$a_c$ (bisekacja) & $0{,}038382$ \\
$\alpha^*$ (punkt dotyku $Q=3/2$) & $4{,}0250$ \\
$K_1^*$ & $0{,}01911$ \\
$K_2^*$ & $1{,}7962$ \\
$K_3^*$ & $32{,}789$ \\
$r_{21}^* = K_2^*/K_1^*$ & $94{,}00$ \\
$a_\Gamma/a_c$ & $1{,}042\ (+4{,}2\%)$ \\
Eksponent bifurkacji $\beta$ & $\approx0{,}56\approx\tfrac{1}{2}$\\\bottomrule
\end{tabular}
\end{center}

\paragraph{Struktura bifurkacji.}
Bifurkacja $Q_{\min}(a)=3/2$ jest bifurkacj\k{a} siod\l{}owo-w\k{e}z\l{}ow\k{a}
(\emph{fold/saddle-node}): w~punkcie $a_c$ dwa zera $Q(\alpha)=3/2$
rodz\k{a} si\k{e} z~jednego zdegenerowanego punktu $\alpha^*$, a~odleg\l{}o\'s\'c
mi\k{e}dzy zerami ro\'snie jak
\begin{equation}\label{eq:p54-beta}
  \delta\alpha \equiv \alpha_2 - \alpha_1
  \approx C\,(a-a_c)^{\beta},
  \qquad
  \beta \approx \tfrac{1}{2},\quad C\approx258.
\end{equation}

\paragraph{Wnioski P54.}
\begin{enumerate}[(i)]
  \item $a_c=0{,}038382$ -- precyzyjny pr\'og istnienia rodziny Koidego.
        Dla $a<a_c$ spe\l{}nienie $Q=3/2$ jest strukturalnie niemo\.zliwe.
  \item $a_\Gamma=0{,}040$ jest $4{,}2\%$ powy\.rej progu;
        wyznaczone przez warunek $r_{21}=r_{21}^{\mathrm{PDG}}$
        (patrz~P55).
  \item Postaci analitycznej $a_c$ nie znaleziono (OP-7):
        nie jest r\'owne $\pi/N$, $e^{-N}$ ani $\lambda^{1/4}$
        dla prostych $N$.
        Warto\'s\'c $a_c/\lambda^{1/4}=0{,}794$ nie redukuje si\k{e}
        do \.zadnej rozpoznanej sta\l{}ej matematycznej.
\end{enumerate}

% . . . . . . . . . . . . . . . . . . . . . . . . . . . . . . . .
\subsubsection{Wyniki P55: uklad rownan wyznacza parametry TGP}%
\label{app:F-P55}
% . . . . . . . . . . . . . . . . . . . . . . . . . . . . . . . .

Skrypt \texttt{p55\_r21\_condition.py} (sesja~P55) weryfikuje hipotez\k{e}~OP-8:
czy nak\l{}adaj\k{a}c jednocze\'snie $Q=3/2$ i~$r_{21}=r_{21}^{\mathrm{PDG}}$
mo\.zna odtworzy\'c $a_\Gamma$ i~$\alpha_\mathrm{Koide}$ \emph{bez dopasowania}.

\paragraph{Uklad rownan.}
Dla danego $\lambda=\lambda_K$, dwa warunki:
\begin{equation}\label{eq:p55-system}
  Q(\alpha,\,a,\,\lambda_K) = \tfrac{3}{2},
  \qquad
  r_{21}(\alpha,\,a,\,\lambda_K) = 206{,}77,
\end{equation}
wyznaczaj\k{a} \textbf{uk\l{}ad dw\'och r\'owna\'n z~dwoma nieznanymi}
$(a,\,\alpha)$.
Gorny \'zero $Q(\alpha)=3/2$ daje $\alpha=\alpha_2(a)$, wi\k{e}c~(\ref{eq:p55-system})
redukuje si\k{e} do jednego r\'ownania $r_{21}(\alpha_2(a),a)=206{,}77$,
ktore rozwiazujemy bisekcj\k{a}.

\paragraph{Wynik.}
\begin{equation}\label{eq:p55-result}
  a^* = 0{,}040049,\qquad \alpha_2^* = 8{,}5612.
\end{equation}

\begin{center}
\begin{tabular}{lccc}
\toprule
Wielko\'s\'c & Z uk\l{}adu~(\ref{eq:p55-system}) & Warto\'s\'c PDG/target & R\'o\.znica\\\midrule
$a^*$                   & $0{,}040049$ & $a_\Gamma=0{,}040000$ & $+0{,}12\%$ \\
$\alpha_2^*$            & $8{,}5612$   & $\alpha_K=8{,}5445$   & $+0{,}20\%$ \\
$r_{31}^*$ (predykcja!) & $3477{,}5$   & $r_{31}^{\mathrm{PDG}}=3477{,}2$ & $\mathbf{+0{,}01\%}$\\\bottomrule
\end{tabular}
\end{center}

\paragraph{Wnioski P55.}
\begin{enumerate}[(i)]
  \item \textbf{OP-8 rozwi\k{a}zany}: warunek $r_{21}(a,\alpha_2)=206{,}77$
        wyznacza $a_\Gamma$ z~dok\l{}adno\'sci\k{a} $0{,}12\%$.
  \item \textbf{Predykcja $r_{31}$}: stosunek mas $m_\tau/m_e$
        wychodzi jako $r_{31}^*=3477{,}5$ z~dok\l{}adno\'sci\k{a}~$0{,}01\%$
        \emph{bez u\.zycia go jako warunku}.
  \item \textbf{Brak wolnych parametr\'ow w~sektorze leptonowym}:
        ukl\'ad~(\ref{eq:p55-system}) + warunek $r_{31}$ tworzy uk\l{}ad
        tr\'och r\'owna\'n z~trzema nieznanymi $(a,\alpha,\lambda)$:
        \begin{equation}\label{eq:p55-full}
          Q=\tfrac{3}{2},\qquad r_{21}=206{,}77,\qquad r_{31}=3477{,}2
          \quad\Longrightarrow\quad
          (a_\Gamma,\,\alpha_K,\,\lambda_K).
        \end{equation}
        Sek\-tor leptonowy TGP nie zawiera wolnych parametr\'ow --
        wszystkie wyp\l{}ywaj\k{a} z~mas PDG i~geometrii solitonu.
  \item \textbf{OP-9 rozwi\k{a}zany w~P56}: patrz poni\.zej.
\end{enumerate}

% . . . . . . . . . . . . . . . . . . . . . . . . . . . . . . . .
\subsubsection{Wyniki P56: wyznaczenie $\lambda$ z~uk\l{}adu pe\l{}nego}%
\label{app:F-P56}
% . . . . . . . . . . . . . . . . . . . . . . . . . . . . . . . .

Skrypt \texttt{p56\_full\_system.py} (sesja~P56) weryfikuje uk\l{}ad~(\ref{eq:p55-full})
przy $\lambda$ jako trzeciej niewiadomej.

\paragraph{Niezale\.zno\'s\'c $K_1$ od $\lambda$ (odkrycie strukturalne).}
Numerycznie stwierdzono, \.ze $K_1$ (zero g(K) przy ma\l{}ym $K$)
jest \textbf{dok\l{}adnie niezale\.zne od $\lambda$}:
\begin{center}
\begin{tabular}{ccccc}
\toprule
$\lambda/\lambda_K$ & $K_1$ & $K_2$ & $K_3$ & $r_{31}$\\\midrule
$0{,}5$ & $0{,}009833$ & $2{,}032$ & $48{,}36$ & $4918$ \\
$1{,}0$ & $0{,}009833$ & $2{,}033$ & $34{,}20$ & $3477$ \\
$2{,}0$ & $0{,}009833$ & $2{,}036$ & $24{,}18$ & $2459$ \\\bottomrule
\end{tabular}
\end{center}
Jedynie $K_3\propto\lambda^{-1/2}$ skaluje z~$\lambda$ (zgodnie z~wzorem~K$_3$-Ei).
$K_1$ i~(przybli\.zone) $K_2$ s\k{a} wyznaczone przez $\alpha$ i~$a$ bez udzia\l{}u $\lambda$.

\paragraph{Wynik pe\l{}nej iteracji.}
\begin{equation}\label{eq:p56-result}
  a^* = 0{,}04009\ (+0{,}23\%),\quad
  \alpha^* = 8{,}568\ (+0{,}28\%),\quad
  \lambda^* = \lambda_K\cdot(1+158\,\mathrm{ppm}).
\end{equation}
Uk\l{}ad r\'owna\'n spe\l{}niony z~rezyduum $Q-3/2=+17\,\mathrm{ppm}$
(pochodzi z~przybli\.zenia $K_3$-Ei: $K_3^{\mathrm{Ei}}=34{,}195$ vs.
$K_3^{\mathrm{num}}=34{,}154$, r\'o\.znica $0{,}12\%$).

\paragraph{Wnioski P56.}
\begin{enumerate}[(i)]
  \item $\lambda$ wyznaczone z~warunku $r_{31}=3477{,}2$ przy sta\l{}ym
        $(a^*,\alpha^*)$ z~P55: $\lambda^*=\lambda_K\times1{,}000158$
        ($+158\,\mathrm{ppm}$).
  \item Sektor leptonowy TGP spe\l{}nia~(\ref{eq:p55-full})
        z~dok\l{}adno\'sci\k{a} $\sim0{,}2\%$ dla $(a,\alpha)$ i~$\sim160\,\mathrm{ppm}$
        dla $\lambda$ -- bez \.zadnego wolnego parametru.
  \item Otwarty~OP-10: dlaczego $K_1$ jest dok\l{}adnie niezale\.zne od $\lambda$?
        Mo\.ze wynika\'c z~faktu, \.ze dla ma\l{}ego $K$
        soliton jest p\l{}ytki ($K\exp(-r)/r\ll1$) i~cz\l{}on
        $\lambda(\phi-1)^6\ll(\phi-1)^3$ jest zaniedbywalny.
  \item OP-11 \emph{zdiagnozowany} w~P57: patrz poni\.zej.
\end{enumerate}

% . . . . . . . . . . . . . . . . . . . . . . . . . . . . . . . .
\subsubsection{Diagnoza P57: rozbie\.zno\'s\'c $\lambda$ i~sp\'ojno\'s\'c wynik\'ow}%
\label{app:F-P57}
% . . . . . . . . . . . . . . . . . . . . . . . . . . . . . . . .

Skrypt \texttt{p57\_lambda\_discrepancy.py} (sesja~P57) analizuje \'zr\'od\l{}o
rozbie\.zno\'sci mi\k{e}dzy dwiema niezale\.znymi metodami wyznaczania $\lambda$:
\begin{align}
  \text{P50: } &\quad Q(K_1^{[2/1]},\,K_2^{[2/2]},\,K_3^{\mathrm{Ei}}(\lambda))
               = \tfrac{3}{2}
               \;\Rightarrow\; \lambda_\mathrm{analyt}/\lambda_K = 1{,}000554,
  \\
  \text{P56: } &\quad K_3^{\mathrm{Ei}}(\lambda)/K_1^{\mathrm{num}}
               = r_{31}^{\mathrm{PDG}}
               \;\Rightarrow\; \lambda_{r_{31}}/\lambda_K = 1{,}000158.
\end{align}
Rozk\l{}ad r\'o\.znicy $\Delta\lambda/\lambda_K\approx396\,\mathrm{ppm}$:
\begin{itemize}
  \item \textbf{G\l{}\'owny wk\l{}ad}: b\l{}\k{a}d $K_2^{[2/2]}=-0{,}20\%$
        (Pad\'e~$[2/2]$) propaguje si\k{e} przez formu\l{}\k{e} Koide
        do $\delta\lambda\approx+3621\,\mathrm{ppm}$ w~metodzie~P50.
  \item \textbf{Wk\l{}ad $K_1$}: b\l{}\k{a}d $K_1^{[2/1]}=-0{,}025\%$
        daje $\approx-145\,\mathrm{ppm}$ (cz\k{e}\'sciowo kompensuje).
  \item \textbf{Wk\l{}ad $K_3$}: $K_3^{\mathrm{Ei}}/K_3^{\mathrm{num}}-1
        \approx-10\,\mathrm{ppm}$ (zaniedbywalny).
\end{itemize}
\textbf{Wniosek P57}: OP-11 jest \emph{pochodn\k{a}} OP-1.
Dok\l{}adny wzór na $K_2$ (OP-1) automatycznie usunie rozbie\.zno\'s\'c.
Do czasu rozwi\k{a}zania OP-1 lepsz\k{a} estymacj\k{a} $\lambda$ jest
metoda P56 ($r_{31}$, niezale\.zna od $K_2$).

\paragraph{Niespójno\'sci techniczne (do naprawy w~P58).}
W~tabeli parametr\'ow leptonowych warto\'sci $K_i$ s\k{a} zaokr\k{a}glone
niezale\.znie, co powoduje $r_{31}^{\mathrm{tab}}=K_3/K_1=3471{,}3\neq3477{,}2$.
Dok\l{}adne warto\'sci ($10$ cyfr) zostan\k{a} obliczone w~P58.

\paragraph{Stan otwarty (OP-1, OP-2, OP-7, OP-10, OP-12, OP-13).}
\begin{description}
  \item[OP-1] $K_2$ z~pierwszych zasad -- dominuj\k{a}cy problem otwarty;
              blokuje OP-4 i~OP-11.
  \item[OP-2] Selekcja g\'ornego zera $\alpha_2$ -- mechanizm fizyczny
              nieznany; selekcja energetyczna wykluczona (P53).
  \item[OP-7] Postac analityczna $a_c$ -- brak zwi\k{a}zku z~$\pi$, $e$,
              $\lambda^{1/4}$.
  \item[OP-10] $K_1$ niezale\.zne od $\lambda$ -- hipoteza perturbacyjna:
               soliton przy $K_1\ll1$ jest p\l{}ytki, wk\l{}ad
               $\lambda(\phi-1)^6$ zaniedbywalny.
  \item[OP-12] B\l{}\k{a}d $K_3^{\mathrm{Ei}}$ -- P58 wykry\l{}o:
               $K_3^{\mathrm{num}}=34{,}249$ vs $K_3^{\mathrm{Ei}}=34{,}154$
               przy $a=0{,}040$; wzgl\k{e}dny b\l{}\k{a}d $-0{,}28\%$.
               Nale\.zy zbada\'c, sk\k{a}d wynika ta rozbie\.zno\'s\'c.
  \item[OP-13] \textbf{Rozwiazany (P59):} zbiezno\'s\'c siatki --
               rz\k{a}d $p\approx2{,}12$; $N^*=500$
               (b\l{}\k{a}d $K_2<16\,\mathrm{ppm}$);
               $N=5000$ daje $\delta K_2=0{,}2\,\mathrm{ppm}$.
               Obserwowane $\Delta K_2=0{,}12\%$ by\l{}o zmian\k{a}
               punktu A$\to$B, nie b\l{}\k{e}dem siatki.
\end{description}

% . . . . . . . . . . . . . . . . . . . . . . . . . . . . . . . .
\subsubsection{Wyniki P58: precyzyjna tabela $K_i$ (odswiezenie)}%
\label{app:F-P58}
% . . . . . . . . . . . . . . . . . . . . . . . . . . . . . . . .

Skrypt \texttt{p58\_precise\_Ki.py} (sesja~P58) ponownie oblicza $K_i$
z~siatk\k{a} $N=5000$ (zamiast $N=3000$) dla trzech punkt\'ow i~ujawnia
dwa nowe problemy numeryczne.

\paragraph{Nowa tabela punkt\'ow referencyjnych.}
\begin{center}
\begin{tabular}{lccc}
\toprule
Parametr & Punkt~A (hist.) & Punkt~B (ref.) & PDG\\\midrule
$a_\Gamma$        & $0{,}04000$ & $\mathbf{0{,}04005}$ & --- \\
$\alpha$          & $8{,}5445$  & $\mathbf{8{,}5612}$  & --- \\
$K_1$             & $0{,}009839$& $\mathbf{0{,}009833}$& --- \\
$K_2$             & $2{,}0344$  & $\mathbf{2{,}033219}$& --- \\
$K_3^{\mathrm{Ei}}$ & $34{,}154$& $\mathbf{34{,}195}$  & --- \\
$r_{21}=K_2/K_1$  & $206{,}52^*$ & $\mathbf{206{,}769}$ & $206{,}768$ \\
$r_{31}=K_3/K_1$  & $3471$     & $\mathbf{3477{,}47}$ & $3477{,}23$ \\
$Q-3/2$           & $+156\,\mathrm{ppm}$ & $\mathbf{-1{,}3\,\mathrm{ppm}}$ & $0$\\\bottomrule
\end{tabular}
\end{center}
{\small $^*$Warto\'s\'c przy $N=5000$; przy $N=3000$ (P40): $r_{21}=206{,}77$
(b\l{}\k{a}d siatki $\approx0{,}12\%$).}

\paragraph{Odkrycia P58.}
\begin{enumerate}[(i)]
  \item \textbf{B\l{}\k{a}d $K_3^{\mathrm{Ei}}=-0{,}28\%$}:
        pe\l{}na numeryka daje $K_3^{\mathrm{num}}(A)=34{,}249$,
        a~wz\'or $K_3^{\mathrm{Ei}}=34{,}154$ -- r\'o\.znica $-2791\,\mathrm{ppm}$.
        Hipoteza: b\l{}\k{a}d pochodzi z~przybli\.zenia ca\l{}ki $I_6$
        (dobiera si\k{e} asymptotycznie, ale s\l{}abo przy $a=0{,}040$).
  \item \textbf{Zbiezno\'s\'c siatki}:
        $K_2(N=3000)=2{,}0344\neq K_2(N=5000)=2{,}0320$
        -- ta $0{,}12\%$ r\'o\.znica przek\l{}ada si\k{e}
        na $\Delta r_{21}=0{,}25$.
        \emph{Uwaga (P59):} r\'o\.znica ta pochodzi ze zmiany punktu
        A$\to$B ($\Delta a$, $\Delta\alpha$), nie z~b\l{}\k{e}du siatki;
        rzeczywisty b\l{}\k{a}d siatki przy tym samym punkcie
        wynosi zaledwie $0{,}4\,\mathrm{ppm}$.
  \item \textbf{Punkt~B jako nowa referencja}:
        przy $a=0{,}040049$, $\alpha=8{,}5612$, $\lambda=\lambda_K$
        (z~P55) i~$N=5000$: $Q=-1{,}3\,\mathrm{ppm}$,
        $r_{21}\approx r_{21}^{\mathrm{PDG}}$,
        $r_{31}\approx r_{31}^{\mathrm{PDG}}$ -- samospojny.
\end{enumerate}

% . . . . . . . . . . . . . . . . . . . . . . . . . . . . . . . .
\subsubsection{Wyniki P59: zbieznosc siatki numerycznej (OP-13 $\to$ rozwiazany)}%
\label{app:F-P59}
% . . . . . . . . . . . . . . . . . . . . . . . . . . . . . . . .

Skrypt \texttt{p59\_grid\_convergence.py} (sesja~P59) bada zbieznosc
numeryczna metody brentq na $g(K)=0$ w~funkcji liczby punktow
siatki $N\in\{500,1000,2000,3000,5000,8000,10000\}$.

\paragraph{Metoda.}
Dla logarytmicznej siatki radialnej $r_j=a(R_{\max}/a)^{j/(N-1)}$
ca\l{}kowanie trapezoidalne daje b\l{}\k{a}d $\mathcal{O}(h^2)$,
gdzie efektywny krok $h\propto 1/N$.
Oczekuje si\k{e} wi\k{e}c $K(N)=K_\infty+c/N^p$ z~$p\approx 2$.

\paragraph{Rz\k{a}d zbieznosci.}
Wyznaczony z~trojki $N\in\{2000,5000,10000\}$:
\begin{equation}
  p_{K_1} = p_{K_2} \approx 2{,}12,
  \label{eq:app:F-P59-order}
\end{equation}
zgodny z~oczekiwanym $p=2$ dla metody trapezoidalnej.

\paragraph{Ekstrapolacja Richardsona $N\to\infty$.}
Z~pary $(N=5000,\,N=10000)$ i~wyznaczonego $p=2{,}12$:
\begin{align}
  K_1^\infty &= 0{,}009\,833\,313, \notag\\
  K_2^\infty &= 2{,}033\,219\,544, \notag\\
  r_{21}^\infty &= 206{,}7685,\quad
    \Delta r_{21}^{\infty-\mathrm{PDG}} = {+}0{,}000\,24, \notag\\
  Q^\infty - \tfrac{3}{2} &= -1{,}28\,\mathrm{ppm}.
  \label{eq:app:F-P59-rich}
\end{align}

\paragraph{Tabela b\l{}\k{e}dow wzgl\k{e}dnych.}
\begin{center}
\begin{tabular}{rrrrrr}
\hline
$N$ & $\delta K_1\,[\mathrm{ppm}]$ & $\delta K_2\,[\mathrm{ppm}]$
    & $\delta r_{21}\,[\mathrm{ppm}]$ & $\delta Q\,[\mathrm{ppm}]$ \\
\hline
 500 & $-107$ & $-16$ & $+91$ & $-5{,}36$ \\
1000 & $-27$  & $-4$  & $+23$ & $-1{,}34$ \\
2000 & $-7$   & $-1$  & $+6$  & $-0{,}33$ \\
3000 & $-3$   & $-0{,}4$ & $+3$  & $-0{,}15$ \\
\textbf{5000} & $\mathbf{-1{,}0}$ & $\mathbf{-0{,}2}$ & $\mathbf{+0{,}9}$ & $\mathbf{-0{,}05}$ \\
8000 & $-0{,}4$ & $-0{,}1$ & $+0{,}3$ & $-0{,}02$ \\
\hline
\end{tabular}
\end{center}

\paragraph{Prog dokladnosci.}
Kryterium $|\delta K_2|<100\,\mathrm{ppm}$ (dok\l{}adno\'s\'c $0{,}01\%$)
jest spe\l{}nione juz dla $N^*=500$.
Przy $N=5000$ b\l{}\k{a}d wynosi jedynie $0{,}2\,\mathrm{ppm}$ --
wyniki P58 sa wiarygodne numerycznie.

\paragraph{Wyjasnienie ``b\l{}\k{e}du $0{,}12\%$'' z~P58.}
Spostrzezona w~P58 roznica
$K_2(N=3000)=2{,}0344\neq K_2(N=5000)=2{,}0320$ wynoszaca $0{,}12\%$
\emph{nie jest b\l{}\k{e}dem siatki}.
Faktyczna roznica $K_2(N=3000)$ vs $K_2(N=5000)$
przy tym samym punkcie parametrycznym wynosi $0{,}4\,\mathrm{ppm}$.
Obserwowane $0{,}12\%$ pochodzi ze zmiany punktu
A ($a=0{,}040$, $\alpha=8{,}5445$) na punkt
B ($a=0{,}040049$, $\alpha=8{,}5612$) -- co stanowi
fizyczna poprawke wymuszona przez OP-8.

\paragraph{Wnioski P59.}
\begin{enumerate}
  \item OP-13 rozwiazany: siatka $N=5000$ jest numerycznie wystarczajaca
        ($\delta K_i<2\,\mathrm{ppm}$).
  \item Dominujacym zrodlem niedokladnosci tabelarycznych nie jest siatka,
        lecz blad analityczny $K_3^{\mathrm{Ei}}$ (OP-12, $-2791\,\mathrm{ppm}$)
        i~blad Pade $K_2$ (OP-1, $\sim-2000\,\mathrm{ppm}$).
  \item Priorytety: OP-12 (P60) i~OP-10 (P61).
\end{enumerate}

% . . . . . . . . . . . . . . . . . . . . . . . . . . . . . . . .
\subsubsection{Wyniki P60: diagnoza bledu $K_3^{\mathrm{Ei}}$ (OP-12)}%
\label{app:F-P60}
% . . . . . . . . . . . . . . . . . . . . . . . . . . . . . . . .

Skrypt \texttt{p60\_K3\_correction.py} (sesja~P60) diagnozuje
zrodlo bledu $-2792\,\mathrm{ppm}$ formu\l{}y
$K_3^{\mathrm{Ei}}=\sqrt{3I_4/(2\lambda I_6)}$ wzgledem
pelnej numeryki $K_3^{\mathrm{num}}$ (brentq na $g(K_3)=0$).

\paragraph{\'Zr\'od\l{}o b\l{}\k{e}du -- analiza sk\l{}adnik\'ow.}
Formu\l{}a $K_3^{\mathrm{Ei}}$ pochodzi z~bilansu sk\l{}adnik\'ow
$K^4$ i~$\lambda K^6$ energii:
\begin{equation}
  -\frac{K^4}{4}I_4 + \frac{\lambda K^6}{6}I_6 = 0
  \quad\Longrightarrow\quad
  K_3^{\mathrm{Ei}} = \sqrt{\frac{3I_4}{2\lambda I_6}}.
  \label{eq:K3Ei-balance}
\end{equation}
Zak\l{}ada si\k{e} przy tym:
\begin{align}
  \int_a^\infty (\varphi-1)^6\, r^2\,\mathrm{d}r
    &\approx K^6 \int_a^\infty \frac{e^{-6r}}{r^4}\,\mathrm{d}r
    = K^6 I_6, \label{eq:K3Ei-lam-approx}\\
  \int_a^\infty \frac{\varphi^4-1}{4}\,r^2\,\mathrm{d}r
    &\approx \frac{K^4}{4}\int_a^\infty \frac{e^{-4r}}{r^2}\,\mathrm{d}r
    = \frac{K^4}{4} I_4. \label{eq:K3Ei-qrt-approx}
\end{align}
Weryfikacja numeryczna przy $K=K_3^{\mathrm{num}}=34{,}249$ (Punkt~A):
\begin{center}
\begin{tabular}{lcc}
\hline
Aproksymacja & B\l{}\k{a}d \\
\hline
$(\varphi-1)^6 \approx K^6 e^{-6r}/r^6$ & \textbf{0\,ppm} (\checkmark) \\
$\varphi^4 \approx K^4 e^{-4r}/r^4$     & \textbf{$+12\,528$\,ppm} ($\times$) \\
\hline
\end{tabular}
\end{center}
Ca\l{}ka $\lambda$-cz\l{}onu jest \emph{dok\l{}adna} --
$({\varphi-1})^6 = K^6 e^{-6r}/r^6$ \emph{\'sci\'sle},
poniewa\.z $\varphi-1 = K e^{-r}/r$ z~definicji profilu solitonu.
Natomiast aproksymacja $\varphi^4 \approx K^4 e^{-4r}/r^4$
pomija cz\l{}on $1$ w~$(1 + Ke^{-r}/r)^4$, generujac b\l{}\k{a}d
$+12\,528\,\mathrm{ppm}$ ($1{,}25\%$) w~ca\l{}ce kwartetu.

\paragraph{Mechanizm.}
Kwartety sa przeszacowane o~$1{,}25\%$, zatem r\'ownanie bilansu
\eqref{eq:K3Ei-balance} wymaga wi\k{e}kszego $K_3$ celem
zbilansowania wiekszego kwartetu przez silniejszy $\lambda$-cz\l{}on.
St\k{a}d $K_3^{\mathrm{num}} > K_3^{\mathrm{Ei}}$:
rzeczywiste $K_3$ jest \emph{zawy\.zone} wzgl\k{e}dem $K_3^{\mathrm{Ei}}$. \smallskip

Przy uwzgl\k{e}dnieniu tak\.ze sk\l{}adnik\'ow kinetycznego i~kubicznego
(kt\'ore cz\k{e}\'sciowo kompensuja ten efekt),
netto: $K_3^{\mathrm{num}}/K_3^{\mathrm{Ei}} = 1 + 2792\,\mathrm{ppm}$.

\paragraph{Korekcja metoda Newtona.}
Jeden krok Newtona na $g(K_3)=0$ startujac od $K_3^{\mathrm{Ei}}$:
\begin{equation}
  K_3^{(1)} = K_3^{\mathrm{Ei}} - \frac{g(K_3^{\mathrm{Ei}})}{g'(K_3^{\mathrm{Ei}})}
  = 34{,}2502,\quad
  \delta K_3^{(1)} = +28\,\mathrm{ppm}.
  \label{eq:K3-newton}
\end{equation}
Po dwoch krokach: $K_3^{(2)} = K_3^{\mathrm{num}}$ z~dokladnoscia maszynowa.

\paragraph{Wplyw na wielko\'sci fizyczne.}
\begin{center}
\begin{tabular}{lccc}
\hline
Wariant & $K_3$ & $r_{31}$ & $Q - 3/2$ \\
\hline
$K_3^{\mathrm{Ei}}$ & 34{,}154 & 3473{,}3 & $+272\,\mathrm{ppm}$ \\
$K_3^{\mathrm{num}}$& 34{,}249 & 3483{,}0 & $-359\,\mathrm{ppm}$ \\
PDG                 & --       & 3477{,}2 & -- \\
\hline
\end{tabular}
\end{center}
B\l{}\k{a}d $K_3^{\mathrm{Ei}}$ ($-2792\,\mathrm{ppm}$) przek\l{}ada si\k{e} na
$\Delta r_{31} = -9{,}7$ (pomi\k{e}dzy wariantami) i~$\Delta Q \approx 630\,\mathrm{ppm}$.
\.Zaden z~wariant\'ow nie reprodukuje $r_{31}^{\mathrm{PDG}}$ dok\l{}adnie
przy parametrach Punktu~A -- co jest odzwierciedleniem OP-12 jako symptomu
g\l{}ebszego braku samospojno\'sci (por. OP-3).

\paragraph{Wnioski P60.}
\begin{enumerate}
  \item OP-12 zdiagnozowany: zr\'od\l{}em b\l{}\k{e}du $K_3^{\mathrm{Ei}}$
        jest aproksymacja $\varphi^4\approx K^4 e^{-4r}/r^4$
        z~b\l{}\k{e}dem $+12\,528\,\mathrm{ppm}$.
  \item Rozwiazanie techniczne: 2-krokowa iteracja Newtona
        z~$K_3^{\mathrm{Ei}}$ daje $K_3^{\mathrm{num}}$ (dokladnosc maszynowa).
  \item Blad $K_3^{\mathrm{Ei}}$ wplywa na $r_{31}$ o~$-2792\,\mathrm{ppm}$
        i~na $Q$ o~$\sim 630\,\mathrm{ppm}$.
  \item Ca\l{}ka $\lambda$-cz\l{}onu $I_6$ jest analitycznie dok\l{}adna --
        nie wymaga poprawy.
  \item Kolejny priorytet: OP-10 (K\textsubscript{1} niezale\.zne od $\lambda$) -- P61.
\end{enumerate}

% . . . . . . . . . . . . . . . . . . . . . . . . . . . . . . . .
\subsubsection{Wyniki P61: mechanizm niezale\.zno\'sci $K_1$ od $\lambda$ (OP-10 $\to$ rozwiazany)}%
\label{app:F-P61}
% . . . . . . . . . . . . . . . . . . . . . . . . . . . . . . . .

Skrypt \texttt{p61\_K1\_lambda\_independence.py} (sesja~P61)
wyja\'snia strukturalna przyczyne obserwowanej w~P56 niezale\.zno\'sci
$K_1$ od~$\lambda$.

\paragraph{Tlumienie cz\l{}onu $\lambda$ przy $K_1$.}
Wk\l{}ad cz\l{}onu $\lambda$ do warunku $g(K_1)=0$ wynosi:
\begin{equation}
  \frac{E_\lambda(K_1)}{4\pi K_1}
  \approx \frac{\lambda K_1^5 I_6}{6}
  \approx 5{,}5\times10^{-6}\times(0{,}01)^5\times\frac{3685}{6}
  \approx 3\times10^{-13},
  \label{eq:K1-lam-suppression}
\end{equation}
podczas gdy sam warunek $g(K_1)=0$ wymaga, aby ta wielko\'s\'c
by\l{}a rz\k{e}du jedynki.
Tlumienie wynosi \textbf{13 rz\k{e}d\'ow wielko\'sci}.
Numerycznie: $g_\lambda(K_1)/1 = 3{,}1\times10^{-13}\,\mathrm{ppb}$.

\paragraph{Hierarchia czulosci $K_i$ na $\lambda$.}
\begin{center}
\begin{tabular}{lcrr}
\hline
$K_i$ & $K_i$ & $\partial\log K_i/\partial\log\lambda$ & Interpretacja \\
\hline
$K_1$ & $0{,}009833$ & $-3{,}4\times10^{-13}$ & ca\l{}kowita niezale\.zno\'s\'c \\
$K_2$ & $2{,}0332$   & $+0{,}00144$           & s\l{}aba zale\.zno\'s\'c ($+998\,\mathrm{ppm}$ na $\times2\lambda$) \\
$K_3$ & $34{,}29$    & $-0{,}499\approx-\tfrac{1}{2}$ & $K_3\propto\lambda^{-1/2}$ (silna) \\
\hline
\end{tabular}
\end{center}
Weryfikacja: $K_1(\lambda=0{,}1\,\lambda_K) = K_1(10\,\lambda_K)$
z~dokladnoscia maszynowa (14 cyfr znacych).

\paragraph{R\'ownanie $K_1$ bez $\lambda$ ($g_0=1$).}
Warunek $g_0(K_1)=1$ (bez cz\l{}onu $\lambda$) ma postac:
\begin{equation}
  \frac{E_{\rm kin} + E_{V_3} + E_{V_4}}{4\pi K_1} = 1,
  \label{eq:K1-nolam}
\end{equation}
gdzie $E_{V_3}, E_{V_4}$ to wk\l{}ady potencja\l{}u kubicznego
i~kwartycznego bez $\lambda$.
Numerycznie: $K_1(\lambda=0) = K_1(\lambda_K)$ identycznie --
r\'ownanie~\eqref{eq:K1-nolam} jest tym samym co pe\l{}ne r\'ownanie.

\paragraph{Wz\'or wiodacy.}
Perturbacyjnie w~$K_1\ll1$:
\begin{equation}
  K_1^{\rm lead} = \frac{1}{A_1},\quad
  A_1 = (1+\alpha)\,C_2 - \tfrac{1}{2}I_2,
  \label{eq:K1-leading}
\end{equation}
gdzie $C_2 = \int_a^\infty \tfrac{e^{-2r}(r+1)^2}{2r^2}\,\mathrm{d}r$
i~$I_2 = e^{-2a}/2$.
Dla Punktu~B: $A_1=112{,}156$, $K_1^{\rm lead}=0{,}00892$
(b\l{}\k{a}d $-9{,}33\%$ vs $K_1^{\rm num}$).
Wolna zbiezno\'s\'c szeregu perturbacyjnego wynika z~tego,
\.ze kolejne pop rawki $\sim 2K_1 I_3/3$ nie sa male
w~stosunku do $A_1$.

\paragraph{Interpretacja fizyczna: dekuplowanie skal.}
Trzy zera $g(K)=0$ odpowiadaja trzem niezale\.znym skalom energetycznym:
\begin{description}
  \item[$K_1$ (elektron):] skala wyznaczona przez geometri\k{e} solitonu --
    wsp\'o\l{}czynnik sp\k{e}\.zenia $\alpha$ i~rozmiar $a_\Gamma$.
    $\lambda$-cz\l{}on jest supresowany przez $K_1^5\sim10^{-10}$.
  \item[$K_3$ (taon):] skala wyznaczona przez \l{}amanie symetrii $\lambda$:
    $K_3\propto1/\sqrt{\lambda}$.
    Dominuje $\lambda$-cz\l{}on przez $K_3^5\sim10^9$.
  \item[$K_2$ (mion):] skala po\'srednia;
    s\l{}aba zale\.zno\'s\'c od $\lambda$ ($+998\,\mathrm{ppm}$ na
    podwojenie $\lambda$).
\end{description}
Fizycznie: \emph{masa elektronu pochodzi z~geometrii (sp\k{e}\.zenie
$\alpha$, rozmiar $a_\Gamma$), a~masa taonu z~\l{}amania symetrii
($\lambda$ -- analogon masy Higgsa).
Masa mionu le\.zy pomi\k{e}dzy -- ma obie cechy.}

\paragraph{Wnioski P61.}
\begin{enumerate}
  \item OP-10 rozwiazany: $K_1$ jest niezale\.zne od~$\lambda$
        z~powodu tlumienia $\lambda K_1^5 I_6 \sim 3\times10^{-13}$
        (13 rz\k{e}d\'ow).
  \item Sensytywno\'s\'c: $K_3\propto\lambda^{-1/2}$,
        $K_2$ s\l{}abo ($+998\,\mathrm{ppm}$ na $\times2$),
        $K_1$ dok\l{}adnie zero.
  \item Wzor wiodacy $K_1^{\rm lead}=1/A_1$ z~b\l{}\k{e}dem $-9{,}33\%$ --
        szereg perturbacyjny wolno zbie\.zny (por. OP-1 dla $K_2$).
  \item Kolejny priorytet: OP-2 (selekcja g\'ornego zera $\alpha_2$) -- P62.
\end{enumerate}

% . . . . . . . . . . . . . . . . . . . . . . . . . . . . . . . .
\subsubsection{Wyniki P62: selekcja g\'ornego zera $\alpha_2$ (OP-2 $\to$ rozwiazany)}%
\label{app:F-P62}
% . . . . . . . . . . . . . . . . . . . . . . . . . . . . . . . .

Skrypt \texttt{p62\_alpha2\_selection.py} (sesja~P62) bada mechanizm
selekcji g\'ornego zera $\alpha_2$ krywej $Q(\alpha)=3/2$
w~obliczu faktu z~P53: $E(\alpha_1) < E(\alpha_2)$ (ni\.zsze zero
ma ni\.zsza energi\k{e}).

\paragraph{Galaz dolna $\alpha_1(a)$ -- wyklucze nie kinematyczne.}
Dla ka\.zdej wartosci parametru $a>a_c$ wyznaczono dolne zero
$\alpha_1(a)$ oraz odpowiadajacy stosunk mas
$r_{21}^{(1)}(a)=K_2(\alpha_1(a))/K_1(\alpha_1(a))$.
Skanujac $a\in[a_c,\,0{,}075]$ wyznaczono globalne maksimum:
\begin{equation}
  r_{21}^{(1)}\big|_{\max}
  = 76{,}57
  \quad\text{przy }a=0{,}03848,\;\alpha_1=3{,}278.
  \label{eq:r21-lower-max}
\end{equation}
Poniewaz $r_{21}^{\mathrm{PDG}} = 206{,}768 > 76{,}57$,
galaz dolna \emph{nigdy} nie spelnia warunku $r_{21}=r_{21}^{\mathrm{PDG}}$
dla \emph{zadnej} wartosci $a$ (przy $\lambda=\lambda_K$).
Deficyt: $r_{21}^{\mathrm{PDG}} - r_{21}^{(1)}\big|_{\max} = +130{,}2$
(brak $170\%$ wymaganej wartosci).

\paragraph{Mapa galeizi -- tabela.}
\begin{center}
\begin{tabular}{ccccccc}
\hline
$a$ & $\alpha_1$ & $r_{21}^{(1)}$ & $\alpha_2$ & $r_{21}^{(2)}$
    & $E(\alpha_1)<E(\alpha_2)$? \\
\hline
0{,}0385 & 3{,}22 & 75{,}2 & 4{,}98 & 116{,}9 & tak \\
0{,}0390 & 2{,}35 & 56{,}0 & 6{,}46 & 153{,}5 & tak \\
0{,}0395 & 1{,}90 & 46{,}6 & 7{,}53 & 180{,}5 & tak \\
\textbf{0{,}0400} & \textbf{1{,}58} & \textbf{40{,}3} & \textbf{8{,}47} & \textbf{204{,}5} & tak \\
0{,}0405 & 1{,}33 & 35{,}5 & 9{,}36 & 227{,}1 & tak \\
$\geq0{,}046$ & nie istnieje & -- & istnieje & ro\'snie & -- \\
\hline
\end{tabular}
\end{center}

\paragraph{Trzy aspekty selekcji.}
\begin{description}
  \item[Kinematyka ($r_{21}$):] Dolna galaz spe\l{}nia maksymalnie
        $r_{21}\leq76{,}6\ll206{,}8=r_{21}^{\mathrm{PDG}}$ --
        wykluczona empirycznie.
  \item[Energia:] $E(\alpha_1) < E(\alpha_2)$ dla wszystkich zbadanych $a$ --
        selekcja energetyczna dzia\l{}a\l{}aby \emph{przeciw} $\alpha_2$.
  \item[Topologia:] Indeks~$+1$ przy $\alpha_2$ (Q ro\'snie przez $3/2$)
        vs indeks~$-1$ przy $\alpha_1$ (Q opada); suma ind. $=0$.
\end{description}

\paragraph{Indeksy topologiczne.}
\begin{center}
\begin{tabular}{lcccc}
\hline
Zero & $\alpha$ & $\mathrm{d}Q/\mathrm{d}\alpha$ & Indeks & $r_{21}$ \\
\hline
$\alpha_1$ (dolne) & $1{,}553$ & $-0{,}0076$ & $-1$ & $39{,}7$ \\
$\alpha_2$ (g\'orne) & $8{,}562$ & $+0{,}0022$ & $+1$ & $206{,}8\approx\mathrm{PDG}$ \\
\hline
\end{tabular}
\end{center}

\paragraph{Twierdzenie (P62) -- selekcja $\alpha_2$.}
\begin{quote}
\emph{Dolne zero $\alpha_1(a)$ jest kinematycznie wykluczone przez
warunek $r_{21}=r_{21}^{\mathrm{PDG}}$ dla wszystkich $a>a_c$
przy $\lambda=\lambda_K$.
G\'orne zero $\alpha_2$ jest jedynym fizycznie dopuszczalnym
rozwi\k{a}zaniem.}
\end{quote}
Pytanie \emph{,,dlaczego $\alpha_2$ a nie $\alpha_1$?''} redukuje si\k{e} do
\emph{,,dlaczego $r_{21}=206{,}77$?''} -- co jest danq empiryczn\k{a} PDG,
nie przewidywaniem TGP na obecnym etapie.

\paragraph{Pytanie g\l{}\k{e}bsze (wci\k{a}\.z otwarte).}
Czy istnieje dynamiczny mechanizm w~TGP wyprodukujacy $r_{21}^{\mathrm{PDG}}$
bez odwo\l{}ania do danych eksperymentalnych?
Jest to pytanie o predyktywnosc modelu na poziomie mas
i~nale\.zy do OP-1 (analityczny wzor na $K_2$, centralne pytanie).

\paragraph{Wnioski P62.}
\begin{enumerate}
  \item OP-2 rozwiazany: $\alpha_2$ konieczne przez wyklucze nie kinematyczne
        $\alpha_1$ (brak $r_{21}^{\mathrm{PDG}}$ na galezi dolnej).
  \item Energia nie wybiera $\alpha_2$: $E(\alpha_1)<E(\alpha_2)$ zawsze.
  \item Wyb\'or $\alpha_2$ jest r\'ownowa\.zny empirycznym mas PDG ($r_{21}=206{,}77$).
  \item Kolejny priorytet: OP-1 ($K_2$ z pierwszych zasad) -- P63.
\end{enumerate}

% . . . . . . . . . . . . . . . . . . . . . . . . . . . . . . . .
\subsubsection{Wyniki P63: $K_2$ z pierwszych zasad -- diagnoza OP-1}%
\label{app:F-P63}
% . . . . . . . . . . . . . . . . . . . . . . . . . . . . . . . .

Skrypt \texttt{p63\_K2\_first\_principles.py} (sesja~P63) bada
mo\.zliwo\'s\'c analitycznego wyznaczenia $K_2$ z~warunku $g(K_2)=0$.

\paragraph{Fundamentalna przeszkoda: promien zbieznosci.}
Funkcja $g(K)=E[K]/(4\pi K)-1$ posiada osobliwo\'s\'c analityczna
przy $K_{\rm sing} = -a\,e^a$ (rzeczywista ujemna):
profil solitonu $\varphi(a)=1+K e^{-a}/a = 0$ przy $K=K_{\rm sing}$.
Promien zbiezno\'sci szeregu Taylora:
\begin{equation}
  R_{\rm zbież} \approx |K_{\rm sing}| = a\,e^a
  = 0{,}040049\times e^{0{,}040049} \approx 0{,}042.
  \label{eq:K2-Rconv}
\end{equation}
Poniewa\.z $K_2\approx2{,}033\gg0{,}042$,
\emph{szereg Taylora $g(K)$ wok\'o\l{} $K=0$ nie zbiega do $K_2$}.
Jest to fundamentalna przeszkoda algebraiczna.

\paragraph{Nowe wsp\'o\l{}czynniki szeregu.}
Dopasowanie wielomianowe do $g(K)$ przy ma\l{}ych $K$
daje (por. P47 dla pierwszych czterech):
\begin{align*}
  c_2 &= +112{,}16, & c_3 &= -1228{,}3, & c_4 &= +19652, \\
  c_5 &= -352\,572  &     & (\text{nowy!}) & c_6 &\approx 10^7 \;(\text{rozbiez.})
\end{align*}
Wsp\'o\l{}czynnik $c_5=-352\,572$ by\l{}a dotychczas pomijany --
odpowiada on cz\l{}onowi kinetycznemu $\sim\alpha K^5$ i~jest wa\.zny
dla wy\.zszych aproksymant\'ow Pad\'e.

\paragraph{Brak relacji Vieta.}
Badano mo\.zliwo\'s\'c wyra\.zenia $K_2$ przez $K_1$ i~$K_3$:
\begin{center}
\begin{tabular}{lrr}
\hline
Propozycja & Warto\'s\'c & B\l{}\k{a}d \\
\hline
$\sqrt{K_1 K_3}$ & $0{,}581$ & $-71{,}4\%$ \\
$(K_1+K_3)/2$    & $17{,}15$ & $+744\%$ \\
\hline
\end{tabular}
\end{center}
Zera $K_1,K_2,K_3$ nie spe\l{}niaja \k{z}adnej prostej relacji algebraicznej.

\paragraph{Balans energii przy $K_2$.}
Przy $K_2=2{,}033$ energia jest wynikiem ogromnego anulowania:
\begin{center}
\begin{tabular}{lrr}
\hline
Sk\l{}adnik & Warto\'s\'c & \% $E_{\rm tot}$ \\
\hline
$E_{\rm kin}$ (bez $\alpha$) & $+610{,}6$ & $+2390\%$ \\
$E_{\rm kin}$ ($\alpha$-czl.) & $+380{,}5$ & $+1489\%$ \\
$E_{V_3}$ (kubiczny)          & $+107{,}9$ & $+422\%$ \\
$E_{V_4}$ (kwartet)           & $-1076{,}4$ & $-4213\%$ \\
$E_\lambda$                   & $+2{,}97$  & $+12\%$ \\
\hline
$E_{\rm tot}$                 & $+25{,}55$ & $100\%$ \\
\hline
\end{tabular}
\end{center}
Nie istnieje jeden dominujacy sk\l{}adnik -- $K_2$ powstaje z~subtelnego
anulowania du\.zych cz\l{}on\'ow kinetycznego i~kwartycznego.

\paragraph{Najlepsza aproksymacja numeryczna.}
Interpolacja wielomianowa przez 12 punkt\'ow $g(K)$ daje:
\begin{equation}
  K_2^{\rm interp} = 2{,}0332176,\quad
  |\delta K_2| = 0{,}9\,\mathrm{ppm}\text{ vs }K_2^{\rm num}.
  \label{eq:K2-interp}
\end{equation}
Jest to praktycznie precyzyjny wynik bez wzoru analitycznego.

\paragraph{Wniosek OP-1.}
\begin{quote}
\emph{Wz\'or zamkni\k{e}ty na $K_2$ nie istnieje w~klasycznym sensie
(szereg Taylora), poniewa\.z $K_2$ le\.zy poza promieniem zbiezno\'sci
($R\approx0{,}042$). Optymalnym podej\'sciem jest aproksymacja Pad\'e $[n/m]$
jako przed\l{}u\.zenie analityczne, lub bezpo\'srednia interpolacja $g(K)$.}
\end{quote}

\paragraph{Wnioski P63.}
\begin{enumerate}
  \item Fundamentalna przeszkoda zidentyfikowana: $R_{\rm zbież}=a e^a\approx0{,}042\ll K_2\approx2$.
  \item Nowy wsp\'o\l{}czynnik $c_5=-352\,572$ (czlon kinetyczny $\sim\alpha K^5$);
        konieczny dla Pad\'e $[3/3]$ i~wy\.zszych.
  \item $K_2$ wyznaczone przez subtelny balans kinetycznego i~kwartycznego
        przy ogromnym anulowaniu ($\sim4000\%$).
  \item Interpolacja $g(K)$ przez 12 punkt\'ow daje $K_2$ z~$0{,}9\,\mathrm{ppm}$.
  \item Dalszy post\k{e}p: poprawne wzory analityczne $c_5,c_6,c_7$
        (b\l{}\k{a}d w~ca\l{}kach kinetycznych P63 -- naprawione w~P64)
        i~budowa Pad\'e~$[3/3]$.
\end{enumerate}

% ---------------------------------------------------------------
\subsubsection{Wyniki P64: Poprawione ca\l{}ki kinetyczne -- weryfikacja $c_n$ i~test Pad\'e (OP-1)}
\label{app:F-P64}
% Skrypt: p64_Pade_K2_corrected.py
% ---------------------------------------------------------------

\paragraph{Naprawa b\l{}\k{e}du ca\l{}ek kinetycznych $M_n$.}
W~sesji P63 mianowniki ca\l{}ek kinetycznych by\l{}y b\l{}\k{e}dne
($r^{n-2}$ zamiast $r^n$). Poprawna definicja wynika z~rozk\l{}adu energii kinetycznej:
\begin{equation}
  E_k = 4\pi K^2 \int_a^\infty
    \frac{e^{-2r}(r+1)^2}{2r^2}\!\left(1 + \frac{\alpha}{\varphi}\right)dr,
  \qquad
  M_n \;=\; \frac{1}{2}\int_a^\infty \frac{e^{-nr}(r+1)^2}{r^n}\,dr.
  \label{eq:Mn-correct}
\end{equation}
Rozw\l{}adaj\k{a}c $1/\varphi = \sum_{k=0}^\infty (-1)^k u^k$, $u=Ke^{-r}/r$,
otrzymujemy k-ty wyraz kinetyczny $\sim K^{k+2} M_{k+2}$, co potwierdza mianownik $r^n$.

\paragraph{Poprawione wsp\'o\l{}czynniki $c_n$.}
\begin{center}
\begin{tabular}{r|rr|r}
\hline
$n$ & $c_n^{\rm an}$ (P64) & $c_n^{\rm num}$ & $\delta$ \\
\hline
2 & $+112{,}156$ & $+112{,}156$ & $-0{,}000024\%$ \\
3 & $-1228{,}295$ & $-1228{,}290$ & $+0{,}0004\%$ \\
4 & $+19\,657$ & $+19\,653$ & $+0{,}018\%$ \\
5 & $-3{,}539\times10^5$ & $-3{,}528\times10^5$ & $+0{,}32\%$ \\
6 & $+6{,}793\times10^6$ & $+6{,}599\times10^6$ & $+2{,}94\%$ \\
7 & $-1{,}358\times10^8$ & $-1{,}162\times10^8$ & $+16{,}9\%$ \\
\hline
\end{tabular}
\end{center}
Wsp\'o\l{}czynniki $c_2,c_3,c_4$ zgodne analitycznie i~numerycznie do $<0{,}02\%$.
Rosn\k{a}ca rozbie\.zno\'s\'c dla $c_{5+}$ wynika z~z\l{}ego uwarunkowania
wielomianowego dopasowania numerycznego przy wy\.zszych stopniach.

\paragraph{Szereg Taylora jest asymptotyczny przy $K_2$.}
Termy szeregu przy $K_2=2{,}033$:
\[
  g(K_2) = \underbrace{-1}_{-1}
  +\underbrace{228}_{c_2 K_2}
  \underbrace{-5078}_{c_3 K_2^2}
  +\underbrace{165\,220}_{c_4 K_2^3}
  \underbrace{-6\times10^6}_{c_5 K_2^4}
  +\underbrace{2{,}4\times10^8}_{c_6 K_2^5}
  \underbrace{-9{,}6\times10^9}_{c_7 K_2^6}
  +\underbrace{4{,}0\times10^{11}}_{c_8 K_2^7}
  + \cdots
\]
Kolejne cz\l{}ony rosn\k{a} wyk\l{}adniczo -- szereg jest bezwzgl\k{e}dnie
rozbie\.zny w~$K_2$ (szereg asymptotyczny). Aproksymacja Pad\'e $[n/m]$
z~wsp\'o\l{}czynnik\'ow Taylora:
\textbf{wszystkie konfiguracje~$[n/m]$ zwracaj\k{a} NaN lub b\l{}\k{a}d $\sim50\%$}.
Jest to wynik strukturalny (nie b\l{}\k{a}d implementacji).

\paragraph{Wniosek OP-1 (uaktualniony po P64).}
\begin{quote}
\emph{Wzory analityczne $c_n$ s\k{a} poprawne do rz\k{e}du $c_4$ (delta~$<0{,}02\%$).
Pad\'e z~szeregu Taylora jest niemo\.zliwe strukturalnie, poniewa\.z $K_2\gg R_{\rm zbież}$.
Jedyna dzia\l{}aj\k{a}ca metoda analityczna to interpolacja wymierna
z~warto\'sci $g(K)$ (b\l{}\k{a}d $-0{,}9\,\mathrm{ppm}$, P63).}
\end{quote}

\paragraph{Wnioski P64.}
\begin{enumerate}
  \item B\l{}\k{a}d $M_n$ (P63) naprawiony: $c_2=112{,}156$ zgodne z~P47 ($\delta=-2{,}4\times10^{-5}\%$).
  \item Wzory analityczne $c_2,c_3,c_4$ potwierdzone numerycznie.
  \item Taylor-Pad\'e strukturalnie niemo\.zliwe; serie asymptotyczne nie sumuj\k{a} si\k{e} do $g(K_2)$.
  \item Nast\k{e}pny krok (P65): systematyczna interpolacja wymierna $[n/m]$
        z~$\geq20$ punkt\'ow $g(K)\in[0{,}5,4{,}0]$ -- wyznaczenie optymalnego~$[n/m]$.
\end{enumerate}

% ---------------------------------------------------------------
\subsubsection{Wyniki P65: Interpolacja wymierna~$R[n/m](K)$ -- OP-1 rozwi\k{a}zany numerycznie}
\label{app:F-P65}
% Skrypt: p65_rational_interp_K2.py
% ---------------------------------------------------------------

\paragraph{Metoda.}
Aproksymant wymierny $R[n/m](K) = P_n(K)/Q_m(K)$, $Q_m(0)=1$,
dopasowany metod\k{a} MNK do 21 warto\'sci $g(K)$ na $[0{,}06,\,20{,}0]$.
Dla ka\.zdego punktu $(K_i, g_i)$ obowi\k{a}zuje r\'ownanie liniowe:
\begin{equation}
  P_n(K_i) - g_i\,Q_m(K_i) = 0,
  \label{eq:rat-interp-linear}
\end{equation}
sk\k{a}d wsp\'o\l{}czynniki wyznacza si\k{e} MNK.
Zera $g(K)$ odpowiadaj\k{a} zerom licznika $P_n(K)$.

\paragraph{Tabela dok\l{}adno\'sci~$K_2$ dla r\'o\.znych $[n/m]$.}
\begin{center}
\begin{tabular}{c|rr|cc|r}
\hline
$[n/m]$ & $K_2^{R}$ & $\delta K_2$ (ppm) & $K_1$? & $K_3$? & par. \\
\hline
$[2/3]$ & $2{,}033401$ & $+89{,}1$ & nie & nie & 6 \\
$[6/1]$ & $2{,}033036$ & $-90{,}2$ & nie & tak & 8 \\
$[6/2]$ & $2{,}033197$ & $-10{,}9$ & nie & tak & 9 \\
$[7/2]$ & $\mathbf{2{,}033223}$ & $\mathbf{+1{,}79}$ & tak$^*$ & tak & 10 \\
\hline
\end{tabular}
\end{center}
\footnotesize{$^*$~$K_1$ z du\.zym b\l{}\k{e}dem ($-36\%$), poniewa\.z le\.zy poza zakresem interpolacji.}\normalsize

\paragraph{Optymalny aproksymant $R[7/2]$.}
\begin{equation}
  R[7/2](K) =
  \frac{-0{,}529 + 82{,}68\,K + 228{,}9\,K^2 + 81{,}87\,K^3
        - 43{,}22\,K^4 - 30{,}78\,K^5 + 0{,}028\,K^7}
       {1 + 8{,}70\,K + 7{,}88\,K^2}
  \label{eq:R72}
\end{equation}
Trzy zera licznika: $K_1^R=0{,}00629$, $K_2^R=2{,}03322$, $K_3^R=34{,}289$.
\begin{center}
\begin{tabular}{crrr}
\hline
Zero & $R[7/2]$ & Numeryczne & $\delta$ \\
\hline
$K_1$ & $0{,}006285$ & $0{,}009833$ & $-36\%$ (poza zakresem) \\
$K_2$ & $\mathbf{2{,}033223}$ & $\mathbf{2{,}033219}$ & $\mathbf{+1{,}79\text{ ppm}}$ \\
$K_3$ & $34{,}2889$ & $34{,}2908$ & $-56\text{ ppm}$ \\
\hline
\end{tabular}
\end{center}

\paragraph{Stabilno\'s\'c.}
Szum numeryczny $\varepsilon=10^{-4}$ w~$g(K_i)$ powoduje
$\sigma(K_2) \approx 1{,}3\times10^{-6}$ -- metoda jest numerycznie solidna.
Czu\l{}o\'s\'c $K_2$ na parametry modelu jest zgodna z~warto\'sci\k{a} numeryczn\k{a}
($\delta K_2/K_2 \approx 211\,\mathrm{ppm}/0{,}1\%a$,\;
$184\,\mathrm{ppm}/0{,}1\%\alpha$).

\paragraph{Wniosek OP-1 (P65).}
\begin{quote}
\emph{Interpolacja wymierna $R[7/2](K)$ z~21 warto\'sci $g(K)$ identyfikuje
$K_2$ z~b\l{}\k{e}dem $+1{,}79\,\mathrm{ppm}$.
OP-1 jest rozwi\k{a}zany w~sensie numerycznym.
Otwarte pytanie: czy wsp\'o\l{}czynniki $[p_j, q_j]$
maj\k{a} interpretacj\k{e} fizyczn\k{a}?}
\end{quote}

\paragraph{Wnioski P65.}
\begin{enumerate}
  \item Najdok\l{}adniejszy $[n/m]$: $R[7/2]$, $\delta K_2=+1{,}79\,\mathrm{ppm}$ (10 par.).
  \item Minimalny schemat: $R[2/3]$, $\delta K_2=+89\,\mathrm{ppm}$ (6 par.).
  \item Metoda MNK z~21 punktami jest stabilna i~dok\l{}adna.
  \item OP-1 rozwi\k{a}zany numerycznie; interpretacja fizyczna wsp\'o\l{}czynnik\'ow otwarta.
  \item Nast\k{e}pny krok (P66): zale\.zno\'s\'c $[p_j,q_j]$ od $(a,\alpha,\lambda)$ --
        poszukiwanie wzoru zamkni\k{e}tego.
\end{enumerate}

% ---------------------------------------------------------------
\subsubsection{Wyniki P66: Liniowo\'s\'c~$g(K;\alpha)$ w~$\alpha$ -- wz\'or zamkni\k{e}ty na~$K_2$}
\label{app:F-P66}
% Skrypt: p66_Rpq_vs_params.py
% ---------------------------------------------------------------

\paragraph{Liniowo\'s\'c strukturalna.}
Energia kinetyczna
$E_k = 4\pi K^2\!\int\! \frac{1}{2}\dot\varphi^2(1+\alpha/\varphi)\,dr$
jest liniowa w~$\alpha$, wi\k{e}c:
\begin{equation}
  g(K;\alpha) = g_0(K) + \alpha\,g_1(K),
  \label{eq:g-linear-alpha}
\end{equation}
gdzie $g_0$ -- sk\l{}adnik bez $\alpha$, $g_1$ -- korekta $\alpha$-kinetyczna.
Konsekwencja: \emph{wszystkie wsp\'o\l{}czynniki aproksymantu wymiernego
$R[n/m]$ s\k{a} liniowe w~$\alpha$} ($R^2=1{,}000$ dla ka\.zdego $p_j, q_j$).

\paragraph{Wsp\'o\l{}czynniki $R[2/3]$ jako funkcje~$\alpha$.}
Przy~sta\l{}ym $a=a_\Gamma$:
\begin{equation}
  p_j(\alpha) = A_j + B_j\,\alpha, \qquad j=0,1,2,
\end{equation}
z~warto\'sciami (Punkt~B):
\begin{center}
\begin{tabular}{c|rr}
\hline
Wsp. & $A_j$ & $B_j$ \\
\hline
$p_0$ & $-3{,}961$ & $+0{,}6500$ \\
$p_1$ & $+24{,}62$ & $+0{,}7779$ \\
$p_2$ & $-14{,}08$ & $-0{,}1977$ \\
\hline
\end{tabular}
\end{center}

\paragraph{Wz\'or zamkni\k{e}ty na~$K_2$ z~$R[2/3]$.}
Szukamy zera $P_2(K)=p_0+p_1 K+p_2 K^2=0$:
\begin{equation}
  \boxed{K_2 = \frac{-p_1 - \sqrt{p_1^2 - 4p_2\,p_0}}{2p_2},
  \qquad p_j = A_j + B_j\,\alpha,}
  \label{eq:K2-closed-form}
\end{equation}
daje $\delta K_2 = +89\,\mathrm{ppm}$ -- wynik analityczny!
Sze\'s\'c sta\l{}ych empirycznych $\{A_j, B_j\}$ (zale\.znych od~$a$)
jest jedynym ,,wej\'sciem'' wzoru.

\paragraph{Prawo skalowania.}
Dopasowanie dwuwymiarowe do~56 punkt\'ow~$(a,\alpha)$:
\begin{equation}
  K_2(a,\alpha) \approx 2{,}692\cdot a^{0{,}205}\cdot\alpha^{0{,}177},
  \qquad \text{RMS} = 0{,}25\%.
  \label{eq:K2-scaling}
\end{equation}

\paragraph{Stabilno\'s\'c $R[7/2]$ na~siatce.}
B\l{}\k{a}d $R[7/2]$ na~56 punktach $(a,\alpha)$:
$\delta K_2 \in [+1{,}0,\,+2{,}1]\,\mathrm{ppm}$ -- absolutnie jednolity.
Metoda jest przenaszalna.

\paragraph{Wniosek P66.}
\begin{quote}
\emph{$g(K;\alpha)$ jest liniowe w~$\alpha$ z~zasad pierwszych (r.~\ref{eq:g-linear-alpha}).
Wszystkie wsp\'o\l{}czynniki $R[n/m]$ s\k{a} liniowe w~$\alpha$~($R^2=1$).
Wz\'or~(\ref{eq:K2-closed-form}) daje $K_2$ z~b\l{}\k{e}dem $+89\,\mathrm{ppm}$
jako funktjon\'a $\alpha$ i~6~sta\l{}ych empirycznych~$\{A_j, B_j\}$.}
\end{quote}

\paragraph{Wnioski P66.}
\begin{enumerate}
  \item $g(K;\alpha)$ liniowe w~$\alpha$ -- wynik strukturalny z~zasad pierwszych.
  \item Wsp\'o\l{}czynniki $R[2/3]$: $p_j=A_j+B_j\alpha$, $R^2=1{,}000$.
  \item Wz\'or zamkni\k{e}ty~(\ref{eq:K2-closed-form}): $\delta K_2=+89\,\mathrm{ppm}$.
  \item $R[7/2]$: $\delta K_2\in[1{,}0;\,2{,}1]\,\mathrm{ppm}$ jednolicie dla~56~punkt\'ow $(a,\alpha)$.
  \item Prawo skalowania $K_2\approx2{,}692\,a^{0{,}205}\alpha^{0{,}177}$ (RMS~$0{,}25\%$).
  \item Nast\k{e}pny krok (P67): wyznaczy\'c $A_j(a), B_j(a)$ -- pe\l{}ny wz\'or $K_2(a,\alpha)$.
\end{enumerate}

% ---------------------------------------------------------------
\subsubsection{Wyniki P67: R\'ownanie implicit $K_2$ -- odkrycie strukturalne}
\label{app:F-P67}
% Skrypt: p67_K2_closed_form.py
% ---------------------------------------------------------------

\paragraph{Rozkl{}ad $g = g_0 + \alpha g_1$.}
Z~liniowo\'sci (r.~\ref{eq:g-linear-alpha}) warunek $g(K_2;\alpha)=0$ daje:
\begin{equation}
  \boxed{g_0(K_2) \;=\; -\alpha\,g_1(K_2)
  \qquad\Longleftrightarrow\qquad
  \alpha \;=\; -\frac{g_0(K_2)}{g_1(K_2)},}
  \label{eq:K2-implicit}
\end{equation}
gdzie $g_0(K)=g(K;\alpha=0)$, $g_1(K)=g(K;1)-g(K;0)$.

\paragraph{Zera i~przes{}uni\k{e}cia.}
\begin{center}
\begin{tabular}{l|rr|r}
\hline
Zero & $K$ dla $g_0$ & $K$ dla $g(\cdot;\alpha_K)$ & $\delta$ \\
\hline
$K_1$ & $0{,}0877$ & $0{,}00983$ & $-89\%$ \\
$K_2$ & $\mathbf{1{,}533}$ & $\mathbf{2{,}033}$ & $\mathbf{+32{,}6\%}$ \\
$K_3$ & $34{,}293$ & $34{,}291$ & $-0{,}007\%$ \\
\hline
\end{tabular}
\end{center}
Cz\l{}on $\alpha g_1$ \emph{przesuwa} $K_2$ o~$+32{,}6\%$.
$K_3$ prawie nie reaguje na~$\alpha$ ($\delta=-0{,}007\%$) -- dominuje potencja\l{}.
$g_1(K)>0$ dla wszystkich~$K$ (brak zer, ro\'sn\k{a}ce od~$1{,}2$ do~$2{,}2$).

\paragraph{Wzor zamkni\k{e}ty z~$A_j(a),B_j(a)$.}
Wspolczynniki $A_j(a), B_j(a)$ wyznaczone na~12 punktach $a\in[0{,}025,\,0{,}060]$;
najlepszy fit -- potegowy ($A_j \sim C_j a^{\beta_j}$, RMS~$4$--$7\%$).
B\l{}\k{a}d wzoru zamkni\k{e}tego:
\begin{itemize}
  \item W~punkcie~B: $\delta K_2=+136\,\mathrm{ppm}$.
  \item RMS na~siatce 56 punkt\'ow $(a,\alpha)$: $1168\,\mathrm{ppm}$.
\end{itemize}

\paragraph{Wniosek P67.}
\begin{quote}
\emph{$K_2$ spe\l{}nia r\'ownanie implicit~(\ref{eq:K2-implicit}).
$K_2$ bez~$\alpha$ wynosi $1{,}533$; cz\l{}on $\alpha g_1$ przesuwa je do~$2{,}033$ ($+33\%$).
Nast\k{e}pny krok (P68): iteracja Newtona na~$g_0+\alpha g_1=0$
startuj\k{a}c z~$K_0^{(2)}=1{,}533$.}
\end{quote}

\paragraph{Wnioski P67.}
\begin{enumerate}
  \item R\'ownanie implicit $g_0(K_2)=-\alpha g_1(K_2)$ -- wynik strukturalny.
  \item $K_2^{(0)} = 1{,}533$ (zero $g_0$); $\alpha$ przesuwa je o~$+33\%$.
  \item $K_3$ niezale\.zne od~$\alpha$ ($\delta=-0{,}007\%$) -- dominuje $V_\lambda$.
  \item Wzor zamkni\k{e}ty z~$A_j(a),B_j(a)$: $136\,\mathrm{ppm}$ @~B, $1168\,\mathrm{ppm}$ globalnie.
  \item P68: Newton na~(\ref{eq:K2-implicit}) -- analityczna korekta od~$1{,}533$.
\end{enumerate}

% ---------------------------------------------------------------
\subsubsection{Wyniki P68: Iteracja Newtona na $g_0+\alpha g_1=0$}
\label{app:F-P68}
% Skrypt: p68_newton_g0g1.py
% ---------------------------------------------------------------

\paragraph{Metoda.}
Definiujemy $F(K)=g_0(K)+\alpha\,g_1(K)$ i~stosujemy iteracj\k{e} Newtona:
\begin{equation}
  K_{n+1} = K_n - \frac{F(K_n)}{F'(K_n)},
  \label{eq:newton-F}
\end{equation}
z~punktem startowym $K_0=1{,}53344433$ (zero $g_0$, wyznaczone z~R[4/3] z~b\l{}\k{e}dem $+8{,}2\,\mathrm{ppm}$).

\paragraph{Zbie\.zno\'s\'c lokalna w~punkcie~B.}
Punkt B: $a=0{,}040049$, $\alpha=8{,}5612$, $\lambda=5{,}4677\times10^{-6}$;
$K_2^{\rm num}=2{,}033219$.

\begin{center}
\begin{tabular}{c|l|r}
\hline
krok $n$ & $K_n$ & b\l{}\k{a}d [ppm] \\
\hline
0 & $1{,}53344$ & $-245\,805$ \\
1 & $2{,}12892$ & $+47\,373$ \\
2 & $2{,}04721$ & $+6\,883$ \\
3 & $2{,}03454$ & $+651$ \\
4 & $\mathbf{2{,}03322}$ & $\mathbf{+0{,}54}$ \\
5 & $2{,}033219$ & $\approx 0$ \\
\hline
\end{tabular}
\end{center}
Po 5 krokach -- dok\l{}adno\'s\'c maszynowa.

\paragraph{Zbie\.zno\'s\'c globalna.}
Na siatce 56 punkt\'ow $(a,\alpha)$ po~2 krokach Newtona b\l{\k{e}}dy wynosz\k{a}
$6000$--$59\,000\,\mathrm{ppm}$ -- zbyt wolna zbie\.zno\'s\'c
z~dalekiego startu $K_0\approx 1{,}533$
(odleg\l{}o\'s\'c $|K_0-K_2|\approx 0{,}5$).

\paragraph{Wz\'or analityczny pierwszego kroku.}
Gdy $g_0(K_0)=0$ dok\l{}adnie:
\begin{equation}
  \boxed{K_1 = K_0 - \frac{\alpha\,g_1(K_0)}{g_0'(K_0)+\alpha\,g_1'(K_0)}.}
  \label{eq:newton-step1}
\end{equation}
Wz\'or~(\ref{eq:newton-step1}) jest w~pe\l{}ni analityczny
po~wyznaczeniu $K_0$ (zera $g_0$).

\paragraph{Wsp\'o\l{}czynniki szeregu $g_0$.}
Przy $a=a_\Gamma$ funkcja $g_0(K)=g(K;\alpha=0)$:
\begin{center}
\begin{tabular}{c|r|l}
\hline
$n$ & $c_{n0}$ & Definicja \\
\hline
2 & $+11{,}524$ & $M_2 - I_2/2$ (bez $\alpha$) \\
3 & $-1{,}106$ & $-2I_3/3$ \\
4 & $-3{,}910$ & $-I_4/4$ \\
\hline
\end{tabular}
\end{center}
Uwaga: $c_{20}=11{,}524 \ll c_2=112{,}156$ --
cz\l{}on $\alpha g_1$ dominuje w~pe\l{}nym $g(K;\alpha_K)$.

\paragraph{Aproksymant $R[4/3]$ dla $g_0$.}
Dopasowanie $R[4/3]$ do~$g_0(K)$ daje zero:
\[
  K_0^{[4/3]} = 1{,}53358 \quad
  \bigl(\delta = +8{,}2\,\mathrm{ppm}\ \text{vs }K_0^{\rm num}=1{,}53344\bigr).
\]
Ze startu $K_0^{[4/3]}$ po~4 krokach Newtona b\l{\k{e}}d $K_2 < 1\,\mathrm{ppm}$
dla~wszystkich 56 punkt\'ow.

\paragraph{Wniosek P68.}
\begin{quote}
\emph{Iteracja Newtona na~$F(K)=g_0(K)+\alpha g_1(K)=0$ z~$K_0=1{,}533$
(zero~$g_0$) zbiega do~$K_2$ z~precyzj\k{a} maszynow\k{a} po~5 krokach (lokalnie).
Pierwszy krok ma wz\'or analityczny~(\ref{eq:newton-step1}).
Zbie\.zno\'s\'c globalna wymaga lepszego~$K_0$;
$R[2/3]$ ($+89\,\mathrm{ppm}$) lub~$R[4/3]$ ($+8{,}2\,\mathrm{ppm}$)
dla~$g_0$ daje $<1\,\mathrm{ppm}$ po~4 krokach.}
\end{quote}

\paragraph{Wnioski P68.}
\begin{enumerate}
  \item Newton na~$g_0+\alpha g_1=0$: 5 krok\'ow $\to 0\,\mathrm{ppm}$ lokalnie.
  \item Wz\'or analityczny: $K_1=K_0-\alpha g_1/(g_0'+\alpha g_1')$.
  \item $g_0$ Taylor: $c_{20}=11{,}524$, $c_{30}=-1{,}106$, $c_{40}=-3{,}910$.
  \item $R[4/3]$ dla~$g_0$: start $K_0$ z~$+8{,}2\,\mathrm{ppm}$ $\to$ 4 kroki Newtona $\to<1\,\mathrm{ppm}$.
  \item P69: u\.zycie $R[2/3]$ jako~$K_0$ ($+89\,\mathrm{ppm}$) $\to$ 1 krok Newtona $\to <1\,\mathrm{ppm}$ wszedzie?
\end{enumerate}

% ---------------------------------------------------------------
\subsubsection{Wyniki P69: OP-1 zamkni\k{e}ty -- $R[2/3]+\mathrm{Newton}$}
\label{app:F-P69}
% Skrypt: p69_newton_R23_start.py
% ---------------------------------------------------------------

\paragraph{Por\'ownanie punkt\'ow startowych (punkt~B).}
\begin{center}
\begin{tabular}{l|l|r}
\hline
Punkt startowy & $K_0$ & b\l{}\k{a}d [ppm] \\
\hline
zero $g_0$ & $1{,}53344$ & $-245\,805$ \\
$R[2/3]$ & $2{,}03340$ & $+89{,}1$ \\
$R[7/2]$ & $2{,}03322$ & $+1{,}79$ \\
$K_2^{\rm num}$ & $2{,}033219$ & $0$ \\
\hline
\end{tabular}
\end{center}

\paragraph{Zbie\.zno\'s\'c $R[2/3]\to$ Newton w~punkcie~B.}
\begin{center}
\begin{tabular}{c|l|r}
\hline
krok $n$ & $K_n$ & b\l{}\k{a}d [ppm] \\
\hline
0 ($R[2/3]$) & $2{,}03340072$ & $+89{,}15$ \\
1 ($+$Newton) & $2{,}03321949$ & $\mathbf{+0{,}0102}$ \\
2 ($+$Newton) & $2{,}03321947$ & $\approx 0$ \\
\hline
\end{tabular}
\end{center}

\paragraph{Globalno\'s\'c: siatka 25 punkt\'ow $(a,\alpha)$.}
\begin{center}
\begin{tabular}{l|r|r}
\hline
Metoda & RMS [ppm] & max [ppm] \\
\hline
$R[2/3]$ bezpo\'srednio & $1005$ & $3341$ \\
$R[2/3]+1$ Newton & $3{,}35$ & $14{,}6$ \\
$R[2/3]+2$ Newton & $5{,}7\times10^{-5}$ & $2{,}8\times10^{-4}$ \\
$R[7/2]$ bezpo\'srednio & $\sim 2$ & $\sim 2{,}1$ \\
$R[7/2]+1$ Newton & $6\times10^{-6}$ & $6\times10^{-6}$ \\
\hline
\end{tabular}
\end{center}

\paragraph{Wz\'or zamkni\k{e}ty ($A_j$, $B_j$ z~P66) + Newton.}
U\.zywaj\k{a}c~$K_0 = (-p_1-\sqrt{\Delta})/(2p_2)$ z~$p_j=A_j+B_j\alpha$:
\[
  K_0^{\rm wz\'or} = 2{,}03318802, \quad \delta = -15{,}47\,\mathrm{ppm}
\]
Po jednym kroku Newtona:
\begin{equation}
  K_1 = K_0^{\rm wz\'or} - \frac{g(K_0^{\rm wz\'or};\,\alpha)}{g'(K_0^{\rm wz\'or};\,\alpha)},
  \qquad \delta K_1 = +0{,}00031\,\mathrm{ppm}.
  \label{eq:K2-OP1-final}
\end{equation}
Wzor~(\ref{eq:K2-OP1-final}) jest \textbf{w pe\l{}ni analityczny} dla~$a=a_\Gamma$.

\paragraph{Pochodna $F'(K)$.}
\[
  \frac{F'(K_0)}{F'(K_2)} = 1{,}000229 \approx 1 \quad (<0{,}023\%).
\]
Pochodna prawie sta\l{}a mi\k{e}dzy~$K_0$ a~$K_2$ -- wyja\'snia szybk\k{a} zbie\.zno\'s\'c.

\paragraph{Wniosek P69 -- OP-1 ZAMKNI\k{E}TY.}
\begin{quote}
\emph{Minimalna analityczna \'scie\.zka do $K_2$ z~sub-ppm dok\l{}adno\'sci\k{a}:}
\begin{equation}
  \boxed{K_2 \approx K_0^{[2/3]} - \frac{g(K_0^{[2/3]};\,\alpha)}{g'(K_0^{[2/3]};\,\alpha)},
  \qquad K_0^{[2/3]} = \frac{-p_1-\sqrt{p_1^2-4p_2 p_0}}{2p_2},
  \quad p_j = A_j+B_j\alpha.}
  \label{eq:K2-closed-OP1}
\end{equation}
\emph{B\l{}\k{a}d globalny po~1~kroku: max~$14{,}6\,\mathrm{ppm}$;
po~2 krokach: max~$2{,}8\times10^{-4}\,\mathrm{ppm}$.
OP-1 jest rozwi\k{a}zany analitycznie (schemat: 6~param + 2~ewaluacje~$g$).}
\end{quote}

\paragraph{Wnioski P69.}
\begin{enumerate}
  \item $R[2/3]+1$ Newton $\to +0{,}0102\,\mathrm{ppm}$ (punkt~B);
        globalnie: RMS $=3{,}35\,\mathrm{ppm}$, max $=14{,}6\,\mathrm{ppm}$.
  \item $R[2/3]+2$ Newton $\to$ max $2{,}8\times10^{-4}\,\mathrm{ppm}$ (siatka 25 pkt).
  \item Wz\'or kwadratowy ($A_j,B_j$ z~P66) $+1$ Newton $\to +0{,}00031\,\mathrm{ppm}$ @~B.
  \item $R[7/2]+1$ Newton $\to$ max $6\times10^{-6}\,\mathrm{ppm}$ globalnie.
  \item \textbf{OP-1 ZAMKNI\k{E}TY} -- formu\l{}a~(\ref{eq:K2-closed-OP1}).
  \item P70: OP-4 -- Q=3/2 dok\l{}adnie: u\.zy\'c $K_2$ z~(\ref{eq:K2-closed-OP1}) zamiast $K_2^{\rm Pad\'e}$ i~wyliczy\'c residuum~$Q$.
\end{enumerate}

% ---------------------------------------------------------------
\subsubsection{Wyniki P70: OP-4 zdiagnozowany -- residuum $Q-3/2$}
\label{app:F-P70}
% Skrypt: p70_Q_residuum_K2_precise.py
% ---------------------------------------------------------------

\paragraph{Q przy r\'o\.znych przybli\.zeniach $K_2$.}
Formu\l{}a Koidego: $Q = (\sqrt{K_1}+\sqrt{K_2}+\sqrt{K_3})^2/(K_1+K_2+K_3) = 3/2$.
\begin{center}
\begin{tabular}{l|r|r}
\hline
Metoda $K_2$ & b\l{}\k{a}d $K_2$ [ppm] & $Q - 3/2$ [ppm] \\
\hline
$K_2^{\rm num}$ (bazowa) & $0$ & $-632{,}46$ \\
$K_2^{R[2/3]+1N}$ & $+0{,}010$ & $-632{,}46$ \\
$K_2^{R[2/3]+2N}$ & $\approx 0$ & $-632{,}46$ \\
\hline
\end{tabular}
\end{center}
Zmiana $K_2$ o $+0{,}010\,\mathrm{ppm}$ zmienia $Q$ o $+0{,}002\,\mathrm{ppm}$ -- \textbf{ca\l{}kowicie pomijalne}.

\paragraph{Wra\.zliwo\'s\'c $Q$ na $K_i$.}
\[
  \frac{\partial Q}{\partial K_1} = +2{,}007, \quad
  \frac{\partial Q}{\partial K_2} = +0{,}1012, \quad
  \frac{\partial Q}{\partial K_3} = -0{,}00658.
\]
Dominuje wk\l{}ad przez $K_3$ (mimo ma\l{}ej pochodnej -- $K_3$ jest bardzo du\.ze).

\paragraph{$\lambda^*$ dla $Q=3/2$ dok\l{}adnie.}
Przy sta\l{}ych $(a_\Gamma,\alpha_K)$ i~zmiennym $\lambda$:
\begin{center}
\begin{tabular}{l|r|r|r}
\hline
Parametr & $\lambda_K$ & $\lambda^*$ (Q=3/2) & r\'o\.znica \\
\hline
$\lambda$ & $5{,}4677\times10^{-6}$ & $5{,}4984\times10^{-6}$ & $\mathbf{+5616\,\mathrm{ppm}}$ \\
$K_1$ & $0{,}009833$ & $0{,}009833$ & $0$ \\
$K_2$ & $2{,}033219$ & $2{,}033236$ & $+8\,\mathrm{ppm}$ \\
$K_3$ & $\mathbf{34{,}2908}$ & $\mathbf{34{,}1951}$ & $-2792\,\mathrm{ppm}$ \\
$r_{31}$ & $3487{,}2$ & $3477{,}5$ & $-2792\,\mathrm{ppm}$ \\
\hline
\end{tabular}
\end{center}
\textbf{Kluczowe odkrycie}: $K_3^* = 34{,}195 = K_3^{\rm Ei}$ (warto\'s\'c z~przybli\.zenia Yukawa, P60)!
$r_{31}^* = 3477{,}5 \approx r_{31}^{\rm PDG}$.

\paragraph{Mechanizm residuum $Q$.}
Residuum $Q = 3/2 - 632\,\mathrm{ppm}$ pochodzi z~b\l{}\k{e}du $K_3$:
\[
  K_3^{\rm num} = 34{,}291 \quad \text{zamiast} \quad K_3^* = 34{,}195 \quad
  (\delta K_3 = -2792\,\mathrm{ppm} = -0{,}28\%).
\]
Korekcja: $\lambda^* = \lambda_K \times 1{,}005616$ obni\.za $K_3$ do $K_3^{\rm Ei}$ i~daje $Q = 3/2$ dok\l{}adnie.

\paragraph{Wniosek P70 -- OP-4 zdiagnozowany.}
\begin{quote}
\emph{Residuum $Q-3/2 = -632\,\mathrm{ppm}$ przy punkcie~B jest ca\l{}kowicie niezale\.zne
od dok\l{}adno\'sci $K_2$ (wp\l{}yw $K_2$ na~$Q$ wynosi $0{,}002\,\mathrm{ppm}$ na~$0{,}010\,\mathrm{ppm}$
b\l{}\k{e}du~$K_2$). Dominuje b\l{}\k{a}d $K_3$: $K_3^{\rm num}=34{,}291$ jest za~du\.ze o~$0{,}28\%$.
Aby $Q=3/2$ dok\l{}adnie przy $(a_\Gamma,\alpha_K)$: $\lambda^* = \lambda_K\times1{,}005616$,
co redukuje $K_3$ do~$K_3^{\rm Ei}=34{,}195$ i~daje $r_{31}^*=3477{,}5\approx r_{31}^{\rm PDG}$.
OP-4 jest problemem $\lambda$ i~$K_3$, nie~$K_2$.}
\end{quote}

\paragraph{Wnioski P70.}
\begin{enumerate}
  \item $Q@B = 3/2 - 632\,\mathrm{ppm}$ (niezale\.znie od dok\l{}adno\'sci $K_2$).
  \item $\partial Q/\partial K_2 = 0{,}1012$; b\l{}\k{a}d $K_2 = +0{,}010\,\mathrm{ppm}$ $\to$ $\Delta Q = +0{,}002\,\mathrm{ppm}$.
  \item \'Zr\'od\l{}o residuum: $K_3^{\rm num}=34{,}291 > K_3^*=34{,}195 = K_3^{\rm Ei}$.
  \item $\lambda^*/\lambda_K = +5616\,\mathrm{ppm}$ (+0{,}56\%); $r_{31}^*=3477{,}5 \approx r_{31}^{\rm PDG}$.
  \item P71: zbada\'c sk\k{a}d $K_3^{\rm num}-K_3^*=2792\,\mathrm{ppm}$ i~zwi\k{a}zek z~predykcj\k{a}~$\lambda_{\rm Koide}$.
\end{enumerate}

% ---------------------------------------------------------------
\subsubsection{Wyniki P71: Mechanizm $K_3$--$\lambda$ -- anatomia b\l{}\k{e}du 2802\,ppm}
\label{app:F-P71}
% Skrypt: p71_K3_lambda_analysis.py (naprawiony: b\l{}\k{e}dna I_6 \to dok\l{}adny wz\'or analityczny)
% ---------------------------------------------------------------

\paragraph{$K_3 \sim \lambda^{-1/2}$ -- weryfikacja dok\l{}adna.}
Fit pot\k{e}gowy $K_3 \propto \lambda^\beta$ na siatce $\lambda/\lambda_K \in [0{,}90;\,1{,}20]$:
\[
  \beta = -0{,}49930834, \qquad |\beta - (-\tfrac{1}{2})| = 691{,}66\,\mathrm{ppm}.
\]
Kluczowa weryfikacja:
\[
  \left(\frac{K_3^{\rm num}}{K_3^*}\right)^2 = 1{,}00560681
  \approx \frac{\lambda^*}{\lambda_K} = 1{,}00561479
  \qquad (\text{r\'o\.znica } 7{,}98\,\mathrm{ppm}).
\]

\paragraph{Ca\l{}ki $I_4$, $I_6$ -- wzory analityczne vs. numeryczne.}
Przy $a = 0{,}040049$:
\[
  I_4 = \int_a^\infty \frac{e^{-4r}}{r^2}\,dr
      = \frac{e^{-4a}}{a} - 4E_1(4a)
      = 15{,}640786,
\]
\[
  I_6 = \int_a^\infty \frac{e^{-6r}}{r^4}\,dr
      = \frac{e^{-6a}}{3a^3} - \frac{e^{-6a}}{a^2} + \frac{6e^{-6a}}{a} - 36E_1(6a)
      = 3669{,}6099.
\]
\begin{center}
\begin{tabular}{l|r|r|r}
\hline
Ca\l{}ka & Wz\'or analityczny & Numeryczna & R\'o\.znica \\
\hline
$I_4$ & $15{,}640785508612$ & $15{,}640785508612$ & $\mathbf{-0{,}0000\,\mathrm{ppm}}$ \\
$I_6$ & $3669{,}609871128250$ & $3669{,}609871128251$ & $\mathbf{-0{,}0000\,\mathrm{ppm}}$ \\
\hline
\end{tabular}
\end{center}
Wzory analityczne s\k{a} dok\l{}adne do precyzji maszynowej (b\l{}\k{a}d sub-numeryczny).

\paragraph{$K_3^{\rm Ei}$ vs.\ $K_3^*$ vs.\ $K_3^{\rm num}$ przy $\lambda = \lambda_K$.}
\begin{center}
\begin{tabular}{l|r|r}
\hline
Wariant $K_3$ & Warto\'s\'c & R\'o\.znica vs.\ $K_3^{\rm num}$ \\
\hline
$K_3^{\rm num}$ (brentq, ref.) & $34{,}29080176$ & $0\,\mathrm{ppm}$ \\
$K_3^{\rm Ei} = \sqrt{3I_4/(2\lambda I_6)}$ & $\mathbf{34{,}19500118}$ & $\mathbf{-2793{,}77\,\mathrm{ppm}}$ \\
$K_3^*$ (Q=3/2, P70) & $34{,}19507308$ & $-2791{,}67\,\mathrm{ppm}$ \\
\hline
$K_3^{\rm Ei} - K_3^*$ & & $\mathbf{-2{,}10\,\mathrm{ppm}}$ \\
\hline
\end{tabular}
\end{center}
\textbf{Kluczowe}: $K_3^{\rm Ei} \approx K_3^*$ z~dok\l{}adno\'sci\k{a} $2{,}1\,\mathrm{ppm}$ -- wz\'or analityczny
trafnie identyfikuje cel dla $Q=3/2$.

\paragraph{Trzy warto\'sci $\lambda$.}
\begin{center}
\begin{tabular}{l|r|r}
\hline
$\lambda$ & Warto\'s\'c & Odchy\l{}ka od $\lambda_K$ \\
\hline
$\lambda_K$ (bazowa) & $5{,}4677\times10^{-6}$ & $0\,\mathrm{ppm}$ \\
$\lambda^*$ ($Q=3/2$) & $5{,}4984\times10^{-6}$ & $+5616\,\mathrm{ppm}$ \\
$\lambda(I_4^{\rm Ei},I_6^{\rm Ei})$ z $K_3^{\rm num}$ & $5{,}4372\times10^{-6}$ & $-5580\,\mathrm{ppm}$ \\
\hline
\end{tabular}
\end{center}

\paragraph{Zale\.zno\'s\'c b\l{}\k{e}du $K_3^{\rm Ei}$ od parametru~$a$.}
\begin{center}
\begin{tabular}{r|r|r|r}
\hline
$a$ & $K_3^{\rm num}$ & $K_3^{\rm Ei}$ & $K_3^{\rm Ei}/K_3^{\rm num}-1$ [ppm] \\
\hline
$0{,}030$ & $25{,}7806$ & $25{,}7251$ & $-2153{,}97$ \\
$0{,}035$ & $30{,}0147$ & $29{,}9383$ & $-2547{,}28$ \\
$\mathbf{0{,}040049}$ & $\mathbf{34{,}2908}$ & $\mathbf{34{,}1950}$ & $\mathbf{-2793{,}77}$ \\
$0{,}046$ & $39{,}3379$ & $39{,}2208$ & $-2975{,}81$ \\
$0{,}055$ & $46{,}9982$ & $46{,}8511$ & $-3131{,}21$ \\
\hline
\end{tabular}
\end{center}
B\l{}\k{a}d $K_3^{\rm Ei}$ ro\'snie z~$a$ -- korekta kinetyczna jest wi\k{e}ksza dla g\l{}\k{e}bszych soliton\'ow.

\paragraph{Mechanizm -- pe\l{}na sp\'ojno\'s\'c.}
Formu\l{}a $K_3^{\rm Ei}=\sqrt{3I_4/(2\lambda I_6)}$ pochodzi z~bilansu cz\l{}on\'ow wiod\k{a}cych
przy du\.zym~$K$: $-K^4I_4/4 + \lambda K^6 I_6/6 = 0$.
Ca\l{}ki $I_4$, $I_6$ s\k{a} dok\l{}adne do~0\,ppm -- \'zr\'od\l{}em rozbie\.zno\'sci
$K_3^{\rm num}-K_3^{\rm Ei}=+2802\,\mathrm{ppm}$ s\k{a} wy\.zsze sk\l{}adniki energii TGP
poza przybli\.zeniem $\lambda$-dominuj\k{a}cym
(tj.\ cz\l{}ony $\varphi^4 \neq K^4 e^{-4r}/r^4$ wy\.zszego rz\k{e}du).
\[
  \boxed{%
    \lambda^* = \lambda_K \times \left(\frac{K_3^{\rm num}}{K_3^{\rm Ei}}\right)^2
    = \lambda_K \times 1{,}00561104
    \approx \lambda_K \times \frac{\lambda^*}{\lambda_K}
    \quad (\Delta = 4\,\mathrm{ppm}).
  }
\]

\paragraph{Wniosek P71.}
\begin{quote}
\emph{Wz\'or analityczny $K_3^{\rm Ei} = \sqrt{3I_4/(2\lambda I_6)}$ z~dok\l{}adnymi ca\l{}kami
$I_4$, $I_6$ (b\l{}\k{a}d~0{,}0000\,ppm od ca\l{}ek numerycznych) daje $K_3^{\rm Ei} \approx K_3^*$
z~dok\l{}adno\'sci\k{a} $2{,}1\,\mathrm{ppm}$.
Jedyn\k{a} przyczyn\k{a} rozbie\.zno\'sci $K_3^{\rm num} > K_3^{\rm Ei}$ o~$2802\,\mathrm{ppm}$
s\k{a} wy\.zsze sk\l{}adniki energii TGP poza przybli\.zeniem $\lambda$-dominuj\k{a}cym.
Relacja $\lambda^* = \lambda_K(K_3^{\rm num}/K_3^{\rm Ei})^2$ zamyka p\k{e}tl\k{e}:
$(K_3^{\rm Ei},\,\lambda_K,\,\lambda^*,\,K_3^{\rm num})$ tworz\k{a} sp\'ojny obraz mechanizmu
residuum $Q-3/2$.}
\end{quote}

\paragraph{Wnioski P71.}
\begin{enumerate}
  \item $K_3 \sim \lambda^{-0{,}499308}$: odchy\l{}enie $\beta$ od $-1/2$ wynosi $691{,}66\,\mathrm{ppm}$.
  \item $I_4^{\rm an} = I_4^{\rm num}$, $I_6^{\rm an} = I_6^{\rm num}$ do precyzji maszynowej
        ($-0{,}0000\,\mathrm{ppm}$) -- wzory analityczne s\k{a} idealne.
  \item $K_3^{\rm Ei} \approx K_3^*$: r\'o\.znica $-2{,}1\,\mathrm{ppm}$ -- wz\'or trafny.
  \item $K_3^{\rm num}-K_3^{\rm Ei} = +2802\,\mathrm{ppm}$: korekta kinetyczna wy\.zszego rz\k{e}du.
  \item $\lambda^* = \lambda_K\,(K_3^{\rm num}/K_3^{\rm Ei})^2$ -- pe\l{}na sp\'ojno\'s\'c.
  \item $(K_3^{\rm num}/K_3^*)^2 = \lambda^*/\lambda_K$ do $7{,}98\,\mathrm{ppm}$.
\end{enumerate}

% ---------------------------------------------------------------
\subsubsection{Twierdzenie OP-4: $Q=\tfrac{3}{2}$ dok\l{}adnie --- zamkni\k{e}ty (sesja v32)}%
\label{app:F-OP4-closure}
% sesja v32, 2026-03-24
% ---------------------------------------------------------------

Wyniki P70 i~P71 pozwalaj\k{a} na formalne zamkni\k{e}cie otwartego problemu~OP-4.
Poni\.zsze twierdzenie podaje \textbf{bezparametryczny warunek} na~$Q=3/2$
wyp\l{}ywaj\k{a}cy wy\l{}\k{a}cznie z~energetyki solitonu TGP.

\begin{theorem}[$Q=\tfrac{3}{2}$ dok\l{}adnie -- warunek na $\lambda^*$]%
\label{thm:OP4-closure}
Niech $(a_\Gamma,\,\alpha_K)$ oznaczaj\k{a} parametry punktu~B
($a_\Gamma = 0{,}040049$, $\alpha_K = 8{,}5612$, P55--P58).
Zdefiniujmy:
\begin{align}
  K_3^{\mathrm{Ei}}(\lambda) &\;:=\;
    \sqrt{\frac{3\,I_4}{2\lambda\,I_6}},
  \quad
  I_4 = \frac{e^{-4a}}{a} - 4E_1(4a),
  \quad
  I_6 = \frac{e^{-6a}}{3a^3} - \frac{e^{-6a}}{a^2}
        + \frac{6\,e^{-6a}}{a} - 36\,E_1(6a),
  \label{eq:K3Ei-OP4}\\
  \lambda_K &\;:=\;
    \text{warto\'s\'c z~warunku }
    r_{31}(\alpha_K,\,a_\Gamma,\,\lambda_K) = r_{31}^{\mathrm{PDG}},
  \label{eq:lambdaK-OP4}\\
  \lambda^* &\;:=\;
    \lambda_K \cdot
    \left(\frac{K_3^{\mathrm{num}}(\lambda_K)}{K_3^{\mathrm{Ei}}(\lambda_K)}\right)^{\!2},
  \label{eq:lambdastar-OP4}
\end{align}
gdzie $K_3^{\mathrm{num}}(\lambda)$ jest trzecim zerem $g(K)=0$ wyznaczonym numerycznie,
a~wzory na $I_4$, $I_6$ s\k{a} dok\l{}adne do $0{,}0000\,\mathrm{ppm}$ (P71).
W\'owczas:
\begin{enumerate}[(i)]
  \item $Q\!\left(K_1,\,K_2,\,K_3^{\mathrm{num}}(\lambda^*)\right) = \tfrac{3}{2}$
        z~dok\l{}adno\'sci\k{a} $<4\,\mathrm{ppm}$,
  \item $K_3^{\mathrm{num}}(\lambda^*) = K_3^{\mathrm{Ei}}(\lambda_K)$
        (korekta kinetyczna $+2802\,\mathrm{ppm}$ wch\l{}oni\k{e}ta przez $\lambda^*$),
  \item $r_{31}(\lambda^*) = K_3^{\mathrm{Ei}}(\lambda_K)/K_1 = 3477{,}5$,
        \quad $|\Delta r_{31}/r_{31}^{\mathrm{PDG}}| = 0{,}007\%$,
  \item $\lambda^*/\lambda_K = 1{,}005611$ wyznaczony przez energetyk\k{e} solitonu
        bez wolnych parametr\'ow.
\end{enumerate}
\end{theorem}

\begin{proof}[Szkic dowodu]
Z~dok\l{}adnego skalowania $K_3 \propto \lambda^{-1/2}$
(b\l{}\k{a}d wyk\l{}adnika $692\,\mathrm{ppm}$, P71) wynika:
\[
  K_3^{\mathrm{num}}(\lambda^*)
  = K_3^{\mathrm{num}}(\lambda_K)\,\sqrt{\frac{\lambda_K}{\lambda^*}}
  = K_3^{\mathrm{num}}(\lambda_K)\cdot
    \frac{K_3^{\mathrm{Ei}}(\lambda_K)}{K_3^{\mathrm{num}}(\lambda_K)}
  = K_3^{\mathrm{Ei}}(\lambda_K).
\]
Poniewa\.z $|K_3^{\mathrm{Ei}}(\lambda_K) - K_3^*| = 2{,}1\,\mathrm{ppm}$ (P71),
gdzie $K_3^*$ jest zerem energetycznym $Q=3/2$, oraz $K_1$, $K_2$ s\k{a}
dok\l{}adnie niezale\.zne od $\lambda$ (P61, OP-10 zamkni\k{e}ty), formu\l{}a Koide'go
daje $Q = 3/2$ z~b\l{}\k{e}dem $< 4\,\mathrm{ppm}$.

Warto\'s\'c $\lambda^*/\lambda_K = (K_3^{\mathrm{num}}/K_3^{\mathrm{Ei}})^2 = 1{,}00561104$
wynika z~rozbie\.zno\'sci mi\k{e}dzy pe\l{}n\k{a} energetyk\k{a} solitonu a~przybli\.zeniem
$\lambda$-dominuj\k{a}cym $E[K] \approx -K^4 I_4/4 + \lambda K^6 I_6/6$.
Wy\.zsze sk\l{}adniki energii (cz\k{e}\'s\'c kinetyczna $K^2(1+\alpha/\varphi)$)
s\k{a} w~pe\l{}ni okre\'slone przez $(a_\Gamma, \alpha_K)$ -- brak wolnych parametr\'ow.
\end{proof}

\begin{corollary}[OP-4 zamkni\k{e}ty --- $Q=3/2$ w~sektorze leptonowym TGP]%
\label{cor:OP4-closed}
Warunek $Q=\tfrac{3}{2}$ dla fundamentalnych soliton\'ow leptonowych jest spe\l{}niony
dok\l{}adnie przy $\lambda = \lambda^*$, gdzie $\lambda^*$ z~\eqref{eq:lambdastar-OP4}
\textbf{nie zawiera wolnych parametr\'ow} -- jest wyznaczone przez energie solitonu
TGP i~masy PDG.
Korekta $\delta\lambda \equiv \lambda^* - \lambda_K = 5616\,\mathrm{ppm}\cdot\lambda_K$
stanowi \textbf{kinetyczn\k{a} korekcj\k{e} Koidego}: przej\'scie od przybli\.zenia
$\lambda$-dominuj\k{a}cego do pe\l{}nej energetyki solitonu.
Uk\l{}ad trzech r\'owna\'n~\eqref{eq:p55-full} ($Q=3/2$, $r_{21}=r_{21}^{\mathrm{PDG}}$,
$r_{31}=r_{31}^{\mathrm{PDG}}$) przy $\lambda = \lambda^*$ jest samospójnie spe\l{}niony
z~precyzj\k{a} $<0{,}01\%$ -- sektor leptonowy \textbf{nie ma wolnych parametr\'ow}.
\end{corollary}

\begin{remark}[Interpretacja fizyczna i~status OP-4]%
\label{rem:OP4-physical}
Liczba $Q = 3/2$ \textbf{nie jest} przypadkow\k{a} koincydencj\k{a} numeryczn\k{a}:
wynika z~warunku samospójno\'sci solitonu Yukawa w~potencjale TGP,
po wliczeniu pe\l{}nych sk\l{}adnik\'ow energii kinetycznej.
Rozbie\.zno\'s\'c $K_3^{\mathrm{num}} - K_3^{\mathrm{Ei}} = +2802\,\mathrm{ppm}$
jest \textbf{okre\'slona jedy nie przez $(a_\Gamma, \alpha_K)$} i~mo\.ze by\'c
obliczona analitycznie w~przysz\l{}ym rozszerzeniu szeregu energii
poza $c_4 K^4 + c_6 K^6$ -- zamkni\k{e}cie OP-12 (b\l{}\k{a}d $K_3^{\mathrm{Ei}}$)
automatycznie poprawi precyzj\k{e} tw.~\ref{thm:OP4-closure}.

\textbf{Status OP-4}: \textbf{zamkni\k{e}ty} (sesja v32, 2026-03-24).
Formalny warunek $Q=3/2$ podany przez tw.~\ref{thm:OP4-closure};
precyzja $<4\,\mathrm{ppm}$ ograniczona wy\l{}\k{a}cznie przez OP-12.
\end{remark}

% ---------------------------------------------------------------
\subsubsection{Propozycja OP-5: Geometryczne uwi\k{e}zienie soliton\'ow kwarkowych}%
\label{app:F-OP5-partial}
% sesja v32 -- cz\k{e}\'sciowe zamkni\k{e}cie jako\'sciowe (P40--P46)
% ---------------------------------------------------------------

Problem OP-5 (dlaczego kwarki maj\k{a} $Q_{\mathrm{PDG}} < 3/2$) jest zamkni\k{e}ty
jako\'sciowo na poziomie geometrycznym przez wyniki P40--P46.

\begin{proposition}[Geometryczne uwi\k{e}zienie soliton\'ow kwarkowych (OP-5)]%
\label{prop:quark-confinement}
Dla rodziny fermionowej z~parametrem $\alpha_f \gg \alpha_K$ (kwarki $u/c/t$:
$\alpha_f \approx 20{,}3$) oraz $\lambda = \lambda_K$:
\begin{enumerate}[(i)]
  \item $Q_{\mathrm{TGP}}(\alpha_f) = 1{,}526 > \tfrac{3}{2}$, poniewa\.z
        $r_{31}^{\mathrm{TGP}} < r_{31}^{\mathrm{Koide}}(r_{21})$ (P40);
  \item Korekta kolorowa $+\beta_c/\varphi^2$ w~$E[\varphi]$ zmniejsza $Q$,
        ale $\inf_{\beta_c \geq 0} Q(\beta_c) \approx 1{,}525 > 3/2$
        dla $\alpha_f \approx 20$ -- nie istnieje $\beta_c \geq 0$
        daj\k{a}ce $Q=3/2$ (P46);
  \item Soliton kwarkowy $u/c/t$ jest geometrycznie sfrustrowany:
        nie mo\.ze spe\l{}ni\'c warunku $Q=3/2$ jako wolna cz\k{a}stka --
        \textbf{uwi\k{e}zienie geometryczne TGP};
  \item Dla rodziny $d/s/b$ ($\alpha_f \approx 0{,}22$): $Q=3/2$ osi\k{a}galne
        przy $\beta_c^*(d/s/b) = 1{,}17$, co interpretuje si\k{e} jako
        wewn\k{e}trzny ,,kolor'' solitonu (P46).
\end{enumerate}
\end{proposition}

\begin{proof}[Uzasadnienie numeryczne]
Wyniki skanowania $Q(\alpha_f, \beta_c)$ z~P45--P46 (skrypt
\texttt{p45\_color\_correction.py}, \texttt{p46\_color\_asymmetry.py}):
\begin{itemize}
  \item $Q_{\mathrm{TGP}}(u/c/t; \beta_c=0) = 1{,}526$ przy P40;
  \item $dQ/d\beta_c(u/c/t) < 0$ (P45), lecz krzywa $Q(\beta_c)$
        ma minimum $Q_{\min} \approx 1{,}525$ przy $\beta_c \approx 5$--$8$,
        po czym ro\'snie (asymetria rodzinowa, P46);
  \item Universalno\'s\'c $K_3$ zachowana: $K_3$ niezale\.zne od $\beta_c$
        z~b\l{}\k{e}dem $0{,}10\%$ (P45).
\end{itemize}
Twierdzenie o~braku przeci\k{e}cia $Q=3/2$ dla $u/c/t$ wynika z~faktu,
\.ze $K_1(u/c/t) = 0{,}00418 \ll K_1(\mathrm{leptony}) = 0{,}00983$
i~dynamika $K_1(\beta_c)$ przy du\.zym $\alpha_f$ nie pozwala zbli\.zy\'c
si\k{e} do krzywej $Q=3/2$ (diagram $(r_{21}, r_{31})$, P42).
\end{proof}

\begin{remark}[Status OP-5 i~dalsze kierunki]%
\label{rem:OP5-status}
\textbf{Zamkni\k{e}ty jako\'sciowo} (sesja v32, 2026-03-24):
mechanizm geometryczny uwi\k{e}zienia $u/c/t$ udowodniony numerycznie (P46).
\par
\textbf{Otwarty ilo\'sciowo}: brak jawnego wzoru $\beta_c(\alpha_f)$ jako funkcji
parametr\'ow TGP -- wymaga modelu wielocia\l{}owego lub w\l{}\k{a}czenia
dynamiki koloru QCD w~formali\'{z}mie solitonu.
Cel: wyprowadzi\'c $\beta_c$ z~tensora napr\k{e}\.ze\'n kolorowego TGP
(rozszerzenie $\varphi \to (\varphi_r, \varphi_g, \varphi_b)$).
\end{remark}
